\documentclass[10pt,a4paper]{article}
\usepackage{fontspec,amsmath,amssymb,amsthm,mathtools,graphicx,bm,bbm,stmaryrd,nicefrac,mathrsfs}
\usepackage[margin=22mm]{geometry}
\usepackage[bookmarks=false]{hyperref}
\setlength{\emergencystretch}{3em}
\title{Witness encryption characterized by honest half-Levin reductions to circuit gap learning}
\author{Halley Goldberg, Valentine Kabanets}
\date{doi:10.4230/LIPIcs.CCC.2025.34; CC BY 4.0}
\begin{document}
\maketitle
\noindent\textit{Editorial note.} Editable TeX converted from the official CC BY HTML prose and math alttext. Original official PDF/HTML retained. Only the specified claim counts; all author context and available source-local supporting proofs are included.\par\medskip


\par\medskip\noindent\textbf{
1 Introduction and Background}\par




It has long been recognized that learning and cryptography are two sides of the same coin. As Valiant already observed in his paper introducing PAC learning, its easiness would preclude the existence of pseudorandom functions [46]. In 1993, building on a work of Impagliazzo and Levin [24], Blum, Furst, Kearns, and Lipton sharpened Valiant’s observation by showing that private-key cryptography can be broken if and only if PAC learning is easy in an average-case setting [6]. A number of more recent results have developed such connections further (eg. [41, 39, 18]).\par




In the present work, we turn to more structured kinds of hardness of learning (namely, ${\mathsf{NP}}$-hardness) and cryptography (namely, witness encryption), which we show to be likewise equivalent.\par




\par\medskip\noindent\textbf{NP-hardness of PAC learning}\par




Learning theory asks whether an efficient entity can hope to learn general truths about reality from necessarily limited input from experience. Valiant formalized this idea in 1984 with his definition of Probably Approximately Correct (PAC) Learning [46]. In this framework, a learner for some representation class $\mathcal{C}$ gets access to labeled examples $(x,f(x))$ for some unknown concept $f\in\mathcal{C}$ and $x$ randomly sampled according to an arbitrary distribution $\mathcal{D}$. The learner is asked to produce, with high probability, a hypothesis that approximates $f$ well over $\mathcal{D}$. One can distinguish between “proper” and “improper” settings: in proper learning, the learner should produce a hypothesis that also belongs to $\mathcal{C}$, whereas in improper learning, there is no restriction on the complexity of the hypothesis. One can also consider various “semi-proper” settings, in which the hypothesis must belong to some specific concept class $\mathcal{C^{\prime}}$ containing $\mathcal{C}$. Lastly, one can relax the model of learning to be “agnostic” in that the labeled examples $(x,b)$ do not necessarily reflect the values of a function $f\in\mathcal{C}$, and the learner is only asked to do as well as the best function in $\mathcal{C}$ does [30].\par




A long line of work has established ${\mathsf{NP}}$-hardness of PAC learning for various concept classes (eg. [43, 1, 16, 32]). Early results obtained ${\mathsf{NP}}$-hardness of proper learning of restricted circuit classes. Over time, progress has pushed toward ${\mathsf{NP}}$-hardness of “less proper” learning of less restricted circuit classes. However, ${\mathsf{NP}}$-hardness of learning \emph{unrestricted} boolean circuits in the \emph{improper} setting remains elusive. By this, we refer to the problem of learning the concept class ${\mathsf{SIZE}}[s(n)]$ for some fixed polynomial $s$ by a hypothesis of arbitrary polynomial size. Proving ${\mathsf{NP}}$-hardness in this setting seems especially challenging in light of a 2008 result of Applebaum et al.:\par



% SOURCE-ID Thmtheorem1


\par\medskip\noindent\textbf{
\textbf{Theorem 1}\textbf{ }([2])\textbf{.}
}\par



\emph{Suppose there is a randomized non-adaptive honest\footnote{Roughly, “honest” means that the input length of the learning instance is polynomially related to the length of the input to the reduction. See Definition 28.} reduction $R$ and a polynomial $s:\mathbb{N}\to\mathbb{N}$ such that, for every constant $c\in\mathbb{N}$, $R$ reduces ${\mathsf{SAT}}$ to agnostically PAC learning ${\mathsf{SIZE}}[s(n)]$ by ${\mathsf{SIZE}}[s(n)^{c}]$. Then, if ${\mathsf{NP}}$ is hard on average, infinitely-often one-way functions exist.}\par



% END-SOURCE-ID Thmtheorem1



Of course, there are various ways in which one could formulate reductions to PAC learning. The aforementioned paper [2] describes two (of many) possibilities. The statement of Theorem 1 refers to one kind of ${\mathsf{NP}}$-hardness reduction described in which a randomized oracle machine on an input of length $n\in\mathbb{N}$ makes non-adaptive queries $\mathcal{E}$ to its oracle, with $\mathcal{E}$ being a ${\mathsf{poly}}(n)$-size circuit sampling example-label pairs $(x,b)\in\{0,1\}^{n^{\Omega(1)}}\times\{0,1\}$. The oracle returns, for each query, a polynomial-size circuit $h$ having high agreement over $\mathcal{E}$ (i.e. $\Pr_{(x,b)\sim\mathcal{E}}[h(x)=b]$ close to 1), which the reduction may use in any way.\par




In the other kind of ${\mathsf{NP}}$-hardness reduction that [2] describes, the reduction only uses its learning oracle to \emph{decide} if a small circuit consistent with $\mathcal{E}$ exists. The relevant decision problem introduced by [2] is the \emph{Computational Gap Learning Problem} (${\sf CGL}$). The authors show that if ${\sf CGL}$ is ${\mathsf{NP}}$-hard under a deterministic many-one reduction, then every language in ${\mathsf{NP}}$ has a statistical zero-knowledge argument.\par



% SOURCE-ID Thmtheorem2


\par\medskip\noindent\textbf{
$\blacktriangleright$ Remark 2.
}\par



\emph{The results in [2] indicate it is a challenge to show ${\mathsf{NP}}$-hardness of agnostic PAC learning, as it would prove breakthrough results in cryptography, which, although conjectured to be true, seem beyond our reach at the moment. However,
we should note that the assumption in Theorem 1 is actually quite strong in its requirement that the same reduction must work for every constant $c\in\mathbb{N}$. Theorem 1 also requires that queries $\mathcal{E}$ made by the reduction sample pairs $(x,b)$ with the length of $x$ polynomially related to the length of the given ${\mathsf{SAT}}$ instance. It is possible that PAC learning is ${\mathsf{NP}}$-hard in the improper setting without either of these conditions being met. For example, if a different reduction exists for each $c\in\mathbb{N}$ (perhaps taking longer time for larger $c$), then the easiness of improper learning would still imply the easiness of ${\mathsf{NP}}$.}\par



% END-SOURCE-ID Thmtheorem2




\par\medskip\noindent\textbf{Witness encryption for NP}\par




First introduced in a 2013 paper of Garg, Gentry, Sahai, and Waters [11], witness encryption (WE) is a cryptographic primitive in which security is contingent on the ability to provide a proof of an answer to a hard problem. As Garg et al. put it,\par



What if we don’t really care if [the receiver] knows a secret key, but we do care if he knows a solution to a crossword puzzle that we saw in the Times? [11]\par



More formally, the functionality of a witness encryption scheme is based on the structure of a particular language $L\in{\mathsf{NP}}$. An instance $x$ of $L$ is made public, and anyone can encrypt a secret message using the witness encryption scheme along with $x$. It is guaranteed that if $x$ belongs to $L$, then anyone with a witness proving $x\in L$ can decrypt, but if $x$ does not belong to $L$, then no one can do so.\par




Witness encryption is intimately connected to another cryptographic primitive known as indistinguishability obfuscation (iO) [3], which has been studied extensively (eg. [33, 10, 23, 37]) and has many candidate constructions (see [28]). iO is known to imply WE, though the converse is not known [10]. (We give a sketch of the proof that iO implies WE in Appendix A.) WE is also closely related to secret-sharing, another important cryptographic tool [4, 34]. In fact, the candidate construction of WE in [11] also yields a Rudich-type secret sharing scheme. As with much of cryptography, we have no unconditional proof that WE is possible.\par




WE alone does not imply the existence of one-way functions, since it is trivially possible if ${\mathsf{NP}}\subseteq{\mathsf{BPP}}$. However, WE \emph{along with} one-way functions becomes very powerful, together yielding public key encryption, as demonstrated in [11]. By combining this result with two subsequent works, one obtains the following.\par



% SOURCE-ID Thmtheorem3


\par\medskip\noindent\textbf{
\textbf{Theorem 3}\textbf{ }([11, 33, 19]; see [35])\textbf{.}
}\par



\emph{Assume ${\mathsf{NP}}\nsubseteq{\mathsf{io\text{-}}}{\mathsf{P}}/{\mathsf{poly}}$. If ${\mathsf{P}}/{\mathsf{poly}}$-secure witness encryption exists for ${\mathsf{NP}}$, then public-key encryption is possible.}\par



% END-SOURCE-ID Thmtheorem3



In this paper, we show that ${\sf CGL}$ is ${\mathsf{NP}}$-hard if and only if witness encryption exists for ${\mathsf{NP}}$. We prove statements along these lines for a few different settings: WE secure against PPT adversaries, WE secure against polynomial-size circuits, and WE with logarithmic-length ciphertexts; the latter roughly characterizes the ${\mathsf{NP}}$-hardness of a promise version of the Minimum Circuit Size Problem [29].\par




In contrast to the aforementioned results of [2] (such as Theorem 1), not only do we show that proving ${\mathsf{NP}}$-hardness of learning is a challenge as it would imply a breakthrough cryptographic construction (namely, WE), but we also show the \emph{converse}: progress in cryptography (showing that WE is possible) would imply that agnostic PAC learning is ${\mathsf{NP}}$-hard. Hence, if one believes in cryptography such as iO or WE, then one should also believe in the ${\mathsf{NP}}$-hardness of improper agnostic learning, even in the extremely restrictive sense of a deterministic, many-one, half-Levin reduction to ${\sf CGL}$. We interpret our results as evidence that improper learning is likely ${\mathsf{NP}}$-hard.\par




As a complement, we prove a number of \emph{unconditional} ${\mathsf{NP}}$-hardness results for agnostic PAC learning. For one, we obtain ${\mathsf{NP}}$-hardness of agnostic learning for unrestricted boolean circuits in a semi-proper setting that improves on the previous state of the art. In other settings, we obtain ${\mathsf{NP}}$-hardness of \emph{improper} agnostic learning. Though we have made some partial progress toward showing ${\mathsf{NP}}$-hardness unconditionally, we feel that it remains a promising direction for further research.\par




Lastly, we give some applications of our results for the possibility of private- and public-key cryptography. We improve Theorem 1 to conditionally obtain infinitely-often one-way functions from \emph{worst-case} hardness of ${\mathsf{NP}}$. In contrast, if learning is ${\mathsf{NP}}$-hard in a more restricted sense (namely, a half-Levin reduction to ${\sf CGL}$), then one obtains infinitely-often public-key encryption assuming ${\mathsf{NP}}\nsubseteq{\mathsf{BPP}}$. Along the way, we prove Theorem 3 under the weaker uniform assumption ${\mathsf{NP}}\nsubseteq{\mathsf{BPP}}$.\footnote{We believe that this last statement follows from a combination of techniques used in prior work ([11, 33, 19]; see [35]), but we have not seen the uniform version stated. In any case, we offer an alternative proof that does not rely on properties of statistical zero-knowledge arguments.}\par





\par\medskip\noindent\textbf{
1.1 Our Results}\par




\par\medskip\noindent\textbf{Connections between witness encryption and NP-hardness of learning}\par




We first discuss a set of results that characterize witness encryption in terms of ${\mathsf{NP}}$-hardness reductions to PAC learning.\par




\emph{Witness encryption} for an ${\mathsf{NP}}$ language $L$ consists of a pair of algorithms $({\mathsf{Enc}},{\mathsf{Dec}})$ such that, for an $L$-instance $x$ and a secret bit $b\in\{0,1\}$, ${\mathsf{Enc}}(x,b)$ outputs a ciphertext string $c$ using some randomness. It is guaranteed that, if $x\in L$ and $w$ is a witness for the membership of $x$ under a fixed witness relation for $L$, then ${\mathsf{Dec}}(c,x,w)$ deterministically recovers the bit $b$. On the other hand, if $x\notin L$, then ${\mathsf{Enc}}(x,0)$ and ${\mathsf{Enc}}(x,1)$ are indistinguishable (within a negligible difference) for polynomial-size circuits (if the witness encryption is ${\mathsf{P}}/{\mathsf{poly}}$-secure) or PPT algorithms (if the witness encryption is ${\mathsf{BPP}}$-secure). See Definition 39.\par




Our first result characterizes ${\mathsf{P}}/{\mathsf{poly}}$-secure witness encryption in terms of a deterministic many-one reduction to a promise problem known as ${\sf CGL}$, introduced in a work of Applebaum et al. [2]. Inputs to ${\sf CGL}$ are distributions $\mathcal{E}$ over string-label pairs $(x,b)\in\{0,1\}^{n}\times\{0,1\}$, where $\mathcal{E}$ is represented by a ${\mathsf{poly}}(n)$-size sampling circuit. For parameters $s_{1}<s_{2}\in\mathbb{N}$ and $0<\varepsilon<1/2$, a distribution $\mathcal{E}$ is a yes-instance of ${\sf CGL}_{s_{1}}^{s_{2}}[\varepsilon]$ if there is a circuit $C$ of size at most $s_{1}$ such that, for $(x,b)\sim\mathcal{E}$, $\Pr[C(x)=b]=1$; $\mathcal{E}$ is a no-instance of ${\sf CGL}_{s_{1}}^{s_{2}}[\varepsilon]$ if for every circuit $C^{\prime}$ of size at most $s_{2}$, for $(x,b)\sim\mathcal{E}$, $\Pr[C^{\prime}(x)=b]<1/2+\varepsilon$.\par




Note that a many-one reduction to ${\sf CGL}$ is quite a restrictive notion of reducibility to agnostic PAC learning: the reduction only uses its learning oracle to \emph{decide} if a small circuit exists for the given distribution. As discussed earlier, one can imagine other kinds of reduction to learning that make more robust use of a learner: for example, actually examining the circuits returned by the learning oracle.\par




We further restrict the definition of a many-one reduction to ${\sf CGL}$ by requiring that it be a \emph{half-Levin reduction}. This is a term that we introduce here for a kind of reduction that is intermediate between a many-one reduction and a Levin reduction. A half-Levin reduction $(R,R_{wit})$ from $L_{1}$ to $L_{2}$ is a pair of polynomial-time machines such that $R$ is a many-one reduction from $L_{1}$ to $L_{2}$, and $R_{wit}$ transforms witnesses for $x\in L_{1}$ into witnesses for $R(x)\in L_{2}$. In contrast to a standard Levin reduction, we do not require a third algorithm transforming $L_{2}$-witnesses into $L_{1}$-witnesses. We emphasize that the half-Levin reductions considered in this work are deterministic unless stated otherwise.\par




We call a reduction to ${\sf CGL}$ \emph{honest} if, on inputs of length $n\in\mathbb{N}$, it outputs a distribution $\mathcal{E}$ supported over $\{0,1\}^{n^{\Omega(1)}}\times\{0,1\}$.\par




We are now ready to state our main result. In general, fixing parameters $s_{1},s_{2}$, and $\varepsilon$, we show that a half-Levin reduction from a language $L$ to ${\sf CGL}_{s_{1}}^{s_{2}}[\varepsilon]$ is equivalent to a witness encryption scheme for $L$ with running time roughly $s_{1}$, security against circuits of size roughly $s_{2}$, and with advantage parameter roughly $\varepsilon$; see Lemmas 52 and 54. For the special case of ${\mathsf{P}}/{\mathsf{poly}}$-secure witness encryption, we get the following.\par




% SOURCE-ID Thmtheorem4


\par\medskip\noindent\textbf{
\textbf{Theorem 4}\textbf{.}
}\par



\emph{Consider any language $L\in{\mathsf{NP}}$. The following are equivalent.}\par


\begin{enumerate}


\item[] 
1. 


${\mathsf{P}}/{\mathsf{poly}}$\emph{-secure witness encryption exists for }$L$\emph{;}\par





\item[] 
2. 


\emph{there exists a polynomial }$s_{1}$\emph{ and a pair of machines }$(R,R_{wit})$\emph{ such that, for all polynomials }$s_{2},\varepsilon^{-1}$\emph{, }$(R,R_{wit})$\emph{ is an honest half-Levin reduction from }$L$\emph{ to }${\sf CGL}^{s_{2}}_{s_{1}}[\varepsilon]$\emph{.}\par





\end{enumerate}



% END-SOURCE-ID Thmtheorem4




We also note that the above characterization is related to a very recent result from Liu, Mazor, and Pass, which shows that ${\mathsf{P}}/{\mathsf{poly}}$-secure witness encryption for an ${\mathsf{NP}}$-language $L$ is equivalent to the existence of a laconic special-honest verifier zero-knowledge argument for $L$ [36]. We refer to that paper for the definition of this kind of protocol, though we mention that the term “laconic” means that the total length of the prover’s messages is bounded by $O(\log n)$. As a corollary of [36] and our Theorem 4 above, we obtain the following.\par



% SOURCE-ID Thmtheorem5


\par\medskip\noindent\textbf{
\textbf{Corollary 5}\textbf{.}
}\par



\emph{Consider any language $L\in{\mathsf{NP}}$. The following are all equivalent.}\par


\begin{enumerate}


\item[] 
1. 


${\mathsf{P}}/{\mathsf{poly}}$\emph{-secure witness encryption exists for }$L$\emph{;}\par





\item[] 
2. 


\emph{there exists a polynomial }$s_{1}$\emph{ and a pair of machines }$(R,R_{wit})$\emph{ such that, for all polynomials }$s_{2},\varepsilon^{-1}$\emph{, }$(R,R_{wit})$\emph{ is an honest half-Levin reduction from }$L$\emph{ to }${\sf CGL}^{s_{2}}_{s_{1}}[\varepsilon]$\emph{;}\par





\item[] 
3. 


\emph{there exists a laconic special-honest verifier zero-knowledge argument for }$L$\emph{.}\par





\end{enumerate}



% END-SOURCE-ID Thmtheorem5



One may also be interested in the possibility that ${\mathsf{NP}}$ has witness encryption schemes secure against PPT algorithms but not against polynomial-size circuits. We therefore present a kind of reduction to learning that is equivalent to ${\mathsf{BPP}}$-secure witness encryption. This kind of reduction is intermediate between a half-Levin reduction to ${\sf CGL}$ and the more general kind of reduction described in [2] wherein the reduction may actually inspect the hypotheses returned by the learner and use them in any way.\par




More specifically, in our case, the reduction makes a single query $\mathcal{E}$ to the \emph{search} version of ${\sf CGL}$, getting back a small circuit consistent with $\mathcal{E}$ if $\mathcal{E}$ is a yes-instance of ${\sf CGL}$. We still require the reduction to be half-Levin. Furthermore, we require the reduction to be ${\mathsf{BPP}}$-black-box;\footnote{The standard notion of a \emph{class-specific black-box} reduction from a language $L_{1}$ to a language $L_{2}$ was defined by [13] as a way to formalize an \emph{intermediate} case between treating $L_{2}$ as an oracle (“black box”) and requiring a source code for an $L_{2}$ algorithm implemented in some complexity class (“white box”); see Definitions 30 and 31.} see Definition 32.\par



% SOURCE-ID Thmtheorem6


\par\medskip\noindent\textbf{
\textbf{Theorem 6}\textbf{.}
}\par



\emph{Consider any language $L\in{\mathsf{NP}}$. The following are equivalent.}\par


\begin{enumerate}


\item[] 
1. 


${\mathsf{BPP}}$\emph{-secure witness encryption exists for }$L$\emph{;}\par





\item[] 
2. 


\emph{there exist a polynomial }$s_{1}$\emph{ and a pair of machines }$(R,R_{wit})$\emph{ such that, for all polynomials }$s_{2}$\emph{ and }$\varepsilon^{-1}$\emph{, }$(R,R_{wit})$\emph{ is an honest }${\mathsf{BPP}}$\emph{-black-box half-Levin reduction from }$L$\emph{ to }${\sf search}\text{-}{\sf CGL}^{s_{2}}_{s_{1}}[\varepsilon]$\emph{.}\par





\end{enumerate}



% END-SOURCE-ID Thmtheorem6



So far, we have only discussed learning in the “polynomial regime”, where the distribution $\mathcal{E}$ over pairs $(x,b)$ is in ${\mathsf{P}}{\mathsf{SAMP}}/{\mathsf{poly}}$. However, the original PAC learning framework of Valiant makes no such restriction: a PAC learner should be able to learn over \emph{arbitrary} distributions [46]. We thus consider reductions to learning that make further use of this capacity. In particular, imagine a polynomial-time reduction that, on input length $n$, queries a learner for a function on $O(\log n)$-length inputs. In this case, it will be possible to ask the learner to work on inefficiently samplable distributions, and a truth-table of the function in question can be given to the learner explicitly. In this particular setup, the distributions queried represent well-defined functions, which may not be the case in the setup above with ${\sf CGL}$. In this sense, we get closer to standard PAC learning rather than agnostic PAC learning.\par




This notion of learning in the “exponential regime” is closely related to ${\mathsf{MCSP}}$, which asks, given a truth-table $T$ and a size threshold $s$, whether there exists a circuit of size $s$ consistent with $T$. We generalize the problem to the context of learning by providing, along with $T$, a probability distribution $\mu$ represented as a table. We call this problem “Gap Distributional ${\mathsf{MCSP}}$”, denoted ${\mathsf{GapDistMCSP}}$. Specifically, $(T,\mu)$ is a yes-instance of ${\mathsf{GapDistMCSP}}_{s_{1}}^{s_{2}}[\varepsilon]$ if there exists a circuit $C$ of size at most $s_{1}$ such $C$ is consistent with $T$, and $(T,\mu)$ is a no-instance if for every circuit $C^{\prime}$ of size at most $s_{2}$, the probability over $x\sim\mu$ that $C(x)=T(x)$ is at most $1/2+\varepsilon$. Note that this definition naturally generalizes ${\sf Ave}{\mathsf{MCSP}}$ as defined in [45] (and in [7], denoted ${\sf MACSP}$), where the definition is the same except the distribution fixed to uniform.\par




We show that ${\mathsf{NP}}$-hardness of learning in this sense is equivalent to witness encryption for ${\mathsf{NP}}$ such that the ciphertexts output by ${\mathsf{Enc}}$ have only \emph{logarithmic} length, along with a circuit lower bound for ${\mathsf{E}}$. Below we use ${\mathsf{SIZE}}[s(n)]$ to denote the class of $n$-input boolean functions computable by circuits of size $s(n)$; for $0\leq\alpha\leq 1$, we use ${\mathsf{SIZE}}_{D}^{\alpha}[s(n)]$ to denote the class of $n$-input boolean functions that can be computed by circuits of size $s(n)$ on average, with probability at least $\alpha$ over inputs $x\in\{0,1\}^{n}$ sampled according to the distribution $D$.\par



% SOURCE-ID Thmtheorem7


\par\medskip\noindent\textbf{
\textbf{Theorem 7}\textbf{.}
}\par



\emph{Consider any language $L\in{\mathsf{NP}}$, constant $c\in\mathbb{N}$, polynomial $s_{1}$, and $s_{2},\varepsilon^{-1}:\mathbb{N}\to\mathbb{N}$. The following are equivalent.}\par


\begin{enumerate}


\item[] 
1. 


\begin{itemize}


\item[] 
$\blacksquare$ 


$\left({\mathsf{SIZE}}[\Omega(s_{2}(n))],O(\varepsilon(n))\right)$\emph{-secure }$(c\log n+O(1))$\emph{-length }$\widetilde{O}(s_{1}(n))$\emph{-decryption-time witness encryption exists for }$L$\emph{; and}\par





\item[] 
$\blacksquare$ 


${\mathsf{E}}\nsubseteq{\mathsf{io\text{-}}}{\mathsf{SIZE}}_{D}^{\frac{1}{2}+O(\varepsilon(n))}[\Omega(s_{2}(n))]$\emph{ for some }${\mathsf{E}}$\emph{-computable distribution family }$D$\emph{.}\par





\end{itemize}





\item[] 
2. 


\emph{There is a half-Levin reduction }$(R,R_{wit})$\emph{ mapping }$n$\emph{-length instances of }$L$\emph{ to }$O(n^{c})$\emph{-length instances of
}${\mathsf{GapDistMCSP}}_{\widetilde{O}(s_{1}(n))}^{\Omega(s_{2}(n))}[O(\varepsilon)]$\emph{, where }$R_{wit}$\emph{ runs in time at most }$\widetilde{O}(s_{1}(n))$\emph{.}\par





\end{enumerate}



% END-SOURCE-ID Thmtheorem7



Note that we obtain correspondences between the output length of the reduction $R$ and the output length of ${\mathsf{Enc}}$, the $s_{1}$ parameter of ${\mathsf{GapDistMCSP}}$ and the running time of ${\mathsf{Dec}}$, the $s_{2}$ parameter of ${\mathsf{GapDistMCSP}}$ and the security of witness encryption, and the $\varepsilon$ parameter of ${\mathsf{GapDistMCSP}}$ and the advantage parameter of witness encryption.\par





\par\medskip\noindent\textbf{Unconditional NP-hardness of semi-proper learning for ${\mathsf{P}}/{\mathsf{poly}}$
}\par




As discussed in the introduction, ${\mathsf{NP}}$-hardness of proper agnostic PAC-learning for polynomial-size circuits is known. However, ${\mathsf{NP}}$-hardness of improper learning is a major open question that remains elusive. In this section, we advance toward ${\mathsf{NP}}$-hardness of “less proper” agnostic PAC learning for polynomial-size circuits.\par




First, we show that for polynomials $s_{1}$ and $\varepsilon^{-1}$, letting $s_{2}(n)=s_{1}(n)\cdot n^{1/(\log\log n)^{O(1)}}$, ${\sf CGL}_{s_{1}}^{s_{2}}[\varepsilon]$ is ${\mathsf{NP}}$-hard under a randomized many-one reduction. In particular,\par



% SOURCE-ID Thmtheorem8


\par\medskip\noindent\textbf{
\textbf{Theorem 8}\textbf{.}
}\par



\emph{For some constant $c\in\mathbb{N}$ and $g:\mathbb{N}\to\mathbb{N}$ with $g(n)=n^{1/(\log\log n)^{O(1)}}$ for $n\in\mathbb{N}$, for any polynomials $\varepsilon^{-1}$ and $s_{1}$ such that $s_{1}(n)\geq\varepsilon^{-c}(n)$ for $n\in\mathbb{N}$, there is a randomized many-one reduction from ${\mathsf{SAT}}$ to ${\sf CGL}_{s_{1}}^{s_{1}\cdot g}[\varepsilon]$.}\par



% END-SOURCE-ID Thmtheorem8



As a straightforward consequence of Theorem 8, we obtain ${\mathsf{NP}}$-hardness of agnostic PAC-learning in the semi-proper setting of learning size-$s(n)$ circuits by size-$s(n)\cdot n^{1/(\log\log n)^{O(1)}}$ circuits.\par



% SOURCE-ID Thmtheorem9


\par\medskip\noindent\textbf{
\textbf{Corollary 9}\textbf{.}
}\par



\emph{For some polynomial $s$, it is ${\mathsf{NP}}$-hard under a randomized one-query reduction to agnostically PAC learn ${\mathsf{SIZE}}[s(n)]$ by ${\mathsf{SIZE}}[s(n)\cdot n^{1/(\log\log n)^{O(1)}}]$ over a flat ${\mathsf{P}}/{\mathsf{poly}}$-samplable distribution.}\par



% END-SOURCE-ID Thmtheorem9



This improves on a prior work of Hirahara [16], which contains two main results. The first shows that, for polynomials $s$, it is ${\mathsf{NP}}$-hard to agnostically learn \emph{programs} of size $s(n)$ by \emph{programs} of size $s(n)\cdot n^{1/(\log\log n)^{O(1)}}$. This is accomplished by reduction from the “minimum monotone satisfying assignment” problem, ${\mathsf{MMSA}}$, which is known to be ${\mathsf{NP}}$-hard to approximate by way of a PCP theorem [8, 9]. We note that Hirahara’s proof makes use of efficient secret sharing schemes, a cryptographic primitive closely related to witness encryption and known to exist unconditionally (with statistical security) for access structures that are monotone formulae [27, 5] (see [4, 11]).\par




The second main result of [16] shows that ${\mathsf{MCSP}}^{*}$, a partial function version of ${\mathsf{MCSP}}$, is likewise ${\mathsf{NP}}$-hard. This proof requires even more sophisticated techniques in order to produce a truth-table on logarithmic-length inputs. In particular, Hirahara reduces from a gap version of the “collective minimum (weight) monotone satisfying assignment problem” (${\sf CMMSA}$), which generalizes ${\mathsf{MMSA}}$. Using the techniques of that proof and padding the inputs to have polynomial length, one can obtain ${\mathsf{NP}}$-hardness of learning ${\mathsf{P}}/{\mathsf{poly}}$, but with a smaller gap.\par



% SOURCE-ID Thmtheorem10


\par\medskip\noindent\textbf{
\textbf{Theorem 10}\textbf{ }(Implicit in [16])\textbf{.}
}\par



\emph{For any constant $\beta>0$, there exists a constant $\alpha>0$ such that for some polynomial $s$, it is ${\mathsf{NP}}$-hard under a randomized one-query reduction to agnostically PAC learn ${\mathsf{SIZE}}[s(n)]$ by ${\mathsf{SIZE}}[s(n)\cdot 2^{\alpha(\log n)^{1-\beta}}]$ over a flat ${\mathsf{P}}/{\mathsf{poly}}$-samplable distribution.}\par



% END-SOURCE-ID Thmtheorem10



However, those ideas do not extend to obtain the larger $s(n)$ vs. $s(n)\cdot n^{1/(\log\log n)^{O(1)}}$ gap for circuits rather than programs as in Corollary 9.\par



% SOURCE-ID Thmtheorem11


\par\medskip\noindent\textbf{
$\blacktriangleright$ Remark 11.
}\par



\emph{We note that Theorem 8 is in fact proved by a randomized half-Levin reduction from ${\mathsf{Gap}}{\mathsf{MMSA}}$ to ${\sf CGL}$! That is, there is a poly-time machine $R_{wit}$ such that, with high probability over $r$, if $\mathcal{E}_{r}$ is the output of the reduction of Theorem 8 run on a ${\mathsf{Gap}}{\mathsf{MMSA}}$ instance $\varphi$ with randomness $r$, and $\alpha$ is a witness for $\varphi$, then $R_{wit}(\varphi,\alpha,r)$ outputs a circuit certifying that $\mathcal{E}_{r}$ is a yes-instance of ${\sf CGL}$. However, unfortunately, we do not know how to extend the results of the previous section to obtain non-trivial witness encryption from a randomized half-Levin reduction, particularly since the amount of randomness required is greater than the parameter $s_{2}$ (which translates to the security of the witness encryption scheme). In any case, our results show that if the reduction of Theorem 8 can be suitably derandomized (with the running time of $R_{wit}$ less than the parameter $s_{2}$) then one would unconditionally obtain non-trivial witness encryption.}\par



% END-SOURCE-ID Thmtheorem11




\par\medskip\noindent\textbf{Unconditional NP-hardness of improper learning for other concept classes}\par




Though in the case of ${\mathsf{P}}/{\mathsf{poly}}$ we are unable to obtain ${\mathsf{NP}}$-hardness with a gap greater than $n^{1/(\log\log n)^{O(1)}}$, we show in this section that for some related concept classes, \emph{improper} learning is ${\mathsf{NP}}$-hard. We demonstrate such results for learning oracle circuits and learning RAM programs.\par




In the following, for a fixed $\lambda:\mathbb{N}\to\mathbb{N}$ and distribution family $D=\{D_{m}\}_{m\in\mathbb{N}}$, where each $D_{m}$ is supported over strings of length $2^{m}$ and samplable uniformly in time ${\mathsf{poly}}(2^{m})$, $D\text{-}{\sf oracle}\text{-}{\sf CGL}_{s_{1}}^{s_{2}}[\varepsilon,\lambda]$ is a generalization of ${\sf CGL}_{s_{1}}^{s_{2}}[\varepsilon]$ in which all circuits and samplers take access to an oracle $O\in\{0,1\}^{2^{\lambda(n)}}$ randomly sampled from $D_{\lambda(n)}$. $\mathcal{E}$ is a yes-instance of $D\text{-}{\sf oracle}\text{-}{\sf CGL}_{s_{1}}^{s_{2}}[\varepsilon,\lambda]$ if there exists a circuit $C$ of size $s_{1}(n)$ such that for every $O$ in the support of $D_{\lambda(n)}$,\par


\[\begin{gathered}
\Pr_{(x,b)\sim\mathcal{E}^{O}}[C^{O}(x)=b]=1,
\end{gathered}\]


and $\mathcal{E}$ is a no-instance of $D\text{-}{\sf oracle}\text{-}{\sf CGL}_{s_{1}}^{s_{2}}[\varepsilon,\lambda]$ if, with probability $1-\varepsilon(n)$ over $O\sim D_{\lambda(n)}$, for any circuit $C$ of size $s_{2}(n)$,\par


\[\begin{gathered}
\Pr_{(x,b)\sim\mathcal{E}^{O}}[C^{O}(x)=b]<\frac{1}{2}+\varepsilon(n).
\end{gathered}\]




We prove the following ${\mathsf{NP}}$-hardness of ${\sf oracle}\text{-}{\sf CGL}$, from which we will derive further corollaries.\par



% SOURCE-ID Thmtheorem12


\par\medskip\noindent\textbf{
\textbf{Theorem 12}\textbf{.}
}\par



\emph{There exist a distribution family $D\in{\mathsf{P}}{\mathsf{SAMP}}$, an ${\mathsf{NP}}$-complete language $L$, and a polynomial $s_{1}$ such that, for all polynomials $s_{2},\varepsilon^{-1}$, there is a function $\lambda:\mathbb{N}\to\mathbb{N}$ such that $\lambda(n)=O(\log n)$ for $n\in\mathbb{N}$ and $L$ reduces to $D\text{-}{\sf oracle}\text{-}{\sf CGL}^{s_{2}}_{s_{1}}[\varepsilon,\lambda]$ under an honest half-Levin reduction.}\par



% END-SOURCE-ID Thmtheorem12



Theorem 12 is proved using a framework developed by Huang, Ilango, and Ren for “witness encryption in oracle worlds” [21]. More specifically, drawing on the candidate witness encryption scheme proposed in [11], the authors show that there exist an ${\mathsf{NP}}$-complete language $L$ and a distribution family $D\in{\mathsf{P}}{\mathsf{SAMP}}$ such that, with high probability over oracles $O$ sampled randomly from $D$, witness encryption for $L$ exists with respect to $O$. That is, ${\mathsf{Enc}}$, ${\mathsf{Dec}}$, and adversaries all get oracle access to $O$.\par




We may then apply the techniques of one direction of our main equivalence, Theorem 4, to obtain ${\mathsf{NP}}$-hardness of learning with respect to $O$. Crucially, in [21], an oracle $O$ with truth-table length $n$ gives rise to a witness encryption scheme secure against non-uniform adversaries of size polynomially related to $n$. So, to achieve ${\mathsf{NP}}$-hardness of ${\sf oracle}\text{-}{\sf CGL}$ with an arbitrary polynomial gap, one only requires oracle truth-tables of polynomial length.\par




As a consequence of Theorem 12, we show that agnostic learning is ${\mathsf{NP}}$-hard in an “oracle-PAC” model that we introduce here, which generalizes standard agnostic PAC learning of ${\mathsf{P}}/{\mathsf{poly}}$. In this model, a learner for ${\mathsf{SIZE}}[s]$ on input length $n$ is explicitly given the whole truth table of an oracle function $\mathcal{O}:\{0,1\}^{O(\log n)}\to\{0,1\}$ sampled randomly from some distribution $D\in{\mathsf{P}}{\mathsf{SAMP}}$. The learner is asked to output an oracle circuit $h$ such that $h^{\mathcal{O}}$ approximates the target function almost as well as the best size-$s$ $\mathcal{O}$-oracle circuit; see Definition 23.\par



% SOURCE-ID Thmtheorem13


\par\medskip\noindent\textbf{
\textbf{Corollary 13}\textbf{.}
}\par



\emph{For every constant $c\in\mathbb{N}$ and sufficiently large polynomial $s$, it is ${\mathsf{NP}}$-hard under a randomized one-query reduction to agnostically oracle-PAC learn ${\mathsf{SIZE}}[s(n)]$ by ${\mathsf{SIZE}}[s(n)^{c}]$.}\par



% END-SOURCE-ID Thmtheorem13



We stress that the above holds under a standard polynomial-time reduction that does not itself require any oracle.\par




As a second consequence of the ${\mathsf{NP}}$-hardness of ${\sf oracle}\text{-}{\sf CGL}$, we show that it is ${\mathsf{NP}}$-hard to agnostically learn \emph{RAM programs} in an improper setting. In particular, we focus on a decision problem based on the problem ${\sf MINLT}$ defined in [31]. In Ko’s definition, the input to ${\sf MINLT}$ consists of pairs $(x_{i},b_{i})_{i\in[m]}\in(\{0,1\}^{n}\times\{0,1\})^{m}$ for some $m={\mathsf{poly}}(n)$ along with size and running time parameters $s,t\in\mathbb{N}$ represented in unary. The input is a yes-instance if and only if there exists a program $M$ of size at most $s$ and running time at most $t$ such that $M(x_{i})=b_{i}$ for all $i\in[m]$. Ko exhibited an oracle relative to which ${\sf MINLT}$ is not ${\mathsf{NP}}$-hard [31].\par




Hirahara, using non-relativizing techniques, showed that a certain generalization of ${\sf MINLT}$ (called ${\sf GapLearn}$) is in fact ${\mathsf{NP}}$-hard [16]. In this promise-problem, the pairs $(x_{i},b_{i})$ are represented implicitly by a sampling circuit (as in the definition of ${\sf CGL}$), and there is a gap, eg., in program size between yes-instances and no-instances. It is easy to see that ${\mathsf{NP}}$-hardness of this problem (Theorem 1.1 in [16]) implies ${\mathsf{NP}}$-hardness of Ko’s ${\sf MINLT}$.\par




We formulate our definition of ${\sf Gap}{\sf RAM}\text{-}{\sf MINLT}$ analogously to the definition of ${\sf GapLearn}$ in [16] but considering programs that are granted random access to their inputs. Namely, the input to ${\sf Gap}_{g,\varepsilon}{\sf RAM}\text{-}{\sf MINLT}$ is a sampler $\mathcal{E}$ over pairs $(x,b)$ along with parameters $s,t\in\mathbb{N}$. The input $(\mathcal{E},1^{s},1^{t})$ is a yes-instance if there exists a program of size at most $s$ and running in time at most $t$ such that $\Pr_{(x,b)\sim\mathcal{E}}[M^{x}=b]=1$, and $(\mathcal{E},1^{s},1^{t})$ is a no-instance if for any program of size at most $g(s)$ and running in time at most $g(t)$, $\Pr_{(x,b)\sim\mathcal{E}}[M^{x}=b]<\frac{1}{2}+\varepsilon(n)$. Here, $M^{x}$ denotes the program $M$ with random access to its input $x$.\par




We obtain the following ${\mathsf{NP}}$-hardness of ${\sf Gap}{\sf RAM}\text{-}{\sf MINLT}$. Note that the parameters we obtain imply ${\mathsf{NP}}$-hardness of learning RAM programs in an improper setting.\par



% SOURCE-ID Thmtheorem14


\par\medskip\noindent\textbf{
\textbf{Corollary 14}\textbf{.}
}\par



\emph{For all polynomials $g$ and $\varepsilon^{-1}$, ${\sf Gap}_{g,\varepsilon}{\sf RAM}\text{-}{\sf MINLT}$ is ${\mathsf{NP}}$-hard under a randomized many-one reduction.}\par



% END-SOURCE-ID Thmtheorem14



In contrast, [16] considers programs that do not require random access to their inputs; the gap in program size achieved there is small (namely, $s(n)$ vs. $s(n)\cdot n^{1/(\log\log n)^{O(1)}}$), though the no-instances extend to time-unbounded programs.\par





\par\medskip\noindent\textbf{Applications for the possibility of cryptography}\par




Finally, we prove a few results demonstrating what cryptography is possible assuming worst-case hardness of ${\mathsf{NP}}$ along with various assumptions concerning witness encryption and ${\mathsf{NP}}$-hardness of learning.\par




Assuming a randomized non-adaptive reduction from ${\mathsf{SAT}}$ to improper agnostic learning as described in [2] (see Definition 28), we show that the \emph{worst-case} assumption ${\mathsf{NP}}\nsubseteq{\mathsf{BPP}}$ is sufficient to imply one-way functions secure infinitely often against uniform PPT adversaries. This improves a main result of [2] (Theorem 1 in the introduction), which requires \emph{average-case} hardness of ${\mathsf{NP}}$. In terms of Impagliazzo’s five worlds [22], the result of [2] conditionally excludes Pessiland. In contrast, we conditionally exclude both Heuristica and Pessiland.\par



% SOURCE-ID Thmtheorem15


\par\medskip\noindent\textbf{
\textbf{Theorem 15}\textbf{.}
}\par



\emph{Suppose there is a randomized non-adaptive honest reduction $R$ and a polynomial $s:\mathbb{N}\to\mathbb{N}$ such that, for every constant $c\in\mathbb{N}$, $R$ reduces ${\mathsf{SAT}}$ to agnostically PAC learning ${\mathsf{SIZE}}[s(n)]$ by ${\mathsf{SIZE}}[s(n)^{c}]$. Then, unless ${\mathsf{NP}}\subseteq{\mathsf{BPP}}$, there exist infinitely-often one-way functions.}\par



% END-SOURCE-ID Thmtheorem15



As discussed in the introduction, Garg et al. show that WE for an ${\mathsf{NP}}$-complete language along with a one-way function is sufficient for public-key encryption [11]. Combining this result, our main equivalence in Theorem 4, and Theorem 15, we obtain that witness encryption simultaneously excludes Heuristica, Pessiland, and Minicrypt.\par



% SOURCE-ID Thmtheorem16


\par\medskip\noindent\textbf{
\textbf{Theorem 16}\textbf{.}
}\par



\emph{Suppose ${\mathsf{P}}/{\mathsf{poly}}$-secure witness encryption exists for ${\mathsf{NP}}$. Then, unless ${\mathsf{NP}}\subseteq{\mathsf{BPP}}$, there exists public-key encryption secure infinitely often against polynomial-time adversaries.}\par



% END-SOURCE-ID Thmtheorem16



We believe that this statement can also be proved by combining prior works [33, 19, 11] (see also [35]), though we have not seen the uniform version stated (ie. the version with the assumption ${\mathsf{NP}}\nsubseteq{\mathsf{BPP}}$ rather than ${\mathsf{NP}}\nsubseteq{\mathsf{P}}/{\mathsf{poly}}$). In any case, we provide an alternative proof that does not rely on properties of zero-knowledge, such as the “SZK/OWF” characterization from [42].\par




Theorem 16 and our our main equivalence in Theorem 4 together imply that if learning is ${\mathsf{NP}}$-hard under a half-Levin reduction to ${\sf CGL}_{s_{1}}^{s_{2}}[\varepsilon]$ working for arbitrary polynomials $s_{2}$ and $\varepsilon^{-1}$, one can rule out Minicrypt in addition to Heuristica and Pessiland. Thus, we obtain a stark distinction between the two kinds of reduction to learning defined in [2]. Namely, assuming worst-case hardness of ${\mathsf{NP}}$, a \emph{randomized non-adaptive} reduction to agnostic learning over ${\mathsf{P}}{\mathsf{SAMP}}/{\mathsf{poly}}$ yields \emph{one-way functions}, whereas a \emph{deterministic many-one} reduction to ${\sf CGL}$, provided it is \emph{half-Levin}, yields \emph{public-key encryption}.\par



% SOURCE-ID Thmtheorem17


\par\medskip\noindent\textbf{
\textbf{Theorem 17}\textbf{.}
}\par



\emph{Suppose there exist a polynomial $s_{1}$ and a pair of machines $(R,R_{wit})$ such that, for all polynomials $s_{2},\varepsilon^{-1}$, $(R,R_{wit})$ is an honest half-Levin reduction from ${\mathsf{SAT}}$ to ${\sf CGL}^{s_{2}}_{s_{1}}[\varepsilon]$. Then, unless ${\mathsf{NP}}\subseteq{\mathsf{BPP}}$, there exists public-key encryption secure infinitely often against polynomial-time adversaries.}\par



% END-SOURCE-ID Thmtheorem17





\par\medskip\noindent\textbf{
1.2 Related Work}\par




As mentioned above, a very recent work of Liu, Mazor, and Pass shows that witness encryption for ${\mathsf{NP}}$ is equivalent to “laconic special-honest
verifier zero-knowledge arguments” for ${\mathsf{NP}}$ [36]. Our main result shows that half-Levin reductions from ${\mathsf{NP}}$ to ${\sf CGL}$ are equivalent to both. This sharpens the result of Applebaum et al. showing that ${\mathsf{NP}}$-hardness of ${\sf CGL}$ under a many-one (not necessarily half-Levin) reduction implies ${\mathsf{NP}}\subseteq{\sf SZKA}$ [2].\par




Huang, Ilango, and Ren [21] used an assumption of the existence of witness encryption secure against subexponential-size circuits to conclude the ${\mathsf{NP}}$-hardness of approximating the time-bounded conditional Kolmogorov complexity ${\mathsf{K}}^{t}(x\mid y)$, in the regime where $t>|y|$, with essentially optimal gap parameters. This is similar to one direction of our Theorem 4 showing an equivalence between the existence of witness encryption (secure against ${\mathsf{P}}/{\mathsf{poly}}$) and ${\mathsf{NP}}$-hardness of a learning problem ${\sf CGL}$ (under restricted reductions). We note that in our framework, proving such a connection between witness encryption and ${\mathsf{NP}}$-hardness of learning is much simpler than in the setting considered by [21].\par




Our Theorem 12 relies on a framework for “witness encryption in oracle worlds” developed in [21] (building on ideas from [11]). The authors of [21] obtain ${\mathsf{NP}}$-hardness of the Minimum \emph{Oracle} Circuit Size Problem with a large circuit size gap [21]. This is somewhat like an “exponential regime” analogue of our oracle-PAC learning model, but we stress that, in [21], an algorithm for ${\mathsf{MOCSP}}$ needs to work in the worst case over oracle truth-tables. In our case, a learner only needs to work with high probability over oracle truth-tables sampled from some distribution in ${\mathsf{P}}{\mathsf{SAMP}}$.\par




An earlier example of an equivalence between ${\mathsf{NP}}$-hardness of a meta-complexity problem and the existence of a secure cryptographic primitive is the result by
Hirahara [17]. It shows an equivalence between the existence of one-way functions and ${\mathsf{NP}}$-hardness of distributional time-bounded Kolmogorov complexity $d{\mathsf{K}}^{t}$ under a certain kind of restricted reductions.\par




An earlier example of an equivalence between hardness of PAC learning and the existence of a cryptographic primitive is due to Nanashima [39]. In that work, Nanashima shows that auxiliary-input one-way functions exist if and only if a certain average-case, or “heuristic”, variant of PAC learning is hard for efficient algorithms. In contrast, we focus on \emph{${\mathsf{NP}}$-hardness} of worst-case agnostic PAC learning.\par





\par\medskip\noindent\textbf{
1.3 Overview of Techniques}\par




\par\medskip\noindent\textbf{Proof of Theorem 4
}\par




We start with an informal overview of the proof of our main equivalence, Theorem 4. In both directions, the constructions employed and the proofs of correctness are straightforward. We first show how to obtain a half-Levin reduction from $L$ to ${\sf CGL}$ from a ${\mathsf{P}}/{\mathsf{poly}}$-secure witness encryption scheme $({\mathsf{Enc}},{\mathsf{Dec}})$ for $L$. In particular, given an $L$-instance $z$, the reduction outputs a distribution $\mathcal{E}$ as follows.\par



Sample $b\sim\mathcal{U}$. Let $x:={\mathsf{Enc}}(z,b)$. Output $(x,b)$.\par



Note that we can assume that the output length of ${\mathsf{Enc}}$ is at least $n^{\Omega(1)}$ without loss of generality, so the reduction described here is honest.\par




In the case that $z\in L$ with witness $w$, $\mathcal{E}$ is a yes-instance of ${\sf CGL}$, since the circuit defined by ${\mathsf{Dec}}(-,z,w)$ recovers the bit $b$ from $x$ with probability 1 over $\mathcal{E}$. This is true by the correctness of $({\mathsf{Enc}},{\mathsf{Dec}})$. Also note that, since ${\mathsf{Dec}}$ is a uniform polynomial-time algorithm, the circuit ${\mathsf{Dec}}(-,z,w)$ can be constructed in polynomial time given $z$ and $w$. Thus, our reduction is half-Levin. In the case that $z\notin L$, we observe that if there were a circuit $C$ of polynomial size such that\par


\[\begin{gathered}
\Pr_{(x,b)\sim\mathcal{E}}[C(x)=b]\geq\frac{1}{2}+\frac{1}{{\mathsf{poly}}(|x|)},
\end{gathered}\]


then by definition of $\mathcal{E}$,\par


\[\begin{gathered}
\Pr[C({\mathsf{Enc}}(z,1))=1]-\Pr[C({\mathsf{Enc}}(z,0))=1]\geq\frac{1}{{\mathsf{poly}}(n)},
\end{gathered}\]


violating the security of ${\mathsf{Enc}}$. So, $\mathcal{E}$ must be a no-instance of ${\sf CGL}$.\par




Now we show how to obtain a ${\mathsf{P}}/{\mathsf{poly}}$-secure witness encryption scheme from a half-Levin reduction $(R,R_{wit})$ to ${\sf CGL}$. We define a witness encryption scheme as follows. Let $z$ be an instance of $L$, and if $z\in L$, let $w$ be a witness. Let $a$ be a secret bit to encrypt.\par



${\mathsf{Enc}}(z,a)$ computes $\mathcal{E}=R(z)$, samples $(x,b)\sim\mathcal{E}$ and then outputs $(x,a\oplus b)$.\par




${\mathsf{Dec}}((x,c),z,w)$ computes $A:=R_{wit}(z,w)$ and then outputs $A(x)\oplus c$.\par



The correctness of $({\mathsf{Enc}},{\mathsf{Dec}})$ follows from the fact that, for $z\in L$, $R_{wit}(z,w)$ is a witness for $\mathcal{E}=R(z)$ being a yes-instance of ${\sf CGL}$: that is, a circuit $A$ of polynomial size such that\par


\[\begin{gathered}
\Pr_{(x,b)\sim\mathcal{E}}[A(x)=b]=1.
\end{gathered}\]


It follows that $A(x)\oplus c=b\oplus(a\oplus b)=a$. The security of ${\mathsf{Enc}}$ follows from Yao’s argument that a distinguisher implies a next-bit predictor. Specifically, for $(x,b)\sim\mathcal{E}$, from a circuit of size $s(n)$ distinguishing ${\mathsf{Enc}}(z,0)=(x,b)$ from ${\mathsf{Enc}}(z,1)=(x,1\oplus b)$, we would obtain a circuit of size $O(s(n))$ predicting (with inverse polynomial advantage) $b$ from $x$. By the honesty of the reduction, we have $O(s(n))={\mathsf{poly}}(|x|)$. This violates the fact that $\mathcal{E}$ is a no-instance of ${\sf CGL}_{s_{1}}^{s_{2}}$ for every polynomial $s_{2}$, a contradiction.\par





\par\medskip\noindent\textbf{Proof of Theorem 6
}\par




To modify the argument above for the case of ${\mathsf{BPP}}$-secure witness encryption, we note that, in the first direction discussed, the security of WE is only violated if the circuit $C$ can be constructed uniformly in polynomial time. Therefore, we consider ${\mathsf{BPP}}$-black-box reductions to \emph{search}-${\sf CGL}$. That is, roughly, the reduction is only guaranteed to work when its oracle is a PPT algorithm, and we require the oracle to actually construct a small circuit with inverse polynomial advantage over the queried distribution $\mathcal{E}$ if one exists. For the other direction, we employ the observation that, from a uniform PPT distinguisher algorithm $A$, one can construct a next-bit predictor $P^{A}$ uniformly.\par





\par\medskip\noindent\textbf{Proof of Theorem 7
}\par




For the case of ${\mathsf{GapDistMCSP}}$ in Theorem 7, a challenge is that an arbitrary witness encryption scheme $({\mathsf{Enc}},{\mathsf{Dec}})$ may not yield a distribution $\mathcal{E}$ that represents a well-defined function. On one hand, if $z\in L$, then there is guaranteed to be a unique bit $b_{x}$ such that $(x,b_{x})$ is in the support of $\mathcal{E}$, by the correctness of the witness encryption scheme. However, if $z\notin L$, this may not be true. Luckily, since $\mathcal{E}$ is supported over logarithmic-length strings, we can check in polynomial time whether each $x$ has a unique label. If not, we use our assumption of a circuit lower bound for ${\mathsf{E}}$ to output a no-instance of ${\mathsf{GapDistMCSP}}$. In the other direction, from a (deterministic) half-Levin reduction from ${\mathsf{SAT}}$ to ${\mathsf{GapDistMCSP}}$, we obtain a circuit lower bound for ${\mathsf{E}}$ by applying the reduction to a trivial no-instance of ${\mathsf{SAT}}$, as in [29].\par





\par\medskip\noindent\textbf{Proof of Theorem 8 and corollary}\par




Our proof of ${\mathsf{NP}}$-hardness of ${\sf CGL}$ with parameters corresponding to semi-proper learning works within the framework developed in [16]. It differs from Theorem 7.2 in that paper mainly in the “completeness” argument. Referring to the notation in the proof of that theorem and of our Lemma 64, we need to construct a circuit of size not exceeding $O(\theta\lambda)$.\par




For one thing, in the proof of [16], using ${\mathsf{DP}}_{k}$ generators to encode the high-complexity truth-tables $f_{i}$, the number of input gates to the circuit alone would exceed the desired threshold. Therefore, we use Nisan-Wigderson generators rather than ${\mathsf{DP}}_{k}$ generators to allow for a short random seed. (A similar idea is used elsewhere in [16].)\par




More importantly, we apply a recent result of Holmgren and Rothblum showing that “multiselection” of $t$ indices from a string of length $m\in\mathbb{N}$ can be done by a circuit of size at most $O(m+t\cdot\log^{3}(m))$ [20]. See Lemma 61. So, we can make ${\mathsf{poly}}(n)$ queries to $\theta$ truth-tables $f_{j}$, each of length $\lambda$, with total circuit size at most $O(\theta\cdot(\lambda+o(\lambda)))$. This allows us to compute values of the amplified functions $\hat{f}_{j}$ and recover the secret labels $b$ without exceeding the desired size threshold. Note that Uhlig’s Theorem as in [16] cannot accommodate ${\mathsf{poly}}(n)$ queries.\par




To obtain Corollary 9, which states that semi-proper agnostic learning is ${\mathsf{NP}}$-hard, we give a simple randomized reduction from ${\sf CGL}$ to learning that works by empirical estimation of the success probability of the hypothesis produced by the learner. See Lemma 33.\par





\par\medskip\noindent\textbf{Proof of Theorem 12 and corollaries}\par




To obtain ${\mathsf{NP}}$-hardness of ${\sf oracle}\text{-}{\sf CGL}$, we apply an argument analogous to the first direction of our main equivalence, along with the fact that, for some distribution family $D\in{\mathsf{P}}{\mathsf{SAMP}}$, witness encryption for ${\mathsf{NP}}$ exists unconditionally (with high probability) with respect to oracles sampled from $D$. This was demonstrated in [21]. More specifically, the authors show the following.\par



% SOURCE-ID Thmtheorem18


\par\medskip\noindent\textbf{
\textbf{Lemma 18}\textbf{ }([21] – informal; see Lemma 41 and definition 40)\textbf{.}
}\par



\emph{For some ${\mathsf{NP}}$-complete language $L$, distribution $D\in{\mathsf{P}}{\mathsf{SAMP}}$, and constant $\ell\in\mathbb{N}$, there is a witness encryption scheme $({\mathsf{Enc}},{\mathsf{Dec}})$ for $L$ such that, for all $m\in\mathbb{N}$, with high probability over $O\sim D_{m}$ (with $O\in\{0,1\}^{m}$), $({\mathsf{Enc}}^{O},{\mathsf{Dec}}^{O})$ is secure against $O$-oracle adversaries of size $m^{1/\ell}$.}\par



% END-SOURCE-ID Thmtheorem18



Our idea is to apply the foregoing “oracle witness encryption” along with our main technique for proving that witness encryption implies ${\mathsf{NP}}$-hardness of ${\sf CGL}$. We modify the definition of ${\sf CGL}$ accordingly so that all circuits and samplers have access to an oracle randomly sampled from $D$. See Definition 25. Crucially, to achieve ${\mathsf{NP}}$-hardness of $D\text{-}{\sf oracle}\text{-}{\sf CGL}_{s(n)}^{s(n)^{c}}$ for a constant $c\in\mathbb{N}$, we require witness encryption with security against adversaries of size $s(n)^{c}$, and thus we require an oracle $O$ with truth-table length only $s(n)^{c\cdot\ell}={\mathsf{poly}}(n)$.\par




Corollary 13 follows from Theorem 12 analogously to Corollary 9 and Theorem 8, as discussed above: a straightforward reduction from ${\sf oracle}\text{-}{\sf CGL}$ to oracle-PAC learning.\par




To prove Corollary 14, we give a randomized reduction from $D\text{-}{\sf oracle}\text{-}{\sf CGL}$ to the problem ${\sf Gap}{\sf RAM}\text{-}{\sf MINLT}$. Given an instance $\mathcal{E}^{(-)}$ of $D\text{-}{\sf oracle}\text{-}{\sf CGL}$, the reduction starts by sampling a random truth-table $O\sim D$. Then, it defines an instance $\mathcal{E}^{\prime}$ of ${\sf Gap}{\sf RAM}\text{-}{\sf MINLT}$ which always outputs the entire truth-table $O$ along with a string $(x,b)$ sampled according to $\mathcal{E}^{O}$; one can think of $\mathcal{E}^{\prime}$ as sampling pairs $(x^{\prime},b)$ with $x^{\prime}=(x,O)$. Clearly, this results in a non-oracle instance of ${\sf Gap}{\sf RAM}\text{-}{\sf MINLT}$. In the proof of correctness, we also rely on the fact that a program of size and running time at most $s^{\prime}$ can be simulated by a circuit of size ${\mathsf{poly}}(s^{\prime})$ (and vice versa). Since the oracles obtained from [21] require length polynomially greater than the size of adversary they are secure against, we need to consider machines running in time less than the length of their inputs. Given these observations and the definitions of ${\sf oracle}\text{-}{\sf CGL}$ and ${\sf Gap}{\sf RAM}\text{-}{\sf MINLT}$, the corollary follows easily.\par





\par\medskip\noindent\textbf{Proof of Theorem 15
}\par




We prove Theorem 15 with an argument from [2] as a starting point.\footnote{Note that we include a proof below for the reader’s convenience; see Lemma 72.} Let $R$ be a randomized non-adaptive reduction from an ${\mathsf{NP}}$-complete language $L$ to agnostic learning over ${\mathsf{P}}{\mathsf{SAMP}}/{\mathsf{poly}}$. Assume that i.o. one-way functions do not exist. Then, for any $D\in{\mathsf{P}}{\mathsf{SAMP}}$ and auxiliary-input function $\{f_{z}\}_{z\in\{0,1\}^{*}}$, there is a machine $I$ that distributionally inverts $f_{z}$ with high probability over auxiliary inputs $z\sim D$. We think of instances $z$ of $L$ as auxiliary inputs for the function $f_{z}$ that, roughly, simulates $R(z)$ to obtain a sampling circuit $\mathcal{E}$ over string/label pairs and then samples from $\mathcal{E}$. Applebaum et al. show that there is a polynomial-time-computable hypothesis $h$, using a distributional inverter for $f_{z}$, having near-maximum agreement over $\mathcal{E}$. Running the reduction $R$ and answering oracle queries with hypotheses $h$ defined in this way, we obtain\par


\[\begin{gathered}
{\mathsf{Dist}}{\mathsf{NP}}\subseteq{\mathsf{Heur}}{\mathsf{BPP}},
\end{gathered}\]


where “${\mathsf{Heur}}$”, informally, refers to average-case heuristics that may err by incorrectly outputting 1 or 0.\par




We next use the fact that the inversion of polynomial-time-computable functions is \emph{verifiable}. That is, if a machine fails to invert on some instance, it can be defined to “realize this” in polynomial time. More concretely, the machine $I$ above can be defined to output $\bot$ on input $z$ if it fails to invert often on some particular $f_{z}$, using empirical estimation. From this observation, we further obtain\par


\[\begin{gathered}
{\mathsf{Dist}}{\mathsf{NP}}\subseteq{\mathsf{Avg}}{\mathsf{BPP}},
\end{gathered}\]


where ${\mathsf{Avg}}$ refers to heuristics that do not err but may output $\bot$ rather than 1 or 0 with inverse polynomial probability. Lastly, we apply a recent result of [12], which, from ${\mathsf{Dist}}{\mathsf{NP}}\subseteq{\mathsf{Avg}}{\mathsf{BPP}}$, yields \emph{worst-case} easiness of agnostic PAC learning over distributions in ${\mathsf{P}}{\mathsf{SAMP}}/{\mathsf{poly}}$. We then run the reduction $R$ \emph{for a second time}, answering oracle queries with the worst-case agnostic learner guaranteed by [12], to obtain ${\mathsf{NP}}\subseteq{\mathsf{BPP}}$.\par





\par\medskip\noindent\textbf{Proofs of Theorem 16 and Theorem 17
}\par




To obtain Theorem 16, start by assuming the existence of witness encryption for an ${\mathsf{NP}}$-complete language. Then, apply the characterization given by Theorem 4 to obtain a half-Levin reduction from ${\mathsf{NP}}$ to ${\sf CGL}$. Since ${\sf CGL}$ reduces to agnostic PAC learning, we may apply Theorem 15 as described above to exclude Heuristica and Pessiland. To summarize, informally,\par


\[\begin{gathered}
WE \displaystyle\implies{\mathsf{NP}}\leq_{hl}{\sf CGL} \\
\displaystyle\implies{\mathsf{NP}}\leq_{tt}\text{ agnostic PAC-learning over }{\mathsf{P}}{\mathsf{SAMP}}/{\mathsf{poly}} \\
\displaystyle\implies({\mathsf{NP}}\nsubseteq{\mathsf{BPP}}\implies\text{ i.o. OWF}),
\end{gathered}\]


where $\leq_{hl}$ refers to a half-Levin reduction and $\leq_{tt}$ refers to a randomized non-adaptive (ie. “truth-table”) reduction. Furthermore, Garg et al. show that witness encryption for ${\mathsf{NP}}$, together with the existence of a one-way function, yields public-key encryption [11]. This, together with the above, conditionally yields public-key encryption. To summarize,\par


\[\begin{gathered}
\displaystyle\text{WE and }{\mathsf{NP}}\nsubseteq{\mathsf{BPP}} \displaystyle\implies\text{WE and i.o. OWF } \\
\displaystyle\implies\text{i.o. PKE}
\end{gathered}\]




Theorem 17 is immediate from the reasoning above together with another application of the characterization in Theorem 4. That is,\par


\[\begin{gathered}
\displaystyle{\mathsf{NP}}\leq_{hl}{\sf CGL} \displaystyle\implies\text{WE } \\
\displaystyle\implies({\mathsf{NP}}\nsubseteq{\mathsf{BPP}}\implies\text{ i.o. PKE})
\end{gathered}\]







\par\medskip\noindent\textbf{
2 Preliminaries}\par



% SOURCE-ID Thmtheorem19


\par\medskip\noindent\textbf{
\textbf{Definition 19}\textbf{.}
}\par



\emph{We say that a distribution is flat if it is equal to the uniform distribution over its support.}\par



% END-SOURCE-ID Thmtheorem19


% SOURCE-ID Thmtheorem20


\par\medskip\noindent\textbf{
\textbf{Definition 20}\textbf{.}
}\par



\emph{For a distribution family $D=\{D_{n}\}_{n\in\mathbb{N}}$, $s:\mathbb{N}\to\mathbb{N}$, and $\delta:\mathbb{N}\to[0,1]$, ${\mathsf{SIZE}}_{D}^{\delta}[s]$ denotes the class of function families $f=\{f_{n}:\{0,1\}^{n}\to\{0,1\}\}_{n\in\mathbb{N}}$ such that, for all sufficiently large $n\in\mathbb{N}$, there exists a circuit $C_{n}$ of size at most $s(n)$ such that}\par


\[\begin{gathered}
\Pr_{x\sim D_{n}}[C_{n}(x)=f_{n}(x)]\geq\delta(n).
\end{gathered}\]



% END-SOURCE-ID Thmtheorem20



\par\medskip\noindent\textbf{
2.1 PAC Learning Problems}\par



% SOURCE-ID Thmtheorem21


\par\medskip\noindent\textbf{
\textbf{Definition 21}\textbf{ }(Optimal hypotheses)\textbf{.}
}\par



\emph{For $n\in\mathbb{N}$, a function $s:\mathbb{N}\to\mathbb{N}$, a distribution family $\mathcal{E}_{n}$ supported over $\{0,1\}^{n}\times\{0,1\}$, a parameter $\varepsilon>0$, and a hypothesis $h:\{0,1\}^{n}\to\{0,1\}$, we say that $h$ is $(s,\varepsilon)$-optimal for $\mathcal{E}_{n}$ if}\par


\[\begin{gathered}
\Pr_{(x,b)\sim\mathcal{E}_{n}}[h(x)=b]\geq\max_{f\in{\mathsf{SIZE}}[s(n)]}\left\{\Pr_{(x,b)\sim\mathcal{E}_{n}}[f(x)=b]\right\}-\varepsilon.
\end{gathered}\]




\emph{For $\lambda\in\mathbb{N}$, an oracle function $O:\{0,1\}^{\lambda}\to\{0,1\}$, and a hypothesis $g^{(-)}:\{0,1\}^{n}\to\{0,1\}$, we say that $g$ is $(s,\varepsilon,O)$-optimal for $\mathcal{E}_{n}$ if}\par


\[\begin{gathered}
\Pr_{(x,b)\sim\mathcal{E}_{n}}[g^{O}(x)=b]\geq\max_{f\in{\mathsf{SIZE}}[s(n)]}\left\{\Pr_{(x,b)\sim\mathcal{E}_{n}}[f^{O}(x)=b]\right\}-\varepsilon.
\end{gathered}\]



% END-SOURCE-ID Thmtheorem21


% SOURCE-ID Thmtheorem22


\par\medskip\noindent\textbf{
\textbf{Definition 22}\textbf{ }(Agnostic PAC learning circuits)\textbf{.}
}\par



\emph{Consider $s_{1},s_{2}:\mathbb{N}\to\mathbb{N}$. We say that an algorithm $A$ agnostically PAC learns ${\mathsf{SIZE}}[s_{1}(n)]$ by ${\mathsf{SIZE}}[s_{2}(n)]$ if, for every $n\in\mathbb{N}$, $\delta>0$, $\varepsilon>0$, and joint distribution family $\mathcal{E}=\{\mathcal{E}_{n}\}_{n\in\mathbb{N}}$ where each $\mathcal{E}_{n}$ is supported over $\{0,1\}^{n}\times\{0,1\}$, the following holds. With probability at least $1-\delta$, $A^{\mathcal{E}_{n}}$ outputs an $(s_{1},\varepsilon)$-optimal hypothesis $h$ for $\mathcal{E}_{n}$ of size at most $s_{2}(n)$.}\par




$A$\emph{ is given the parameters $n,1/\varepsilon$, and $1/\delta$ in unary.}\par




\emph{We say that $A$ agnostically PAC learns ${\mathsf{SIZE}}[s_{1}]$ in the improper setting if the above is true except there is no restriction on the complexity of $h$ (except that it cannot exceed size the run time of $A$).}\par



% END-SOURCE-ID Thmtheorem22


% SOURCE-ID Thmtheorem23


\par\medskip\noindent\textbf{
\textbf{Definition 23}\textbf{ }(Agnostic Oracle-PAC learning circuits)\textbf{.}
}\par



\emph{The definition is analogous to Definition 22, except we also consider distribution families $D=\{D_{\lambda}\}_{\lambda\in\mathbb{N}}$ for sampling oracle truth-tables, where $D_{\lambda}$ is supported over $\{0,1\}^{2^{\lambda}}$ and uniformly samplable in time ${\mathsf{poly}}(2^{\lambda})$. We say that an algorithm $A$ agnostically oracle-PAC learns ${\mathsf{SIZE}}[s_{1}]$ by ${\mathsf{SIZE}}[s_{2}]$ if, for any such $D$, for every $\varepsilon,\delta>0$ and $n,\lambda\in\mathbb{N}$, the following holds with probability at least $1-\delta$ over $O\sim D_{\lambda}$.}\par




\emph{For any joint distribution family $\mathcal{E}=\{\mathcal{E}_{n}\}_{n\in\mathbb{N}}$ where each $\mathcal{E}_{n}$ is supported over $\{0,1\}^{n}\times\{0,1\}$, with probability at least $1-\delta$ over the randomness of $A$, $A^{\mathcal{E}_{n}}({\sf tt}(O))$ outputs an $(s_{1},\varepsilon,O)$-optimal hypothesis $h$ for $\mathcal{E}_{n}$ of circuit size at most $s_{2}(n)$.}\par




$A$\emph{ is also given the parameters $n,1/\varepsilon$, and $1/\delta$ in unary.}\par



% END-SOURCE-ID Thmtheorem23



Applebaum et al. define the following decision version of the PAC learning problem for circuits.\par



% SOURCE-ID Thmtheorem24


\par\medskip\noindent\textbf{
\textbf{Definition 24}\textbf{ }(Computational gap-learning problem (CGL) [2])\textbf{.}
}\par



\emph{The computational gap-learning problem with size parameters $s_{1},s_{2}:\mathbb{N}\to\mathbb{N}$ and security parameter $\varepsilon:\mathbb{N}\to[0,1]$, denoted ${\sf CGL}_{s_{1}}^{s_{2}}[\varepsilon]$, is the promise problem defined as follows.\footnote{We note that [2] define ${\sf CGL}$ for a given completeness parameter $\alpha:\mathbb{N}\to[0,1]$ (ie. a yes-instance has $C$ such that $\Pr[C(x)=b]\geq\alpha(n)$). In this work, we only consider $\alpha=1$.} The input consists of a distribution $\mathcal{E}$ over pairs $(x,b)\in\{0,1\}^{n}\times\{0,1\}$ represented as a circuit of size ${\mathsf{poly}}(n)$.\footnote{In this work, we abuse notation by taking $\mathcal{E}$ to refer both to the distribution and the circuit sampling it.}}\par


\begin{itemize}


\item[] 
$\blacksquare$ 


$\mathcal{E}$\emph{ is a yes-instance if there exists a circuit }$C$\emph{ of size at most }$s_{1}(n)$\emph{ such that}\par


\[\begin{gathered}
\Pr_{(x,b)\sim\mathcal{E}}\left[C(x)=b\right]=1.
\end{gathered}\]





\item[] 
$\blacksquare$ 


$\mathcal{E}$\emph{ is a no-instance if for every circuit }$C$\emph{ of size at most }$s_{2}(n)$\emph{,}\par


\[\begin{gathered}
\Pr_{(x,b)\sim\mathcal{E}}\left[C(x)=b\right]<\frac{1}{2}+\varepsilon(n).
\end{gathered}\]





\end{itemize}



% END-SOURCE-ID Thmtheorem24


% SOURCE-ID Thmtheorem25


\par\medskip\noindent\textbf{
\textbf{Definition 25}\textbf{ }(Oracle computational gap-learning problem (${\sf oracle}\text{-}{\sf CGL}$))\textbf{.}
}\par



\emph{Let $D=\{D_{\lambda}\}_{\lambda\in\mathbb{N}}$ be a polynomial-time-samplable distribution family such that each $D_{\lambda}$ is supported over truth-tables of oracle functions $O_{\lambda}:\{0,1\}^{\lambda}\to\{0,1\}$. The $D$-oracle computational gap-learning problem with size parameters $s_{1},s_{2}:\mathbb{N}\to\mathbb{N}$, security parameter $\varepsilon:\mathbb{N}\to[0,1]$, and oracle parameter $\lambda:\mathbb{N}\to\mathbb{N}$, denoted $D\text{-}{\sf oracle}\text{-}{\sf CGL}_{s_{1}}^{s_{2}}[\varepsilon,\lambda]$, is defined analogously to ${\sf CGL}^{s_{2}}_{s_{1}}[\varepsilon]$ except all circuits (including the sampler $\mathcal{E}$) have access to an oracle randomly drawn from $D_{\lambda(n)}$. More specifically,}\par


\begin{itemize}


\item[] 
$\blacksquare$ 


$\mathcal{E}^{(-)}$\emph{ is a yes-instance if there exists an oracle circuit }$C^{(-)}$\emph{ of size at most }$s_{1}(n)$\emph{ such that, for every oracle }$O$\emph{ in the support of }$D_{\lambda(n)}$\emph{,}\par


\[\begin{gathered}
\Pr_{(x,b)\sim\mathcal{E}^{O}}\left[C^{O}(x)=b\right]=1.
\end{gathered}\]





\item[] 
$\blacksquare$ 


$\mathcal{E}^{(-)}$\emph{ is a no-instance if, with probability at least }$1-\varepsilon(n)$\emph{ over }$O\sim D_{\lambda(n)}$\emph{, for every oracle circuit }$C^{(-)}$\emph{ of size at most }$s_{2}(n)$\emph{,}\par


\[\begin{gathered}
\Pr_{(x,b)\sim\mathcal{E}^{O}}\left[C^{O}(x)=b\right]<\frac{1}{2}+\varepsilon(n).
\end{gathered}\]





\end{itemize}



% END-SOURCE-ID Thmtheorem25


% SOURCE-ID Thmtheorem26


\par\medskip\noindent\textbf{
\textbf{Definition 26}\textbf{ }(${\sf Gap}_{g,\varepsilon}{\sf RAM}\text{-}{\sf MINLT}$)\textbf{.}
}\par



${\sf Gap}{\sf RAM}\text{-}{\sf MINLT}$\emph{ with gap parameter $g:\mathbb{N}\to\mathbb{N}$ and soundness parameter $\varepsilon^{-1}:\mathbb{N}\to\mathbb{N}$, denoted ${\sf Gap}_{g,\varepsilon}{\sf RAM}\text{-}{\sf MINLT}$, is the decision problem defined as follows. The input consists of a distribution $\mathcal{E}$ over pairs $(x,b)\in\{0,1\}^{n}\times\{0,1\}$ represented as a circuit of size ${\mathsf{poly}}(n)$. The input also includes a threshold parameter $s\in\mathbb{N}$ and a running time parameter $t\in\mathbb{N}$ represented in unary.}\par


\begin{itemize}


\item[] 
$\blacksquare$ 


$(\mathcal{E},1^{s},1^{t})$\emph{ is a yes-instance of }${\sf Gap}_{g,\varepsilon}{\sf RAM}\text{-}{\sf MINLT}$\emph{ if there exists a }$t$\emph{-time program }$M$\emph{ of size at most }$s$\emph{ such that }$M^{x}=b$\emph{ for all }$(x,b)$\emph{ in the support of }$\mathcal{E}$\emph{.}\par





\item[] 
$\blacksquare$ 


$(\mathcal{E},1^{s},1^{t})$\emph{ is a no-instance of }${\sf Gap}_{g,\varepsilon}{\sf RAM}\text{-}{\sf MINLT}$\emph{ if for every }$g(t)$\emph{-time program }$M$\emph{ of size at most }$g(s)$\emph{,}\par


\[\begin{gathered}
\Pr_{(x,b)\sim\mathcal{E}}[M^{x}=b]<\frac{1}{2}+\varepsilon(n)
\end{gathered}\]





\end{itemize}



% END-SOURCE-ID Thmtheorem26




\par\medskip\noindent\textbf{
2.2 Reductions to Learning}\par




Here we define a form of reduction that is intermediate between the standard notions of a Karp reduction and a Levin reduction.\par



% SOURCE-ID Thmtheorem27


\par\medskip\noindent\textbf{
\textbf{Definition 27}\textbf{ }(Half-Levin reduction)\textbf{.}
}\par



\emph{Consider languages $L_{1},L_{2}\in{\mathsf{NP}}$ with witness relations $V_{L_{1}},V_{L_{2}}$. A half-Levin reduction from $L_{1}$ to $L_{2}$ consists of polynomial-time machines $R$ and $R_{wit}$ with the following properties.}\par


\begin{itemize}


\item[] 
$\blacksquare$ 


\emph{For any }$L_{1}$\emph{-instance }$x$\emph{, }$x\in L_{1}\iff R(x)\in L_{2}$\emph{;}\par





\item[] 
$\blacksquare$ 


\emph{For any }$(x,w)\in V_{L_{1}}$\emph{, it holds that }$(R(x),R_{wit}(x,w))\in V_{L_{2}}$\emph{.}\par





\end{itemize}



% END-SOURCE-ID Thmtheorem27



A number of the results in this paper involve half-Levin reductions to ${\sf CGL}$, which belongs to (a promise version of) ${\mathsf{MA}}$ but not necessarily ${\mathsf{NP}}$. Note that the definition of a half-Levin reduction can be extended to ${\sf CGL}$ in the natural way. Specifically, in a half-Levin reduction from $L\in{\mathsf{NP}}$ to ${\sf CGL}_{s_{1}}^{s_{2}}[\varepsilon]$, the machine $R_{wit}$, given some $z\in L$ and witness $w$, is required to produce a circuit $C$ of size at most $s_{1}$ such that\par


\[\begin{gathered}
\Pr_{(x,b)\sim\mathcal{E}}[C(x)=b]=1,
\end{gathered}\]


where $\mathcal{E}$ is the output of $R(z)$. In addition, we say that a reduction to ${\sf CGL}$ is \emph{honest} if, on an input $z$ of length $n$, it outputs a distribution $\mathcal{E}$ supported over $\{0,1\}^{n^{\Omega(1)}}\times\{0,1\}$.\par



% SOURCE-ID Thmtheorem28


\par\medskip\noindent\textbf{
\textbf{Definition 28}\textbf{ }(Reduction to agnostic PAC learning circuits [2])\textbf{.}
}\par



\emph{For polynomials $s_{1}$ and $s_{2}$ and $\delta:\mathbb{N}\to[0,1]$, a reduction from a language $L$ to agnostic PAC learning ${\mathsf{SIZE}}[s_{1}]$ by ${\mathsf{SIZE}}[s_{2}]$ over ${\mathsf{P}}{\mathsf{SAMP}}/{\mathsf{poly}}$ with failure probability $\delta$ consists of a PPT machine $R$, which on input $z\in\{0,1\}^{n}$, makes queries that are joint distributions $\mathcal{E}^{(i)}$ over $\{0,1\}^{m(n)}\times\{0,1\}$ for some $m:\mathbb{N}\to\mathbb{N}$, each represented by a sampling circuit. The oracle provides a hypothesis $h_{i}$, represented as a ${\mathsf{poly}}(n)$-size circuit, as a response to each query. Based on the hypotheses $h_{i}$ and its own computation, $R$ accepts or rejects the input $z$.}\par




\emph{It is guaranteed that, for some inverse polynomial $\varepsilon$, if the oracle returns an $(s_{1},\varepsilon)$-optimal hypothesis $h$ of size $s_{2}(m(n))$ for each query $\mathcal{E}^{(i)}$, then with probability at least $1-\delta(n)$, $R(z)=L(z)$.}\par




\emph{If the failure probability is omitted, we take it to be $1/3$.}\par




\emph{If $m(n)\geq n^{\Omega(1)}$, we call the reduction honest.}\par



% END-SOURCE-ID Thmtheorem28


% SOURCE-ID Thmtheorem29


\par\medskip\noindent\textbf{
\textbf{Definition 29}\textbf{ }(Reduction to agnostic oracle-PAC learning circuits)\textbf{.}
}\par



\emph{We define reductions to agnostic oracle-PAC learning analogously with Definition 28. Specifically, the reduction samples random $O^{(i)}\sim D_{\lambda(n)}$ and then queries pairs $(\mathcal{E}^{(i)},{\sf tt}(O^{(i)}))$ to its oracle. The guarantee of the reduction holds if the oracle responds with an $(s_{1},\varepsilon,O^{(i)})$-optimal hypothesis of size at most $s_{2}(m(n))$ for each query $\mathcal{E}^{(i)}$.}\par



% END-SOURCE-ID Thmtheorem29



For our results concerning ${\mathsf{BPP}}$-secure witness encryption, we will require non-black-box notions of reduction. We start by recalling the standard definition of a class-specific black-box reduction.\par



% SOURCE-ID Thmtheorem30


\par\medskip\noindent\textbf{
\textbf{Definition 30}\textbf{ }(Class-specific black-box reduction [13])\textbf{.}
}\par



\emph{For a complexity class $C$ and languages $L_{1}$ and $L_{2}$, a PPT oracle machine $R$ is a $C$-black-box reduction from $L_{1}$ to $L_{2}$ if, for any oracle $\mathcal{O}\in C$ deciding $L_{2}$, $R^{\mathcal{O}}$ decides $L_{1}$.}\par



% END-SOURCE-ID Thmtheorem30



We will actually require a slightly less strict kind of reduction: rather than insisting $L_{2}\in C$ everywhere, to be correct on a given input $z$, the reduction will only require that an oracle $\mathcal{O}\in C$ agree with $L_{2}$ on the queries actually made given that input $z$. For our purposes, we only need to consider the case of a reduction that generates its queries deterministically.\par



% SOURCE-ID Thmtheorem31


\par\medskip\noindent\textbf{
\textbf{Definition 31}\textbf{ }(Instance-wise class-specific black-box reduction)\textbf{.}
}\par



\emph{Consider a complexity class $C$, languages $L_{1}$ and $L_{2}$, and a non-adaptive oracle machine $R$ that generates its queries deterministically but may use randomness elsewhere. $R$ is an instance-wise $C$-black-box reduction from $L_{1}$ to $L_{2}$ if the following holds. For any oracle $\mathcal{O}\in C$ and instance $z$ of $L_{1}$, if $q_{1},...,q_{m}$ are the oracle queries made by $R$ on input $z$ and $\mathcal{O}(q_{i})=L_{2}(q_{i})$ for all $i\in[m]$, then $R^{\mathcal{O}}(z)$ correctly decides $L_{1}(z)$.}\par



% END-SOURCE-ID Thmtheorem31



We now explicitly define our required notion of an instance-wise ${\mathsf{BPP}}$-black-box half-Levin reduction to the search version of ${\sf CGL}$. Note that we may drop the qualifier “instance-wise” elsewhere in the presentation. Also note that the definition is extended from those above to the case of a reduction to a search problem.\par



% SOURCE-ID Thmtheorem32


\par\medskip\noindent\textbf{
\textbf{Definition 32}\textbf{ }(Instance-wise ${\mathsf{BPP}}$-black-box half-Levin reduction to ${\sf search}\text{-}{\sf CGL}_{s_{1}}^{s_{2}}$$[\varepsilon]$)\textbf{.}
}\par



\emph{For polynomials $s_{1}$ and $s_{2}$ and a language $L$, an instance-wise ${\mathsf{BPP}}$-black-box half-Levin reduction from $L$ to ${\sf search}\text{-}{\sf CGL}_{s_{1}}^{s_{2}}$$[\varepsilon]$ is an oracle machine $R$ as follows. On an input $z\in\{0,1\}^{*}$, $R$ deterministically makes one query, $\mathcal{E}_{z}$: a joint distribution supported on $\{0,1\}^{m(|z|)}\times\{0,1\}$ for some polynomial $m$, represented by a sampling circuit.}\par




\textbf{(Instance-wise ${\mathsf{BPP}}$-black-box)}\emph{
Consider an input }$z\in\{0,1\}^{n}$\emph{ for some }$n\in\mathbb{N}$\emph{, and suppose the oracle }$A$\emph{ is a PPT algorithm meeting the following condition:}\par



\emph{Given the query }$(z,\mathcal{E}_{z})$\emph{ made by }$R$\emph{ on input }$z$\emph{, if }$\mathcal{E}_{z}$\emph{ is not a no-instance of }${\sf CGL}_{s_{1}}^{s_{2}}[\varepsilon]$\emph{, then with probability at least }$1-2^{-\Omega(n)}$\emph{ over its internal randomness, }$A$\emph{ returns a hypothesis }$h$\emph{ of size at most }$s_{2}(m(n))$\emph{ such that}\par


\[\begin{gathered}
\Pr_{(x,b)\sim\mathcal{E}_{z}}[h(x)=b]\geq\frac{1}{2}+\frac{\varepsilon(m(n))}{2}.
\end{gathered}\]



\emph{Then, the following are guaranteed for such }$z$\emph{.}\par


\begin{itemize}


\item[] 
$\blacksquare$ 


\emph{if }$z\in L$\emph{, }$\mathcal{E}_{z}$\emph{ is a yes-instance of }${\sf CGL}_{s_{1}}^{s_{2}}[\varepsilon]$\emph{, and }$R^{A}(z)$\emph{ accepts with probability }$1-2^{-\Omega(n)}$\emph{;}\par





\item[] 
$\blacksquare$ 


\emph{if }$z\notin L$\emph{, }$\mathcal{E}_{z}$\emph{ is a no-instance of }${\sf CGL}_{s_{1}}^{s_{2}}[\varepsilon]$\emph{, and }$R^{A}(z)$\emph{ rejects with probability at least }$n^{-O(1)}$\emph{.}\par





\end{itemize}




\textbf{(Half-Levin)}\emph{ Further, there exists a polynomial-time computable function }$R_{wit}$\emph{ such that, for any input }$z^{\prime}\in\{0,1\}^{*}$\emph{, if }$z^{\prime}$\emph{ is a yes-instance of }$L$\emph{ with witness }$w$\emph{, then }$R_{wit}(z^{\prime},w)$\emph{ outputs a witness for }$\mathcal{E}_{z^{\prime}}\in{\sf CGL}_{s_{1}}^{s_{2}}[\varepsilon]$\emph{: that is, a circuit }$C$\emph{ of size at most }$s_{1}(m(n))$\emph{ such that }$\Pr_{(x,b)\sim\mathcal{E}_{z^{\prime}}}[C(x)=b]=1$\emph{.}\par



% END-SOURCE-ID Thmtheorem32



The lemma below shows that a reduction to ${\sf CGL}$ implies a reduction to agnostic PAC learning as in Definition 28.\par



% SOURCE-ID Thmtheorem33


\par\medskip\noindent\textbf{
\textbf{Lemma 33}\textbf{.}
}\par



\emph{For a language $L$, $\delta:\mathbb{N}\to[0,1]$, and polynomials $s_{1},s_{2},\varepsilon^{-1}$ with $\varepsilon(n)\leq 1/8$ for $n\in\mathbb{N}$, suppose there is a randomized many-one reduction from $L$ to ${\sf CGL}_{s_{1}}^{s_{2}}[\varepsilon]$ with failure probability $\delta$. Then, there is a randomized poly-time reduction from $L$ to agnostic PAC learning ${\mathsf{SIZE}}[s_{1}]$ by ${\mathsf{SIZE}}[s_{2}]$ with failure probability $\delta^{\prime}$, where $\delta^{\prime}(n):=\delta(n)+2^{-\Omega(n)}$ for $n\in\mathbb{N}$.}\par



% END-SOURCE-ID Thmtheorem33



\par\medskip\noindent\textbf{Proof.}\par




Let $R$ be the assumed reduction to ${\sf CGL}_{s_{1}}^{s_{2}}[\varepsilon]$. Define a reduction $R^{\prime}$ as follows:\par



On input an instance $z\in\{0,1\}^{n}$ of $L$ and a choice of randomness $r$, simulate $R(z,r)$ to produce a polynomial-size sampling circuit $\mathcal{E}$ supported over $\{0,1\}^{m}\times\{0,1\}$. Query $\mathcal{E}$ to the oracle, obtaining a hypothesis $h$ in return. If the size of $h$ is greater than $s_{2}(m)$, reject. By sampling ${\mathsf{poly}}(\varepsilon^{-1}(m)\cdot n)$ times from $\mathcal{E}$, obtain an empirical estimate of $\Pr_{(x,b)\sim\mathcal{E}}[h(x)=b]$. If the value obtained is less than $1-2\cdot\varepsilon(m)$, reject. Otherwise, accept.\par



We will show that $R^{\prime}$ constitutes a reduction to agnostic PAC learning ${\mathsf{SIZE}}[s_{1}]$ by ${\mathsf{SIZE}}[s_{2}]$.\par




Consider an $L$-instance $z$. First suppose $z\in L$. Then, with probability at least $1-\delta(n)$ over $r$, $\mathcal{E}$ is a yes-instance of ${\sf CGL}_{s_{1}}^{s_{2}}[\varepsilon]$. In this case,\par


\[\begin{gathered}
\max_{C\in{\mathsf{SIZE}}[s_{1}(m)]}\left\{\Pr_{(x,b)\sim\mathcal{E}}[C(x)=b]\right\}=1.
\end{gathered}\]


Assume that the oracle $A$ returns an $(s_{1},\varepsilon(m))$-optimal hypothesis $h$ of size $s_{2}(m)$: that is,\par


\[\begin{gathered}
\Pr_{(x,b)\sim\mathcal{E}}[h(x)=b]\geq 1-\varepsilon(m).
\end{gathered}\]


Then, by a Chernoff bound, with probability at least $1-2^{-\Omega(n)}$, $R^{\prime}$ accepts. Overall, for $z\in L$, $R^{\prime}$ accepts with probability greater than $1-(\delta(n)+2^{-\Omega(n)})$.\par




Now suppose $z$ is a no-instance of $L$. In this case, with probability at least $1-\delta(n)$ over $r$, $\mathcal{E}$ is a no-instance of ${\sf CGL}_{s_{1}}^{s_{2}}[\varepsilon]$: that is, for every circuit $C$ of size at most $s_{2}(m)$,\par


\[\begin{gathered}
\Pr_{(x,b)\sim\mathcal{E}}[C(x)=b]<\frac{1}{2}+\varepsilon(m).
\end{gathered}\]


Thus, $A$ \emph{cannot} return a hypothesis $h\in{\mathsf{SIZE}}[s_{2}(m)]$ with success better than $(1/2)+\varepsilon(m)<1-2\cdot\varepsilon(m)$ over $\mathcal{E}$. So $R^{\prime}$ rejects with probability at least $1-2^{-\Omega(n)}$. Overall, for $z\notin L$, $R^{\prime}$ accepts with probability at most $\delta(n)+2^{-\Omega(n)}$, as desired.
$\blacktriangleleft$\par




By a very similar argument, we obtain the following for the oracle-PAC setting. We omit the proof here.\par



% SOURCE-ID Thmtheorem34


\par\medskip\noindent\textbf{
\textbf{Lemma 34}\textbf{.}
}\par



\emph{For a language $L$, $\delta:\mathbb{N}\to[0,1]$, $D\in{\mathsf{P}}{\mathsf{SAMP}}$, $\lambda:\mathbb{N}\to\mathbb{N}$ with $\lambda(n)=O(\log n)$ for $n\in\mathbb{N}$, and polynomials $s_{1},s_{2},\varepsilon^{-1}$ with $\varepsilon(n)\leq 1/8$ for $n\in\mathbb{N}$, suppose there is a randomized many-one reduction from $L$ to $D\text{-}{\sf oracle}\text{-}{\sf CGL}_{s_{1}}^{s_{2}}[\varepsilon,\lambda]$ with failure probability $\delta$. Then, there is a randomized poly-time reduction from $L$ to agnostic oracle-PAC learning ${\mathsf{SIZE}}[s_{1}]$ by ${\mathsf{SIZE}}[s_{2}]$ with failure probability $\delta^{\prime}$, where $\delta^{\prime}(n):=\delta(n)+\varepsilon(n)+2^{-\Omega(n)}$ for $n\in\mathbb{N}$.}\par



% END-SOURCE-ID Thmtheorem34





\par\medskip\noindent\textbf{
2.3 Meta-complexity}\par




The following is a natural generalization of ${\mathsf{MCSP}}$ to the distributional version, where in addition to the truth table of a function $f:\{0,1\}^{n}\to\{0,1\}$, one is also given an explicit distribution $\mu$ over $\{0,1\}^{n}$ (as a table of probabilities), and is asked to distinguish “easy” $f$’s that can be well approximated by small boolean circuits from “hard” $f$’s that are almost impossible to approximate even by much larger boolean circuits.\par



% SOURCE-ID Thmtheorem35


\par\medskip\noindent\textbf{
\textbf{Definition 35}\textbf{ }(${\mathsf{GapDistMCSP}}$)\textbf{.}
}\par



\emph{The gap Distributional Minimum Circuit Size problem with size parameters $s_{1},s_{2}:\mathbb{N}\to\mathbb{N}$ and advantage parameter $\varepsilon:\mathbb{N}\to[0,1]$, denoted ${\mathsf{GapDistMCSP}}_{s_{1}}^{s_{2}}[\varepsilon]$, is the decision problem defined as follows. The input consists of the truth table of a function $f:\{0,1\}^{n}\to\{0,1\}$ (represented as a truth-table of length $2^{n}$ bits) and a probability distribution $\mu:\{0,1\}^{n}\to[0,1]$ (represented as a table of length $2^{O(n)}$ bits).}\par


\begin{itemize}


\item[] 
$\blacksquare$ 


$(f,\mu)$\emph{ is a yes-instance if there exists a circuit }$C$\emph{ of size at most }$s_{1}(n)$\emph{ such that}\par


\[\begin{gathered}
\Pr_{x\sim\mu}\left[C(x)=f(x)\right]=1.
\end{gathered}\]





\item[] 
$\blacksquare$ 


$(f,\mu)$\emph{ is a no-instance if for every circuit }$C$\emph{ of size at most }$s_{2}(n)$\emph{,}\par


\[\begin{gathered}
\Pr_{x\sim\mu}\left[C(x)=f(x)\right]<\frac{1}{2}+\varepsilon(n).
\end{gathered}\]





\end{itemize}



% END-SOURCE-ID Thmtheorem35



We note that, for an appropriate choice of the approximation parameters and distribution $\mu$, ${\mathsf{GapDistMCSP}}$ generalizes other variants of ${\mathsf{MCSP}}$, e.g., ${\mathsf{MCSP}}^{*}$ of [16] and an average-case version ${\mathsf{Ave}}{\mathsf{MCSP}}$ of [45].\par





\par\medskip\noindent\textbf{
2.4 Secret Sharing}\par




See [4] for a survey.\par



% SOURCE-ID Thmtheorem36


\par\medskip\noindent\textbf{
\textbf{Definition 36}\textbf{ }(Access Structures Represented by Monotone Formulae)\textbf{.}
}\par



\emph{An access structure is a monotone collection $\mathcal{A}\subseteq 2^{[n]}$ of non-empty subsets of $[n]$. For $T\subseteq[n]$, we say that $T$ is authorized if it belongs to $\mathcal{A}$ and unauthorized otherwise. We say that a monotone formula $\varphi:\{0,1\}^{n}\to\{0,1\}$ represents $\mathcal{A}$ if $\varphi(\chi_{T})=1$ exactly when $T\in\mathcal{A}$, where $\chi_{T}$ is the characteristic vector of $T$.}\par



% END-SOURCE-ID Thmtheorem36


% SOURCE-ID Thmtheorem37


\par\medskip\noindent\textbf{
\textbf{Definition 37}\textbf{ }(Efficient Secret Sharing Scheme)\textbf{.}
}\par



\emph{An efficient secret sharing scheme for an access structure $\mathcal{A}\subseteq 2^{[n]}$ represented by a monotone formula $\varphi$ consists of a pair of poly-time algorithms $Share$ and $Rec$ with the following properties. $Share$ takes three inputs; the first is the formula $\varphi:\{0,1\}^{n}\to\{0,1\}$, the second is a “secret” $b\in\{0,1\}$, and the third is a choice of “randomness” $r\in\{0,1\}^{|\varphi|}$. $Rec$ also takes three inputs: the formula $\varphi$, a subset of “shares” $v_{T}=(v_{i})_{i\in T}\in\left(\{0,1\}^{|\varphi|}\right)^{|T|}$, and the set $T\subseteq[n]$.}\par


\begin{itemize}


\item[] 
$\blacksquare$ 


\emph{Correctness: }\emph{Authorized parties can recover the secret. Consider any }$b\in\{0,1\}$\emph{ and }$r\in\{0,1\}^{|\varphi|}$\emph{, let }$(v_{1},...,v_{n})$\emph{ be the output of }$Share(\varphi,b;r)$\emph{, and let }$T$\emph{ be an authorized set. Then }$Rec(\varphi,v_{T},T)=b$\emph{.}\par





\item[] 
$\blacksquare$ 


\emph{Privacy: }\emph{Unauthorized parties cannot learn anything about the secret, in an information theoretic sense. Let }$T$\emph{ be an unauthorized set. Then, for }$b\sim\{0,1\}$\emph{ and }$r\sim\{0,1\}^{|\varphi|}$\emph{, the random variable }$Share(\varphi,b;r)_{T}$\emph{ is statistically independent from the random variable }$b$\emph{.}\par





\end{itemize}



% END-SOURCE-ID Thmtheorem37


% SOURCE-ID Thmtheorem38


\par\medskip\noindent\textbf{
\textbf{Lemma 38}\textbf{ }([27, 5])\textbf{.}
}\par



\emph{For every access structure $\mathcal{A}$, there exists an efficient secret sharing scheme as in Definition 37. The circuit size required to simulate Rec is at most $O(|v_{T}|^{2})$.}\par



% END-SOURCE-ID Thmtheorem38



\par\medskip\noindent\textbf{Proof sketch..}\par




We will describe the construction given in [5]. Let $\varphi:\{0,1\}^{n}\to\{0,1\}$ be a monotone formula representing $\mathcal{A}$. Starting at the root of the tree representation of $\varphi$, the \emph{Share} algorithm assigns values to the nodes in a recursive manner as follows. For a secret $b\in\{0,1\}$, assign $b$ to the root, and instantiate all shares $v_{i}$ for $i\in[n]$ to be empty. Then:\par


\begin{itemize}


\item[] 
$\blacksquare$ 


If the current gate is $\lor$ and the current value is $s$, assign $s$ to both children.\par





\item[] 
$\blacksquare$ 


if the current gate is $\land$ and the current value is $s$, randomly choose $t\sim\{0,1\}$. Assign $t$ to one child and $t\oplus s$ to the other.\par





\item[] 
$\blacksquare$ 


If the current gate is a variable $i\in[n]$ and the current value is $s$, include $s$ in the share $v_{i}$.\par





\end{itemize}


Note that the amount of randomness required is at most $|\varphi|$. The \emph{Rec} algorithm proceeds from the leaves of the tree back to the root. Start by assigning shares $(v_{i})_{i\in T}$ to leaves using the set $T$. Then:\par


\begin{itemize}


\item[] 
$\blacksquare$ 


If the current gate is $\lor$ and both children have value $s$, assign $s$ to the parent.\par





\item[] 
$\blacksquare$ 


If the current gate is $\land$ and the children have values $s$ and $t$, assign $t\oplus s$ to the parent.\par





\item[] 
$\blacksquare$ 


If the current gate is the root, return the value assigned.\par





\end{itemize}


$\blacktriangleleft$\par






\par\medskip\noindent\textbf{
2.5 Witness Encryption}\par



% SOURCE-ID Thmtheorem39


\par\medskip\noindent\textbf{
\textbf{Definition 39}\textbf{ }(Witness Encryption [11])\textbf{.}
}\par



\emph{Consider a language $L\in{\mathsf{NP}}$ with witness relation $R_{L}$, a class of function families $\Gamma$, a security function $\varepsilon:\mathbb{N}\to[0,1]$, and a length function $\ell:\mathbb{N}\to\mathbb{N}$. A $(\Gamma,\varepsilon)$-secure $\ell$-length witness encryption scheme for $L$ consists of polynomial-time algorithms ${\mathsf{Enc}}$ and ${\mathsf{Dec}}$ as follows.}\par


\begin{itemize}


\item[] 
$\blacksquare$ 


\emph{Given an }$L$\emph{-instance }$x\in\{0,1\}^{n}$\emph{, a message bit }$b\in\{0,1\}$\emph{, and randomness }$r\sim\mathcal{U}_{{\mathsf{poly}}(n)}$\emph{, }${\mathsf{Enc}}(x,b;r)$\emph{ outputs a ciphertext }$c$\emph{ of length at most }$\ell(n)$\emph{.}\par





\item[] 
$\blacksquare$ 


\emph{Given a ciphertext }$c$\emph{ and an instance/witness pair }$(x,w)\in R$\emph{, }${\mathsf{Dec}}(c,x,w)$\emph{ deterministically outputs a bit }$b\in\{0,1\}$\emph{.}\par





\end{itemize}


\emph{Moreover, ${\mathsf{Enc}}$ and ${\mathsf{Dec}}$ have the properties below.}\par


\begin{itemize}


\item[] 
$\blacksquare$ 


\emph{Correctness:}\emph{ For all sufficiently large }$n\in\mathbb{N}$\emph{, any }$b\in\{0,1\}$\emph{, }$x\in\{0,1\}^{n}\cap L$\emph{, and }$w$\emph{ such that }$(x,w)\in R_{L}$\emph{,}\par


\[\begin{gathered}
\Pr_{r}\left[{\mathsf{Dec}}({\mathsf{Enc}}(x,b;r),x,w)=b\right]=1.
\end{gathered}\]





\item[] 
$\blacksquare$ 


\emph{Security:}\emph{ For any }$A=\{A_{n}\}_{n\in\mathbb{N}}\in\Gamma$\emph{, sufficiently large }$n\in\mathbb{N}$\emph{, and }$x\in\{0,1\}^{n}\backslash L$\emph{,}\par


\[\begin{gathered}
\left|\Pr[A_{n}(x,{\mathsf{Enc}}(x,1;r))=1]-\Pr[A_{n}(x,{\mathsf{Enc}}(x,0;r))=1]\right|<\varepsilon(n).
\end{gathered}\]





\end{itemize}


\emph{If $\varepsilon$ is omitted, we assume it is some negligible function. If $\ell$ is omitted, we allow the output length of ${\mathsf{Enc}}$ to be any polynomial.}\par



% END-SOURCE-ID Thmtheorem39


% SOURCE-ID Thmtheorem40


\par\medskip\noindent\textbf{
\textbf{Definition 40}\textbf{ }(Witness Encryption in Oracle Worlds [21])\textbf{.}
}\par



\emph{Consider a language $L\in{\mathsf{NP}}$ with witness relation $R_{L}$, a class of function families $\Gamma$, a security function $\varepsilon:\mathbb{N}\to[0,1]$, and a family of distributions $D=\{D_{\lambda}\}_{\lambda\in\mathbb{N}}$, where each $D_{\lambda}$ is supported over $\{0,1\}^{2^{\lambda}}$. A $(\Gamma,\varepsilon)$-secure witness encryption scheme for $L$ with respect to $D$ is defined as in Definition 39, except the algorithms ${\mathsf{Enc}}$ and ${\mathsf{Dec}}$ both take an oracle, and they satisfy the following properties.}\par


\begin{itemize}


\item[] 
$\blacksquare$ 


\emph{Correctness:}\emph{ For all sufficiently large }$n,\lambda\in\mathbb{N}$\emph{, }$O$\emph{ in the support of }$D_{\lambda}$\emph{, }$b\in\{0,1\}$\emph{, }$x\in\{0,1\}^{n}\cap L$\emph{, and }$w$\emph{ such that }$(x,w)\in R_{L}$\emph{,}\par


\[\begin{gathered}
\Pr_{r}\left[{\mathsf{Dec}}^{O}(1^{\lambda},{\mathsf{Enc}}^{O}(1^{\lambda},x,b;r),x,w)=b\right]=1.
\end{gathered}\]





\item[] 
$\blacksquare$ 


\emph{Security:}\emph{ For all sufficiently large }$n,\lambda\in\mathbb{N}$\emph{, with probability }$1-\varepsilon(\lambda)$\emph{ over }$O\sim D_{\lambda}$\emph{, for any }$A=\{A_{n,\lambda}\}_{n,\lambda\in\mathbb{N}}\in\Gamma$\emph{ and any }$x\in\{0,1\}^{n}\backslash L$\emph{,}\par


\[\begin{gathered}
\left|\Pr_{r}[A_{n,\lambda}^{O}({\mathsf{Enc}}^{O}(1^{\lambda},x,1;r))=1]-\Pr_{r}[A_{n,\lambda}^{O}({\mathsf{Enc}}^{O}(1^{\lambda},x,0;r))=1]\right|<\varepsilon(\lambda).
\end{gathered}\]





\end{itemize}



% END-SOURCE-ID Thmtheorem40



The following result from Huang, Ilango, and Ren was originally stated with security against oracle \emph{programs} of polynomial size and query complexity; we observe that the same holds for polynomial-size circuits.\par



% SOURCE-ID Thmtheorem41


\par\medskip\noindent\textbf{
\textbf{Lemma 41}\textbf{ }([21])\textbf{.}
}\par



\emph{There exist an ${\mathsf{NP}}$-complete language $L$, a constant $\ell\in\mathbb{N}$, and a family of distributions $D=\{D_{\lambda}\}_{\lambda\in\mathbb{N}}$, where each $D_{\lambda}$ is supported over $\{0,1\}^{2^{\lambda}}$ and samplable in time ${\mathsf{poly}}(2^{\lambda})$, such that there is a $({\mathsf{SIZE}}[2^{\lambda/\ell}],2^{-\lambda/\ell})$-secure witness encryption scheme for $L$ with respect to $D$.}\par



% END-SOURCE-ID Thmtheorem41




\par\medskip\noindent\textbf{
2.6 Hardness of Approximation}\par



% SOURCE-ID Thmtheorem42


\par\medskip\noindent\textbf{
\textbf{Definition 42}\textbf{ }(${\mathsf{MMSA}}$)\textbf{.}
}\par



\emph{For $g,\theta,\Delta:\mathbb{N}\to\mathbb{N}$, ${\sf Gap}_{g}{\mathsf{MMSA}}^{\Delta}_{\theta}$ denotes the following promise problem.}\par




\emph{The input consists of a monotone formula $\varphi$ over variables $[n]$ and of length at most $\Delta(n)$ as a binary string. Let ${\mathsf{MMSA}}(\varphi)$ denote $\min\left\{\;\sum_{i\in[n]}\alpha(i)\;\middle|\;\alpha:[n]\to\{0,1\}\text{ satisfies }\varphi\;\right\}$.}\par


\begin{itemize}


\item[] 
$\blacksquare$ 


$\Pi_{Y}=\{\varphi\mid{\mathsf{MMSA}}(\varphi)\leq\theta(n)\}$\emph{}\par





\item[] 
$\blacksquare$ 


$\Pi_{N}=\{\varphi\mid{\mathsf{MMSA}}(\varphi)>g(n)\cdot\theta(n)\}$\emph{}\par





\end{itemize}



% END-SOURCE-ID Thmtheorem42


% SOURCE-ID Thmtheorem43


\par\medskip\noindent\textbf{
\textbf{Lemma 43}\textbf{ }([8, 9, 16])\textbf{.}
}\par



\emph{There exist $g,\theta,\Delta:\mathbb{N}\to\mathbb{N}$ with $g(n)=n^{(\log\log n)^{-O(1)}}$, $\theta(n)\leq n$, and $\Delta$ a polynomial with $\Delta(n)\geq n$ for $n\in\mathbb{N}$ as follows. For every language $L\in{\mathsf{NP}}$, there is a half-Levin reduction from $L$ to ${\sf Gap}_{g}{\mathsf{MMSA}}^{\Delta}_{\theta}$.}\par



% END-SOURCE-ID Thmtheorem43


% SOURCE-ID Thmtheorem44


\par\medskip\noindent\textbf{
\textbf{Definition 44}\textbf{ }(${\mathsf{CMMSA}}$ [16])\textbf{.}
}\par



\emph{For $g,\varepsilon^{-1},\theta,\Delta:\mathbb{N}\to\mathbb{N}$, ${\sf Gap}_{g,\varepsilon}{\mathsf{CMMSA}}^{\Delta}_{\theta}$ denotes the following promise problem. The input consists of a collection $\Phi=\{\varphi_{1},...,\varphi_{m}\}$ of monotone formulas over variables $[n]$, where each $\varphi_{i}$ has length at most $\Delta(n)$ as a binary string, along with a weight function $w:[n]\to\mathbb{N}$. For an assignment $\alpha:[n]\to\{0,1\}$, let $w(\alpha)$ denote $\sum_{i\in[n]}\alpha(i)\cdot w(i)$.}\par


\begin{itemize}


\item[] 
$\blacksquare$ 


$\Pi_{Y}=\{(\Phi,w)\mid\text{ there exists }\alpha:[n]\to\{0,1\}\text{ such that }$\emph{
}

$\phantom{x}\hskip 68.88867ptw(\alpha)\leq\theta(n)\text{ and }\Pr_{\varphi\sim\Phi}[\varphi(\alpha)=1]=1\}$\emph{}\par





\item[] 
$\blacksquare$ 


$\Pi_{N}=\{(\Phi,w)\mid\text{ for all }\alpha:[n]\to\{0,1\}\text{ such that }w(\alpha)\leq\theta(n)\cdot g(n)\text{, }$\emph{
}

$\phantom{x}\hskip 68.88867pt\Pr_{\varphi\sim\Phi}[\varphi(\alpha)=1]<\varepsilon\}$\emph{}\par





\end{itemize}



% END-SOURCE-ID Thmtheorem44


% SOURCE-ID Thmtheorem45


\par\medskip\noindent\textbf{
\textbf{Lemma 45}\textbf{ }([9, 16])\textbf{.}
}\par



\emph{There exists a constant $\gamma>0$ and a function $\theta:\mathbb{N}\to\mathbb{N}$ such that for all $\Delta:\mathbb{N}\to\mathbb{N}$ with $\Delta(n)=O(\log n)$, for every language $L\in{\mathsf{NP}}$, there is a half-Levin reduction from $L$ to ${\sf Gap}_{\Delta^{\gamma},\Delta^{-\gamma}}{\mathsf{CMMSA}}^{\Delta}_{\theta}$.}\par



% END-SOURCE-ID Thmtheorem45




\par\medskip\noindent\textbf{
2.7 One-way functions and Inversion}\par



% SOURCE-ID Thmtheorem46


\par\medskip\noindent\textbf{
\textbf{Definition 46}\textbf{ }($\delta$-invertible (auxiliary-input) functions)\textbf{.}
}\par



\emph{For $n,m\in\mathbb{N}$, a function $f:\{0,1\}^{n}\to\{0,1\}^{m}$, a PPT algorithm $A$, and $\delta:\mathbb{N}\to[0,1]$, we say that $A$ is a $\delta$-inverter for $f$ if}\par


\[\begin{gathered}
\Pr_{x\sim\mathcal{U}_{n}}[A(1^{n},f(x))\in f^{-1}(f(x))]\geq 1-\delta(n).
\end{gathered}\]




\emph{In the case of an auxiliary-input function $f_{z}:\{0,1\}^{\ell}\to\{0,1\}^{m}$, where $\ell,m\in\mathbb{N}$, we say that $A$ is a $\delta$-inverter for $f_{z}$ if}\par


\[\begin{gathered}
\Pr_{x\sim\mathcal{U}_{n}}[A(z,f_{z}(x))\in f_{z}^{-1}(f_{z}(x))]\geq 1-\delta(|z|).
\end{gathered}\]



% END-SOURCE-ID Thmtheorem46


% SOURCE-ID Thmtheorem47


\par\medskip\noindent\textbf{
\textbf{Definition 47}\textbf{ }(Infinitely-often one-way function)\textbf{.}
}\par



\emph{A polynomial-time computable function family $f=\{f_{n}:\{0,1\}^{n}\to\{0,1\}^{m(n)}\}_{n\in\mathbb{N}}$ is an infinitely-often one-way function if, for every PPT algorithm $A$, there exists an inverse polynomial $\delta$ such that, for infinitely many $n\in\mathbb{N}$, $A$ is not a $\delta$-inverter for $f_{n}$.\footnote{It is known that a weak one-way function is equivalent to a standard (strong) one-way function [47], so we will not distinguish between the two in this work.}}\par



% END-SOURCE-ID Thmtheorem47


% SOURCE-ID Thmtheorem48


\par\medskip\noindent\textbf{
\textbf{Definition 48}\textbf{ }(Auxiliary-input one-way function)\textbf{.}
}\par



\emph{A polynomial-time computable auxiliary-input function family $f=\{f_{z}:\{0,1\}^{\ell(|z|)}\to\{0,1\}^{m(|z|)}\}_{z\in\{0,1\}^{*}}$ is an auxiliary-input one-way function if, for every PPT algorithm $A$, there exists an inverse polynomial $\delta$ such that, for infinitely many $z\in\{0,1\}^{*}$, $A$ is not a $\delta$-inverter for $f_{z}$.}\par



% END-SOURCE-ID Thmtheorem48


% SOURCE-ID Thmtheorem49


\par\medskip\noindent\textbf{
\textbf{Definition 49}\textbf{ }(Statistical Distance)\textbf{.}
}\par



\emph{The statistical distance between two probability distributions $D^{(1)}$ and $D^{(2)}$ over domain $\{0,1\}^{n}$, denoted $\Delta(D^{(1)},D^{(2)})$, is defined as}\par


\[\begin{gathered}
\frac{1}{2}\sum_{x\in\{0,1\}^{n}}\left|D^{(1)}(x)-D^{(2)}(x)\right|,
\end{gathered}\]


\emph{or equivalently,}\par


\[\begin{gathered}
\mathop{\max}_{T\subseteq\{0,1\}^{n}}\left|\Pr_{x\sim D^{(1)}}[x\in T]-\Pr_{x\sim D^{(2)}}[x\in T]\right|.
\end{gathered}\]



% END-SOURCE-ID Thmtheorem49


% SOURCE-ID Thmtheorem50


\par\medskip\noindent\textbf{
\textbf{Definition 50}\textbf{ }(Distributionally invertible auxiliary-input functions)\textbf{.}
}\par



\emph{Consider an auxiliary input function $f_{z}$ computable uniformly in polynomial time given $z\in\{0,1\}^{n}$ and an input $x\in\{0,1\}^{\ell(n)}$. The function $f_{z}$ is said to be distributionally invertible if for every constant $b>0$ there is a PPT algorithm $I$ such that}\par


\[\begin{gathered}
\Delta((x,f_{z}(x)),(I(z,f_{z}(x)),f_{z}(x)))\leq 1/n^{b},
\end{gathered}\]


\emph{where $x\sim\mathcal{U}_{\ell(n)}$. We refer to the machine $I$ as an $n^{-b}$-distributional inverter for $f_{z}$.}\par



% END-SOURCE-ID Thmtheorem50


% SOURCE-ID Thmtheorem51


\par\medskip\noindent\textbf{
\textbf{Lemma 51}\textbf{ }(Strong inversion implies distributional inversion [25])\textbf{.}
}\par



\emph{For every polynomial-time computable auxiliary-input function $f_{z}$ and constant $b\in\mathbb{N}$, there is a constant $c\in\mathbb{N}$ such that the following holds. If there exists an $n^{-c}$-inverter for $f_{z}$, then there exists an $n^{-b}$-distributional inverter for $f_{z}$.}\par



% END-SOURCE-ID Thmtheorem51





\par\medskip\noindent\textbf{
3 Proofs of Main Results}\par




\par\medskip\noindent\textbf{
3.1 Connections between WE and learning}\par




\par\medskip\noindent\textbf{
3.1.1 Proof of Theorem 4
}\par



% SOURCE-ID Thmtheorem52


\par\medskip\noindent\textbf{
\textbf{Lemma 52}\textbf{.}
}\par



\emph{Consider a language $L\in{\mathsf{NP}}$. For $s_{2},\varepsilon^{-1},\ell:\mathbb{N}\to\mathbb{N}$, suppose a $({\mathsf{SIZE}}[s_{2}],\varepsilon)$-secure $\ell$-length witness encryption scheme $({\mathsf{Enc}},{\mathsf{Dec}})$ exists for $L$. Let the polynomial $s_{1}$ be such that $s_{1}(n)$ upper bounds the running times of ${\mathsf{Enc}}$ and ${\mathsf{Dec}}$ on $L$-instances of length $n$. Let $s^{\prime}_{1}(\ell(n))=s_{1}(n)^{2}$, $s^{\prime}_{2}(\ell(n))=s_{2}(n)$, and $2\cdot\varepsilon^{\prime}(\ell(n))=\varepsilon(n)$.}\par




\emph{Then, there is a half-Levin reduction from $L$ to ${\sf CGL}_{s_{1}^{\prime}}^{s_{2}^{\prime}}[\varepsilon^{\prime}]$ mapping $L$-instances of length $n$ to distributions supported over $\{0,1\}^{\ell(n)}\times\{0,1\}$.}\par



% END-SOURCE-ID Thmtheorem52



\par\medskip\noindent\textbf{Proof.}\par




Given an $L$-instance $z\in\{0,1\}^{n}$, the reduction to ${\sf CGL}_{s_{1}^{\prime}}^{s_{2}^{\prime}}[\varepsilon^{\prime}]$ will output a distribution $\mathcal{E}$ as follows.\par



Sample $b\sim\mathcal{U}$. Let $x:={\mathsf{Enc}}(z,b)\in\{0,1\}^{\ell(n)}$. Output $(x,b)$.\par





First consider the case that $z\in L$. Let $w$ be a witness for $z$. Let $C$ be the circuit that, with $z$ and $w$ hard-wired, given $x$ as input, outputs $b:={\mathsf{Dec}}(x,z,w)$. Note that the size of $C$ is less than $s_{1}(n)^{2}=s^{\prime}_{1}(\ell(n))$. Moreover, since $C$ is efficiently computable given $z$ and $w$, the reduction is a half-Levin reduction. By the correctness of the witness encryption scheme, we have\par


\[\begin{gathered}
\Pr_{(x,b)\sim\mathcal{E}}[C(x)=b]=1.
\end{gathered}\]


Thus, $\mathcal{E}$ is a yes-instance of ${\sf CGL}_{s_{1}^{\prime}}^{s_{2}^{\prime}}[\varepsilon^{\prime}]$.\par




Now consider the case that $z\notin L$. Toward a contradiction, suppose there exists a circuit $C$ of size at most $s_{2}^{\prime}(\ell(n))=s_{2}(n)$ such that\par


\[\begin{gathered}
\Pr_{(x,b)\sim\mathcal{E}}[C(x)=b]\geq\frac{1}{2}+\varepsilon^{\prime}(\ell(n))=\frac{1}{2}+\frac{\varepsilon(n)}{2}.
\end{gathered}\]


Then,\par


\[\begin{gathered}
\displaystyle\frac{1}{2}+\frac{\varepsilon(n)}{2} \displaystyle\leq\frac{1}{2}\cdot\Pr_{{\mathsf{Enc}},b}[C(x)=0\mid b=0]+\frac{1}{2}\cdot\Pr_{{\mathsf{Enc}},b}[C(x)=1\mid b=1] \\
\displaystyle=\frac{1}{2}+\frac{1}{2}\cdot\left(\Pr[C(x)=1\mid b=1]-\Pr[C(x)=1\mid b=0]\right),
\end{gathered}\]


which implies that\par


\[\begin{gathered}
\left|\Pr[C({\mathsf{Enc}}(z,1))=1]-\Pr[C({\mathsf{Enc}}(z,0))=1]\right|\geq\varepsilon({n}),
\end{gathered}\]


contradicting the security of the witness encryption scheme. Thus, $\mathcal{E}$ is a no-instance of ${\sf CGL}_{s_{1}^{\prime}}^{s_{2}^{\prime}}[\varepsilon^{\prime}]$.\par




This completes the proof of the lemma.
$\blacktriangleleft$\par




Observe that a ${\mathsf{P}}/{\mathsf{poly}}$-secure encryption scheme is $({\mathsf{SIZE}}[s],\varepsilon)$-secure for all polynomials $s$ and $\varepsilon^{-1}$. Thus, we obtain the corollary below, which gives the first direction of Theorem 4.\par



% SOURCE-ID Thmtheorem53


\par\medskip\noindent\textbf{
\textbf{Corollary 53}\textbf{.}
}\par



\emph{Consider any language $L\in{\mathsf{NP}}$, and suppose a ${\mathsf{P}}/{\mathsf{poly}}$-secure witness encryption scheme $({\mathsf{Enc}},{\mathsf{Dec}})$ exists for $L$. Then, there exist a polynomial $s_{1}$ and a pair of machines $(R,R_{wit})$ such that, for all polynomials $s_{2},\varepsilon^{-1}$, $(R,R_{wit})$ is an honest half-Levin reduction from $L$ to ${\sf CGL}^{s_{2}}_{s_{1}}[\varepsilon]$.}\par



% END-SOURCE-ID Thmtheorem53


% SOURCE-ID Thmtheorem54


\par\medskip\noindent\textbf{
\textbf{Lemma 54}\textbf{.}
}\par



\emph{Consider any polynomials $s_{1},m:\mathbb{N}\to\mathbb{N}$, any $s_{2},\varepsilon^{-1}:\mathbb{N}\to\mathbb{N}$, and any language $L\in{\mathsf{NP}}$. Suppose there is a half-Levin reduction $(R,R_{wit})$ from $L$ to ${\sf CGL}_{s_{1}}^{s_{2}}[\varepsilon]$, where on $L$-instances of length $n\in\mathbb{N}$, $R$ outputs distributions supported over $\{0,1\}^{m(n)}\times\{0,1\}$ in time $t_{R}(n)$, and $R_{wit}$ runs in time $t_{R_{wit}}(n)$.}\par




\emph{Then, there is a $({\mathsf{SIZE}}[s_{2}^{\prime}],\varepsilon^{\prime})$-secure $(m(n)+1)$-length witness encryption scheme for $L$, where $s_{2}^{\prime}(n)=s_{2}(m(n))/2$, $\varepsilon^{\prime}(n)=2\cdot\varepsilon(m(n))$, ${\mathsf{Enc}}$ runs in time $O(t_{R}(n))$, and ${\mathsf{Dec}}$ runs in time $O(t_{R_{wit}}(n))$.}\par



% END-SOURCE-ID Thmtheorem54




\par\medskip\noindent\textbf{Proof.}\par




We define a witness encryption scheme for $L$ as follows.\par



For $z\in\{0,1\}^{n}$ and $a\in\{0,1\}$, ${\mathsf{Enc}}(z,a)$ computes $\mathcal{E}=R(z)$, samples $(x,b)\sim\mathcal{E}$ and then outputs $(x,a\oplus b)$.\par




${\mathsf{Dec}}((x,c),z,w)$ computes $A:=R_{wit}(z,w)$ and then outputs $A(x)\oplus c$.\par





We first show the correctness of $({\mathsf{Enc}},{\mathsf{Dec}})$. Suppose $z$ is a yes-instance of $L$ with witness $w$. Then, by the correctness of the half-Levin reduction, $A=R_{wit}(z,w)$ is a witness for $\mathcal{E}=R(z)\in{\sf CGL}_{s_{1}}^{s_{2}}[\varepsilon]$; that is, $A$ is a circuit of size at most $s_{1}(m(n))$ such that\par


\[\begin{gathered}
\Pr_{(x,b)\sim\mathcal{E}}[A(x)=b]=1.
\end{gathered}\]


Then, for any $a\in\{0,1\}$ and any random choice of $(x,b)$ made by ${\mathsf{Enc}}$, it holds that $A(x)=b$, so ${\mathsf{Dec}}((x,a\oplus b),z,w)$ correctly outputs $a=b\oplus(a\oplus b)$.\par




We now argue for security. Suppose $z$ is a no-instance of $L$. Then $\mathcal{E}=R(z)$ is a no-instance of ${\sf CGL}_{s_{1}}^{s_{2}}[\varepsilon]$; that is, for any circuit $A$ of size at most $s_{2}(m(n))$,\par


\[\begin{gathered}
\Pr_{(x,b)\sim\mathcal{E}}[A(x)=b]<\frac{1}{2}+\varepsilon(m(n)).
\end{gathered}\]


Let $A^{\prime}$ be any circuit of size at most $s_{2}^{\prime}(n)=s_{2}(m(n))/2$. Toward a contradiction, suppose\footnote{We have removed the absolute value signs without loss of generality: a similar argument can be applied if the left-hand side is at most $-\varepsilon^{\prime}(n)$.}\par


\[\begin{gathered}
\Pr[A^{\prime}({\mathsf{Enc}}(z,0))=1]-\Pr[A^{\prime}({\mathsf{Enc}}(z,1))=1]\geq\varepsilon^{\prime}(n)=2\cdot\varepsilon(m(n)).
\end{gathered}\]


By definition of ${\mathsf{Enc}}$, the above can be written as\par


\[\begin{gathered}
\Pr_{(x,b)\sim\mathcal{E}}[A^{\prime}(x,b)=1]-\Pr_{(x,b)\sim\mathcal{E}}[A^{\prime}(x,1\oplus b)=1]\geq 2\cdot\varepsilon(m(n)).
\end{gathered}\]


Let $P^{A^{\prime}}$ denote Yao’s next-bit predictor applied to $A^{\prime}$. In particular,\par



On input $x\in\{0,1\}^{n}$, $P^{A^{\prime}}$ samples a random bit $c\sim\mathcal{U}$ and outputs $c$ iff $A^{\prime}(x,c)=1$.\par



By a standard argument, it holds for some fixed choice of $c\in\{0,1\}$ that\par


\[\begin{gathered}
\Pr_{(x,b)\sim\mathcal{E}}[P^{A^{\prime}}(x)=b]\geq\frac{1}{2}+\varepsilon(m(n)).
\end{gathered}\]


Finally, observe that $P^{A^{\prime}}$ with randomness fixed as above can be implemented as a circuit of size at most $s_{2}(m(n))$. This contradicts the assumption that $\mathcal{E}$ is a no-instance of ${\sf CGL}_{s_{1}}^{s_{2}}[\varepsilon]$. We conclude that ${\mathsf{Enc}}$ is $({\mathsf{SIZE}}[s_{2}^{\prime}],\varepsilon^{\prime})$-secure.\par




This completes the proof of the lemma.
$\blacktriangleleft$\par




Since a witness encryption scheme that is $({\mathsf{SIZE}}[s],\varepsilon)$-secure for all polynomials $s$ and $\varepsilon^{-1}$ is ${\mathsf{P}}/{\mathsf{poly}}$-secure, we obtain the corollary below, which gives the second direction of Theorem 4.\par



% SOURCE-ID Thmtheorem55


\par\medskip\noindent\textbf{
\textbf{Corollary 55}\textbf{.}
}\par



\emph{Consider any language $L\in{\mathsf{NP}}$. Suppose there exist a polynomial $s_{1}$ and a pair of machines $(R,R_{wit})$ such that, for all polynomials $s_{2},\varepsilon^{-1}$, $(R,R_{wit})$ is an honest half-Levin reduction from $L$ to ${\sf CGL}^{s_{2}}_{s_{1}}[\varepsilon]$. Then, there is a ${\mathsf{P}}/{\mathsf{poly}}$-secure witness encryption scheme for $L$.}\par



% END-SOURCE-ID Thmtheorem55





\par\medskip\noindent\textbf{
3.1.2 Proof of Theorem 6
}\par



% SOURCE-ID Thmtheorem56


\par\medskip\noindent\textbf{
\textbf{Lemma 56}\textbf{.}
}\par



\emph{Consider any language $L\in{\mathsf{NP}}$, and suppose a ${\mathsf{BPP}}$-secure witness encryption scheme $({\mathsf{Enc}},{\mathsf{Dec}})$ exists for $L$.}\par




\emph{Then, there exist a polynomial $s_{1}$ and a pair of machines $(R,R_{wit})$ such that, for all polynomials $s_{2}$ and $\varepsilon^{-1}$, $(R,R_{wit})$ is an honest ${\mathsf{BPP}}$-black-box half-Levin reduction from $L$ to ${\sf search}\text{-}{\sf CGL}^{s_{2}}_{s_{1}}[\varepsilon]$}\par



% END-SOURCE-ID Thmtheorem56



\par\medskip\noindent\textbf{Proof.}\par




Consider a language $L\in{\mathsf{NP}}$ with witness relation $R_{L}\subseteq\{0,1\}^{n}\times\{0,1\}^{n^{c}}$ for some constant $c\in\mathbb{N}$ and an $L$-instance $z\in\{0,1\}^{n}$. Let (${\mathsf{Enc}}$, ${\mathsf{Dec}}$) be a ${\mathsf{BPP}}$-secure $\ell$-length witness encryption scheme for $L$. Let the polynomial $s_{1}$ be such that $s_{1}(\ell(n))$ upper bounds both $n^{c}$ and the running time of ${\mathsf{Dec}}$ on $L$-instances of length $n$. As in Lemma 52, we consider the following distribution $\mathcal{E}$.\par



Sample $b\sim\mathcal{U}$. Let $x:={\mathsf{Enc}}(z,b)\in\{0,1\}^{\ell(n)}$ and output $(x,b)$.\par



Our reduction $R$ will make one query to its oracle: namely, $(z,\mathcal{E})$. After receiving a hypothesis $h$ in return, $R$ will sample from $\mathcal{E}$ and simulate $h$ to obtain an empirical estimate of $\Pr_{(x,b)\sim\mathcal{E}}[h(x)=b]$. $R$ will accept iff the estimate of this value is greater than $1/2+\varepsilon(\ell(n))/3$.\par




Let $A$ be a PPT algorithm, $z\in\{0,1\}^{n}$ an instance of $L$, and $s_{2},\varepsilon^{-1}$ arbitrary polynomials. Assume that $A$ meets the condition in Definition 32 on $z$: namely, letting $(z,\mathcal{E})$ be the query made by $R(z)$, if $\mathcal{E}$ is not a no-instance of ${\sf CGL}_{s_{1}}^{s_{2}}[\varepsilon]$, then with probability at least $1-2^{-\Omega(n)}$ over its internal randomness, $A$ returns a hypothesis $h$ with agreement at least $1/2+\varepsilon(\ell(n))/2$ over $\mathcal{E}$.\par




First consider the case that $z\in L$. Let $w\in\{0,1\}^{n^{c}}$ be a witness for $z$. Let $C$ be the circuit that, with $z$ and $w$ hard-wired, given $x$ as input, outputs $b:={\mathsf{Dec}}(x,z,w)$. Note that the size of $C$ is at most $s_{1}(\ell(n))^{2}$. Without loss of generality, assume $s_{2}(\ell(n))\geq s_{1}(\ell(n))^{2}$. By the correctness of the witness encryption scheme, we have\par


\[\begin{gathered}
\displaystyle\max_{C^{\prime}\in{\mathsf{SIZE}}[s_{2}(\ell(n))]}\left\{\Pr_{(x,b)\sim\mathcal{E}}[C^{\prime}(x)=b]\right\} \displaystyle\geq\max_{C^{\prime}\in{\mathsf{SIZE}}[s_{1}(\ell(n))^{2}]}\left\{\Pr_{(x,b)\sim\mathcal{E}}[C^{\prime}(x)=b]\right\} \\
\displaystyle\geq\Pr_{(x,b)\sim\mathcal{E}}[C(x)=b] \\
\displaystyle=1.
\end{gathered}\]


This implies that $\mathcal{E}$ is a yes-instance of ${\sf CGL}_{(s_{1})^{2}}^{s_{2}}[\varepsilon]$. Moreover, by the assumption on $A$, $A$ outputs a hypothesis with agreement at least $1/2+\varepsilon(\ell(n))/2$ over $\mathcal{E}$ with probability at least $1-2^{-\Omega(n)}$, so the reduction accepts with probability at least $1-2^{-\Omega(n)}$ overall.\par




Now consider the case that $z\notin L$. Toward a contradiction, suppose that with probability at least $1-\varepsilon(\ell(n))/4$, $A$ returns a hypothesis $h$ such that\par


\[\begin{gathered}
\Pr_{(x,b)\sim\mathcal{E}}[h(x)=b]\geq\frac{1}{2}+\frac{\varepsilon(\ell(n))}{4}.
\end{gathered}\]


Let $A^{\prime}$ be the algorithm that, on input $(z,x)$, constructs the distribution $\mathcal{E}$ from $z$ as above, simulates $A$ on $(z,\mathcal{E})$ to obtain a hypothesis $h$, and then outputs $h(x)$. Then,\par


\[\begin{gathered}
\Pr_{A\;;\;(x,b)\sim\mathcal{E}}[A^{\prime}(z,x)=b]\geq\left(1-\frac{\varepsilon(\ell(n))}{4}\right)\cdot\left(\frac{1}{2}+\frac{\varepsilon(\ell(n))}{4}\right)>\frac{1}{2}+\frac{\varepsilon(\ell(n))}{16}.
\end{gathered}\]


By the reasoning in Lemma 52 and the definition of $\mathcal{E}$, this implies that\par


\[\begin{gathered}
\left|\Pr[A^{\prime}(z,{\mathsf{Enc}}(z,1))=1]-\Pr[A^{\prime}(z,{\mathsf{Enc}}(z,0))=1]\right|\geq\frac{\varepsilon(\ell(n))}{8},
\end{gathered}\]


contradicting the security of the witness encryption scheme on input $z$. So, with probability greater than $\varepsilon(\ell(n))/4$, $A$ returns a hypothesis $h$ with agreement over $\mathcal{E}$ less than $1/2+\varepsilon(\ell(n))/4$, in which case the reduction rejects with probability $1-2^{-\Omega(n)}$. Overall, the reduction rejects with probability greater than $\varepsilon(\ell(n))/8=n^{-O(1)}$, as desired. Finally, by the assumption on $A$, $\mathcal{E}$ must be a no-instance of ${\sf CGL}_{(s_{1})^{2}}^{s_{2}}[\varepsilon]$.\par




Also note: it is easy to see that the reduction is half-Levin, since a circuit $C$ simulating ${\mathsf{Dec}}$ can always be constructed from an instance $z^{\prime}$ and witness $w^{\prime}$ in polynomial time.\par




The correctness of ${\mathsf{BPP}}$-black-box half-Levin reduction follows. This completes the proof of the lemma.
$\blacktriangleleft$\par



% SOURCE-ID Thmtheorem57


\par\medskip\noindent\textbf{
\textbf{Lemma 57}\textbf{.}
}\par



\emph{Consider any language $L\in{\mathsf{NP}}$. Suppose there exist a polynomial $s_{1}$ and a pair of machines $(R,R_{wit})$ such that, for all polynomials $s_{2},\varepsilon$, $(R,R_{wit})$ is a ${\mathsf{BPP}}$-black-box half-Levin reduction from $L$ to ${\sf search}\text{-}{\sf CGL}_{s_{1}}^{s_{2}}[\varepsilon]$. Then, ${\mathsf{BPP}}$-secure witness encryption exists for $L$.}\par



% END-SOURCE-ID Thmtheorem57




\par\medskip\noindent\textbf{Proof.}\par




Let $(R,R_{wit})$ be a ${\mathsf{BPP}}$-black-box half-Levin reduction from $L$ to ${\sf search}\text{-}{\sf CGL}$, and let the polynomial $m:\mathbb{N}\to\mathbb{N}$ be such that $R$ maps $L$-instances of length $n$ to distributions $\mathcal{E}$ supported over $\{0,1\}^{m(n)}\times\{0,1\}$. We define a witness encryption scheme for $L$ as in Lemma 54. In particular:\par



For $z\in\{0,1\}^{n}$ and $a\in\{0,1\}$, ${\mathsf{Enc}}(z,a)$ computes the query $\mathcal{E}$ made by $R(z)$, samples $(x,b)\sim\mathcal{E}$ and then outputs $(x,a\oplus b)\in\{0,1\}^{m(n)+1}$.\par




${\mathsf{Dec}}((x,b^{\prime}),z,w)$ computes $C:=R_{wit}(z,w)$ and then outputs $C(x)\oplus b^{\prime}$.\par



We will first show the correctness of $({\mathsf{Enc}},{\mathsf{Dec}})$. Suppose $z\in\{0,1\}^{n}$ is a yes-instance of $L$ with witness $w$. By the definition of a half-Levin reduction to ${\sf search}\text{-}{\sf CGL}$, $C=R_{wit}(z,w)$ is a circuit of size at most $\{0,1\}^{s_{1}(m(n))}$ such that\par


\[\begin{gathered}
\Pr_{(x,b)\sim\mathcal{E}}[C(x)=b]=1.
\end{gathered}\]


Then, for any $a\in\{0,1\}$,\par


\[\begin{gathered}
\Pr_{(x,b)\sim\mathcal{E}}[{\mathsf{Dec}}((x,a\oplus b),z,w)=a]=1,
\end{gathered}\]


as desired.\par




We will now show the security of $({\mathsf{Enc}},{\mathsf{Dec}})$. Toward a contradiction, suppose that for some PPT algorithm $A$ and polynomial $q$, for infinitely many $n\in\mathbb{N}$ and \emph{no}-instances $z\in\{0,1\}^{n}$ of $L$,\par


\[\begin{gathered}
\Pr[A(z,{\mathsf{Enc}}(z,0))=1]-\Pr[A(z,{\mathsf{Enc}}(z,1))=1]\geq\frac{1}{q(n)}.
\end{gathered}\]


Let the polynomial $s_{2}$ be such that $s_{2}(m(n))^{1/2}$ upper bounds the running time of $A$ on inputs of length $m(n)+1+n$, and let $\varepsilon(m(n)):=1/(2\cdot q(n))$.\par




Consider the PPT algorithm $A^{\prime}$ that, on input $(z,\mathcal{E})$, outputs a circuit describing $P^{A}$ as in Lemma 54, fixing the bit $c$ that maximizes success probability over $\mathcal{E}$, which $A^{\prime}$ determines by sampling from $\mathcal{E}$ and taking empirical estimates.
By the same reasoning as in Lemma 54, with probability at least $1-2^{-\Omega(n)}$, $A^{\prime}$ outputs a hypothesis $P^{A}$ of size at most $s_{2}(m(n))$ such that\par


\[\begin{gathered}
\Pr_{(x,b)\sim\mathcal{E}}[P^{A}(x)=b]\geq\frac{1}{2}+\varepsilon(m(n)).
\end{gathered}\]


Thus, for infinitely many inputs $z$, $A^{\prime}$ meets the condition on the oracle of an instance-wise ${\mathsf{BPP}}$-black-box reduction to ${\sf search}$-${\sf CGL}^{s_{2}}_{s_{1}}[\varepsilon]$: namely, given the query $(z,\mathcal{E})$ made by $R(z)$, $A^{\prime}$ returns a hypothesis with non-trivial success probability over $\mathcal{E}$. However, $\mathcal{E}$ is \emph{not} a no-instance of ${\sf CGL}_{s_{1}}^{s_{2}}[\varepsilon]$ for the no-instance $z$ of $L$. This contradicts the correctness of $R$.
$\blacktriangleleft$\par






\par\medskip\noindent\textbf{
3.1.3 Proof of Theorem 7
}\par



% SOURCE-ID Thmtheorem58


\par\medskip\noindent\textbf{
\textbf{Lemma 58}\textbf{.}
}\par



\emph{Consider a language $L\in{\mathsf{NP}}$. For $s_{2},\varepsilon^{-1}:\mathbb{N}\to\mathbb{N}$ and a constant $c\in\mathbb{N}$, assume that $({\mathsf{SIZE}}[s_{2}],\varepsilon)$-secure $c\log n$-length witness encryption exists for $L$. Let the polynomial $s_{1}$ be such that $s_{1}(n)$ upper bounds the running time of ${\mathsf{Dec}}$ on $L$-instances of length $n$. Let $s_{1}^{\prime}(c\log n)=4\cdot s_{1}(n)\cdot\log n$, $s_{2}^{\prime}(c\log n)=s_{2}(n)$, and $2\cdot\varepsilon^{\prime}(c\log n)=\varepsilon(n)$. Assume, for some ${\mathsf{E}}$-computable distribution family $D$, that ${\mathsf{E}}\nsubseteq{\mathsf{io\text{-}}}{\mathsf{SIZE}}_{D}^{\frac{1}{2}+\varepsilon^{\prime}(n)}[s_{2}^{\prime}(n)]$.}\par




\emph{Then, there is a half-Levin reduction $(R,R_{wit})$ from $L$ to ${\mathsf{GapDistMCSP}}_{s_{1}^{\prime}}^{s_{2}^{\prime}}[\varepsilon^{\prime}]$ mapping $L$-instances of length $n$ to truth-tables of length $n^{c}$, where $R_{wit}$ runs in time at most $s_{1}^{\prime}(c\log n)$.}\par



% END-SOURCE-ID Thmtheorem58



\par\medskip\noindent\textbf{Proof.}\par




The reduction $R$ will first define a distribution $\mathcal{E}$ as in the proof of Lemma 52: namely,\par



Sample $b\sim\mathcal{U}$. Let $x:={\mathsf{Enc}}(z,b)\in\{0,1\}^{c\log n}$. Output $(x,b)$.\par



Let $\mathcal{E}^{\prime}$ denote the marginal distribution of $\mathcal{E}$ over $x\in\{0,1\}^{c\log n}$.\par




The reduction $R$ then proceeds as follows. If, for all $x\in{\mathsf{supp}}(\mathcal{E}^{\prime})$, there is a unique $b_{x}\in\{0,1\}$ such that $(x,b_{x})\in{\mathsf{supp}}(\mathcal{E})$, $R$ constructs a truth-table $f:\{0,1\}^{c\log n}\to\{0,1\}$ such that\par


\[\begin{gathered}
f(x)=\begin{cases}b_{x}&x\in{\mathsf{supp}}(\mathcal{E}^{\prime})\\
0&\text{otherwise}\end{cases}
\end{gathered}\]


and outputs $(f,\mathcal{E}^{\prime})$. If, for any $x$, both $(x,0)$ and $(x,1)$ are in ${\mathsf{supp}}(\mathcal{E})$, then $R$ outputs $(f^{\prime}_{c\log n},D_{c\log n})$, where $f^{\prime}\in{\mathsf{E}}$ and $f^{\prime}\notin{\mathsf{SIZE}}_{D}^{\frac{1}{2}+\varepsilon^{\prime}(n)}[s_{2}^{\prime}(n)]$. This function exists by our assumption, and both it and a table representing $D_{c\log n}$ can be constructed uniformly in time ${\mathsf{poly}}(n)$. Note that $(f_{c\log n}^{\prime},D_{c\log n})$ is a no-instance of ${\mathsf{GapDistMCSP}}_{s_{1}^{\prime}}^{s_{2}^{\prime}}[\varepsilon^{\prime}]$.\par




In the case that $z\in L$, for all $x\in{\mathsf{supp}}(\mathcal{E}^{\prime})$, there is a unique $b_{x}\in\{0,1\}$ such that $(x,b_{x})\in{\mathsf{supp}}(\mathcal{E})$; this follows from the correctness of the witness encryption scheme. Thus, using a circuit for ${\mathsf{Dec}}$ as in the proof of Lemma 52, we get that $(f,\mathcal{E}^{\prime})$ is a yes-instance of ${\mathsf{GapDistMCSP}}_{s_{1}^{\prime}}^{s_{2}^{\prime}}[\varepsilon^{\prime}]$. Let $R_{wit}$ be the function that, given $z$ and a witness $w$ for $z\in L$, outputs a circuit describing ${\mathsf{Dec}}(-,z,w)$. It follows that $(R,R_{wit})$ is a half-Levin reduction. Note that $R_{wit}$ runs in time at most $4\cdot s_{1}(n)\cdot\log n$.\par




In the case that $z\notin L$, if both $(x,0),(x,1)\in{\mathsf{supp}}(\mathcal{E})$, then $R$ correctly outputs a no-instance of ${\mathsf{GapDistMCSP}}$. Now supposing that the labels associated with strings $x$ are unique, let $C$ be any circuit of size $s_{2}^{\prime}(c\log n)=s_{2}(n)$. Toward a contradiction, suppose\par


\[\begin{gathered}
\Pr_{(x,b_{x})\sim\mathcal{E}}[C(x)=b_{x}]=\Pr_{x\sim\mathcal{E^{\prime}}}[C(x)=f(x)]\geq\frac{1}{2}+\varepsilon^{\prime}(c\log n).
\end{gathered}\]


As in Lemma 52, this implies\par


\[\begin{gathered}
\left|\Pr_{r}[C({\mathsf{Enc}}(z,1,r))=1]-\Pr_{r}[C({\mathsf{Enc}}(z,0,r))=1]\right|\geq\varepsilon(n),
\end{gathered}\]


contradicting the security of the witness encryption scheme. Thus, $(f,\mathcal{E})$ is a no-instance of ${\mathsf{GapDistMCSP}}_{s_{1}^{\prime}}^{s_{2}^{\prime}}[\varepsilon^{\prime}]$.\par




This completes the proof of the lemma.
$\blacktriangleleft$\par



% SOURCE-ID Thmtheorem59


\par\medskip\noindent\textbf{
\textbf{Lemma 59}\textbf{.}
}\par



\emph{Consider a constant $c\in\mathbb{N}$ and $s_{1},s_{2},\varepsilon^{-1}:\mathbb{N}\to\mathbb{N}$. Suppose that there is a half-Levin reduction $(R,R_{wit})$ from ${\mathsf{SAT}}$ to ${\mathsf{GapDistMCSP}}_{s_{1}}^{s_{2}}[\varepsilon]$, where $R$ maps ${\mathsf{SAT}}$-instances of length $n\in\mathbb{N}$ to truth-tables of length $n^{c}$, and $R_{wit}$ runs in time at most $\widetilde{O}(s_{1}(c\log n))$ on $L$-instances of length $n$. Let $s_{1}^{\prime}(n)=s_{1}(c\log n)$, $s_{2}^{\prime}(n)=s_{2}(c\log n)/2$, and $\varepsilon^{\prime}(n)=2\cdot\varepsilon(c\log n)$.}\par




\emph{Then, there is a $({\mathsf{SIZE}}[s_{2}^{\prime}],\varepsilon^{\prime})$-secure $\widetilde{O}(s_{1}^{\prime}(n))$-decryption-time $(c\log n+1)$-length witness encryption scheme for ${\mathsf{SAT}}$. Moreover, there exists an ${\mathsf{E}}$-computable distribution family $D$ such that ${\mathsf{E}}\nsubseteq{\mathsf{io\text{-}}}{\mathsf{SIZE}}_{D}^{\frac{1}{2}+\varepsilon(n)}[s_{2}(n)]$.}\par



% END-SOURCE-ID Thmtheorem59




\par\medskip\noindent\textbf{Proof.}\par




We argue as in the proof of Lemma 54. In particular, the witness encryption scheme for ${\mathsf{SAT}}$ is defined as follows.\par



For $\varphi\in\{0,1\}^{n}$ and $a\in\{0,1\}$, ${\mathsf{Enc}}(\varphi,a)$ computes $(f,D)=R(\varphi)$, samples $x\sim D$ and then outputs $(x,a\oplus f(x))$.\par




${\mathsf{Dec}}((x,c),\varphi,w)$ computes $A:=R_{wit}(\varphi,w)$ and then outputs $A(x)\oplus c$.\par



Note that, the output length of ${\mathsf{Enc}}$ is $c\log n+1$ and that the running time of ${\mathsf{Dec}}$ is at most $\widetilde{O}(s_{1}^{\prime}(n))$.\par




For $\varphi\in{\mathsf{SAT}}$, $A=R_{wit}(\varphi,w)$ is a circuit of size at most $s_{1}(c\log n)$ such that\par


\[\begin{gathered}
\Pr_{x\sim D}[A(x)=f(x)]=1.
\end{gathered}\]


The correctness of $({\mathsf{Enc}},{\mathsf{Dec}})$ follows as in Lemma 54.\par




For $\varphi\notin{\mathsf{SAT}}$, by the same reasoning as in Lemma 54, we have that ${\mathsf{Enc}}$ is $({\mathsf{SIZE}}[s_{2}^{\prime}],\varepsilon^{\prime})$-secure.\par




To see the second conclusion of the lemma, let $\varphi\in\{0,1\}^{n}$ be a trivial no-instance of ${\mathsf{SAT}}$, constructible uniformly in time $O(n)$. Let $(f,D)=R(\varphi)$. By the correctness of the reduction $(R,R_{wit})$ and the definition of ${\mathsf{GapDistMCSP}}$, we have that, for all circuits $C$ of size $s_{2}(c\log n)$,\par


\[\begin{gathered}
\Pr_{x\sim D}[C(x)=f(x)]\leq\frac{1}{2}+\varepsilon(c\log n).
\end{gathered}\]


Also note that tables representing both $f$ and $D$ are uniformly constructible in time $2^{O(c\log n)}$, as desired.
$\blacktriangleleft$\par







\par\medskip\noindent\textbf{
3.2 NP-hardness of learning polynomial-size circuits}\par




We start with the following auxiliary lemmas.\par



% SOURCE-ID Thmtheorem60


\par\medskip\noindent\textbf{
\textbf{Lemma 60}\textbf{ }([40, 44])\textbf{.}
}\par



\emph{For all sufficiently large $\ell,m\in\mathbb{N}$ with $m<2^{\ell}$, for some $d=O(\ell^{2}\log m)$ there are sets $\mathcal{S}_{1},...,\mathcal{S}_{m}\subseteq[d]$ constructible in time ${\mathsf{poly}}(m,d)$ such that $|\mathcal{S}_{i}|=\ell$ and}\par


\[\begin{gathered}
\sum_{j<i}2^{|S_{i}\cap S_{j}|}\leq m-1
\end{gathered}\]


\emph{for all $i\in[m]$.}\par



% END-SOURCE-ID Thmtheorem60


% SOURCE-ID Thmtheorem61


\par\medskip\noindent\textbf{
\textbf{Lemma 61}\textbf{ }([20])\textbf{.}
}\par



\emph{For $n,t\in\mathbb{N}$, consider the “multiselector” function $Sel^{n\to t}$ that, given a string $x\in\{0,1\}^{n}$ and indices $i_{1},...,i_{t}\in[n]$, outputs $(x_{i_{1}},...,x_{i_{t}})$.}\par




\emph{For all $n,t\in\mathbb{N}$, there is a logspace-uniform boolean circuit of size $O(n+t\cdot\log^{3}(n))$ computing $Sel^{n\to t}$.}\par



% END-SOURCE-ID Thmtheorem61


% SOURCE-ID Thmtheorem62


\par\medskip\noindent\textbf{
\textbf{Lemma 62}\textbf{ }([26]; see [16] and [15])\textbf{.}
}\par



\emph{There exist constants $a,c\in\mathbb{N}$ as follows. Consider a function $f:\{0,1\}^{n}\to\{0,1\}$ and a parameter $\varepsilon=o(1)$. There is a function $\hat{f}:\{0,1\}^{a\cdot n}\to\{0,1\}$ such that the truth-table of $\hat{f}$ can be computed from that of $f$ in time ${\mathsf{poly}}(2^{n}/\varepsilon)$ given $\varepsilon$, and $\hat{f}$ has the following properties.}\par


\begin{itemize}


\item[] 
$\blacksquare$ 


\textbf{Local encodability:}\emph{ There is a non-adaptive }$f$\emph{-oracle circuit of size at most }$(n/\varepsilon)^{c}$\emph{ computing }$\hat{f}$\emph{, making at most }$1/\varepsilon^{2}$\emph{ queries to }$f$\emph{.}\par





\item[] 
$\blacksquare$ 


\textbf{Decodability:}\emph{ For any oracle }$\mathcal{O}:\{0,1\}^{*}\to\{0,1\}$\emph{ and sufficiently large time bound }$t\in\mathbb{N}$\emph{,}\par


\[\begin{gathered}
{\mathsf{K}}^{(2^{n}\cdot t/\varepsilon)^{c},\;\mathcal{O}}(f)\leq{\mathsf{K}}^{t,\;\mathcal{O}}_{\frac{1}{2}+\varepsilon}(\hat{f})+2^{n/2}\cdot(1/\varepsilon)^{c}+o(2^{n}).
\end{gathered}\]





\end{itemize}



% END-SOURCE-ID Thmtheorem62



The following is the “algorithmic information extraction” lemma of [16] modified to give a set $B$ defined in terms of ${\mathsf{K}}^{{\mathsf{poly}}}$ rather than ${\mathsf{K}}$. It is easy to see from the proof in [16] that $B$ can be defined in this way.\par



% SOURCE-ID Thmtheorem63


\par\medskip\noindent\textbf{
\textbf{Lemma 63}\textbf{.}
}\par



\emph{Consider an algorithm $D$ running in time $t_{D}$, parameters $n,\Delta,\varepsilon^{-1},\ell,d\in\mathbb{N}$, and $g_{i}:\{0,1\}^{\ell}\to\{0,1\}$ for $i\in[n]$. Let}\par


\[\begin{gathered}
\left(\left(\mathcal{S}^{(1)}_{1},...,\mathcal{S}^{(1)}_{\Delta}\right),...,\left(\mathcal{S}^{(n)}_{1},...,\mathcal{S}^{(n)}_{\Delta}\right)\right)
\end{gathered}\]


\emph{be a family of $\ell$-sized subsets of $[d]$. Define}\par


\[\begin{gathered}
G(z)=\left(z,w^{(1)}(z),...,w^{(n)}(z)\right),
\end{gathered}\]


\emph{where for $z\in\{0,1\}^{d}$,}\par


\[\begin{gathered}
w^{(i)}(z)=\left(g_{i}\left(z|_{\mathcal{S}^{(i)}_{1}}\right),...,g_{i}\left(z|_{\mathcal{S}^{(i)}_{\Delta}}\right)\right)\in\{0,1\}^{\Delta}.
\end{gathered}\]


\emph{Define a set}\par


\[\begin{gathered}
B:=\left\{\;i\in[n]\;\;\middle|\;\;{\mathsf{K}}^{4\cdot 2^{\ell}\cdot t_{D},\;D}_{\frac{1}{2}+\varepsilon}(g_{i})\leq 4\cdot(\Delta n)\cdot\sum_{j<i}2^{|S_{i}\cap S_{j}|}\;\right\}.
\end{gathered}\]


\emph{Then, for any advice string $\beta\in\{0,1\}^{O(\Delta n)}$,}\par


\[\begin{gathered}
\left|\Pr\left[D(\beta,G(z))\right]-\Pr\left[D(\beta,G(z)|_{B})\right]\right|<2\Delta n\cdot\varepsilon,
\end{gathered}\]


\emph{where $G(z)|_{B}=(z,u^{(1)},...,u^{(n)})$ is such that $u^{(i)}=w^{(i)}(z)$ for $i\in B$ and $u^{(i)}\sim\mathcal{U}_{\Delta}$ for $i\notin B$. The probabilities above are over $z\sim\mathcal{U}_{d}$ and $u^{(i)}\sim\mathcal{U}_{\Delta}$.}\par



% END-SOURCE-ID Thmtheorem63


% SOURCE-ID Thmtheorem64


\par\medskip\noindent\textbf{
\textbf{Lemma 64}\textbf{.}
}\par



\emph{For some polynomial $p$ and constant $k\in\mathbb{N}$, for any polynomials $g,\Delta,\theta,\gamma^{-1}$ with $\theta(n)\leq n$, $\Delta(n)\geq n$, and $\gamma(n)\leq 4n^{-4}$ for $n\in\mathbb{N}$, there exists a polynomial-time machine $R$ satisfying the following. For a polynomial $s_{1}$ with $s_{1}(n)\geq\gamma^{-1}(\Delta(n))^{k}$ for $n\in\mathbb{N}$, let}\par


\[\begin{gathered}
\lambda(n)=\frac{s_{1}(n)}{\theta(n)\cdot\log n},
\end{gathered}\]


\emph{and consider an advice string $f=(f_{i}\in\{0,1\}^{\lambda(n)})_{i\in[n]}$ such that, for every subset $B\subseteq[n]$,}\par


\[\begin{gathered}
{\mathsf{K}}^{p(\lambda(n))}(f_{B})\geq|B|\cdot\lambda(n)-2n\cdot\log n.
\end{gathered}\]


\emph{Then, given as input $f$ and an instance $\varphi$ of ${\sf Gap}_{g}{\mathsf{MMSA}}^{\Delta}_{\theta}$ over $n$ variables, $R$ outputs a circuit $\mathcal{E}$ sampling a flat distribution supported over $\{0,1\}^{2\Delta(n)\cdot n}\times\{0,1\}$ as follows.}\par


\begin{itemize}


\item[] 
$\blacksquare$ 


\emph{If }$\varphi$\emph{ is a yes-instance of }${\sf Gap}_{g}{\mathsf{MMSA}}^{\Delta}_{\theta}$\emph{, then there exists a circuit }$C$\emph{ of size at most }$s_{1}(n)$\emph{ such that}\par


\[\begin{gathered}
\Pr_{(x,b)\sim\mathcal{E}}[C(x)=b]=1.
\end{gathered}\]





\item[] 
$\blacksquare$ 


\emph{If }$\varphi$\emph{ is a no-instance of }${\sf Gap}_{g}{\mathsf{MMSA}}^{\Delta}_{\theta}$\emph{, then for every circuit }$C^{\prime}$\emph{ of size at most}\par


\[\begin{gathered}
s_{2}(n):=s_{1}(n)\cdot g(n)/(\log n)^{3},
\end{gathered}\]


\emph{it holds that}\par


\[\begin{gathered}
\Pr_{(x,b)\sim\mathcal{E}}[C^{\prime}(x)=b]<\frac{1}{2}+\gamma(n).
\end{gathered}\]





\end{itemize}



% END-SOURCE-ID Thmtheorem64



\par\medskip\noindent\textbf{Proof.}\par




Let $\varphi$ be a given instance of ${\sf Gap}_{g}{\mathsf{MMSA}}^{\Delta}_{\theta}$, where $\varphi$ is over $n$ variables. Let $a,c\in\mathbb{N}$ be as in Lemma 62. Consider an arbitrary polynomial $s_{1}$ such that $s_{1}(n)\geq\gamma(\Delta)^{-8c}$. Define $\ell:=a\log(\lambda)$. Let $f=(f_{i}\in\{0,1\}^{\lambda})_{i\in[n]}$ satisfy the condition stated in the lemma.\par





\par\medskip\noindent\textbf{The reduction $R$.}\par




Identifying each $f_{i}$ with a function $f_{i}:\{0,1\}^{\log\lambda}\to\{0,1\}$, let $\hat{f_{i}}:\{0,1\}^{\ell}\to\{0,1\}$ be the encoded function as in Lemma 62 with $\varepsilon=\gamma(\Delta)/(2\Delta^{2})$. Let $\left(\mathcal{S}^{(1)}_{1},...,\mathcal{S}^{(1)}_{\Delta}\right),...,$ $\left(\mathcal{S}^{(n)}_{1},...,\mathcal{S}^{(n)}_{\Delta}\right)$ be the sets of size $\ell$ guaranteed by Lemma 60 with $d=O(\ell^{2}\cdot\log(\Delta n))=o(\Delta n)$.\par




For any choice of $z\in\{0,1\}^{d}$, $r\in\{0,1\}^{\Delta}$, and $b\in\{0,1\}$, define $(v_{1},...,v_{n})=Share(\varphi,b;r)\in(\{0,1\}^{\Delta})^{n}$, and define the string $x(z,r,b)$ as\par


\[\begin{gathered}
\left(\;z\;,\;\left(\hat{f}_{1}\left(z|_{\mathcal{S}^{(1)}_{1}}\right),...,\hat{f}_{1}\left(z|_{\mathcal{S}^{(1)}_{\Delta}}\right)\right)\oplus v_{1}\;,\;...\;,\;\left(\hat{f}_{n}\left(z|_{\mathcal{S}^{(n)}_{1}}\right),...,\hat{f}_{n}\left(z|_{\mathcal{S}^{(n)}_{\Delta}}\right)\right)\oplus v_{n}\;\right),
\end{gathered}\]


padded to have length $2\Delta n$. The output of the reduction is a circuit sampling the following distribution $\mathcal{E}$ :\par



For uniformly random $z\sim\mathcal{U}_{d}$, $r\sim\mathcal{U}_{\Delta}$, and $b\sim\{0,1\}$, output $(x(z,r,b),b)$.\par



Note that $\mathcal{E}$ has a ${\mathsf{poly}}(n)$-size sampling circuit and that the marginal distribution of $\mathcal{E}$ over $x$ is flat.\par





\par\medskip\noindent\textbf{Completeness.}\par




To see the completeness of the reduction, suppose ${\mathsf{MMSA}}(\varphi)\leq\theta$, and let $T_{\alpha}=\{j\in[n]\mid\alpha(j)=1\}$ where $\alpha:[n]\to\{0,1\}$ is an assignment of Hamming weight at most $\theta$ satisfying $\varphi$.\par




We will construct a circuit $C$ of size at most $s_{1}(n)$ such that $C(x)=b$ for all $(x,b)$ in the support of $\mathcal{E}$. $C$ takes as input a string $x=(z,w_{1},...,w_{n})$ with $z\in\{0,1\}^{d}$ and $w_{i}\in\{0,1\}^{\Delta}$ for $i\in[n]$, padded to have length $2n\Delta$. $C$ uses $\theta\Delta$ encoding circuits as in Lemma 62 to compute the values\par


\[\begin{gathered}
\left\{\;\hat{f}_{j}\left(z|_{\mathcal{S}^{(j)}_{m}}\right)\;\middle|\;j\in T_{\alpha}\;,\;m\in[\Delta]\;\right\}
\end{gathered}\]


by making at most $\theta\Delta\cdot(1/\varepsilon)^{2}<\gamma^{-4}(\Delta)$ total queries to $f_{T_{\alpha}}$. $C$ answers these queries with the string $f_{T_{\alpha}}$ hard-wired along with a multiselector circuit as in Lemma 61. $C$ then XORs the outputs of the encoding circuits with the appropriate inputs $w_{j}$ to obtain an authorized set of shares $(v_{j})_{j\in T_{\alpha}}$. Finally, $C$ applies $Rec$ with $\varphi$ and $T_{\alpha}$ hard-wired to obtain the secret $b$.\par




$C$ consists of $2n\Delta$ input gates, $\theta\lambda$ gates hard-wiring $f_{T_{\alpha}}$, $O(\theta\lambda+\gamma^{-4}(\Delta)\cdot\log^{3}(\theta\lambda))=O(\theta\lambda)$ gates for multiselection, $\theta\Delta$ encoding circuits each of size $(\log\lambda/\varepsilon)^{c}=o(\lambda)$, $O(\theta\Delta)$ gates to XOR, and $O((\theta\Delta)^{2})$ gates for Rec. This amounts to at most $O(\theta\lambda)<s_{1}(n)$ gates in total.\par




Thus, there exists a circuit $C$ of size at most $s_{1}(n)$ such that\par


\[\begin{gathered}
\Pr_{(x,b)\sim\mathcal{E}}[C(x)=b]=1.
\end{gathered}\]





\par\medskip\noindent\textbf{Soundness.}\par




We now argue for the soundness of the reduction. Suppose ${\mathsf{MMSA}}(\varphi)>\theta\cdot g(n)$.
For a fixed circuit $C^{\prime}$ of size $s_{2}(n)=s_{1}(n)\cdot g(n)/(\log n)^{3}$ , define a program $D$ as follows. $D$ takes as advice a string $\beta(b,r):=(b,v_{1},...,v_{n})\in\{0,1\}^{O(n\cdot\Delta)}$, where $(v_{1},...,v_{n})=Share(\varphi,b;r)$, and accepts its input $(z,w_{1},...,w_{n})$ if and only if\par


\[\begin{gathered}
C^{\prime}(z,w_{1}\oplus v_{1},...,w_{n}\oplus v_{n})=b.
\end{gathered}\]


Note that $D$ runs in time at most $t_{D}:=4\cdot s_{2}(n)\cdot\log(s_{2}(n))$. Now, define a set $B$ as in Lemma 63, applying the lemma with $g_{i}=\hat{f}_{i}$ for $i\in[n]$. In particular,\par


\[\begin{gathered}
B:=\left\{\;i\in[n]\;\;\middle|\;\;{\mathsf{K}}^{4\cdot\lambda^{a}\cdot t_{D},\;D}_{\frac{1}{2}+\varepsilon}(\hat{f}_{i})\leq(\Delta n)^{2}\;\right\}.
\end{gathered}\]


By Lemma 63,\par


\[\begin{gathered}
\displaystyle\left|\Pr[D(\beta,G(z))]-\Pr[D(\beta,G(z)|_{B})]\right| \displaystyle<2\Delta^{2}\cdot\varepsilon \\
\displaystyle=\gamma(\Delta),
\end{gathered}\]


where $G$ and $G(z)|_{B}$ are as defined in that lemma.\par




We claim that $B$ is not too large. Note that for $i\in B$, by the decodability statement of Lemma 62, for some polynomial $q$,\par


\[\begin{gathered}
\displaystyle{\mathsf{K}}^{q(\lambda),\;D}(f_{i}) \displaystyle\leq{\mathsf{K}}^{4\cdot\lambda^{a}\cdot t_{D},\;D}_{\frac{1}{2}+\varepsilon}(\hat{f}_{i})+\lambda^{1/2}\cdot(1/\varepsilon)^{c}+o(\lambda) \\
\displaystyle<o(\lambda).
\end{gathered}\]


Thus, for $p(\lambda):=n^{2}\cdot q(\lambda)\cdot t_{D}$,\par


\[\begin{gathered}
\displaystyle|f_{B}|=|B|\cdot\lambda \displaystyle\leq{\mathsf{K}}^{p(\lambda)}(f_{B})+2n\cdot\log n \\
\displaystyle\leq{\mathsf{K}}^{n^{2}\cdot q(\lambda),\;D}(f_{B})+|D|+2n\cdot\log
n \\
\displaystyle\leq\sum_{i\in B}\left({\mathsf{K}}^{q(\lambda),\;D}(f_{i})\right)+O(|C^{\prime}|\cdot\log|C^{\prime}|)+|\beta|+O(n\cdot\log n) \\
\displaystyle<\frac{s_{1}(n)\cdot g(n)}{2\log n}+o(|B|\cdot\lambda) \\
\displaystyle=\frac{\lambda\cdot\theta\cdot g(n)}{2}+o(|B|\cdot\lambda).
\end{gathered}\]


It follows that\par


\[\begin{gathered}
\displaystyle|B| \displaystyle\leq\frac{\theta\cdot g(n)}{2\cdot(1-o(1))} \\
\displaystyle<\theta\cdot g(n),
\end{gathered}\]


so $B$ is not an authorized set.\par




Now observe that\par


\[\begin{gathered}
\displaystyle\Pr_{z,r,b}[C^{\prime}(x(z,r,b))=b] \displaystyle=\Pr[D(\beta(b,r),G(z))=1] \\
\displaystyle<\Pr[D(\beta(b,r),G(z)|_{B})=1]+\gamma(\Delta) \\
\displaystyle\leq\frac{1}{2}+\gamma(\Delta),
\end{gathered}\]


where the last line follows from the security of $Share$. We conclude that for any circuit $C^{\prime}$ of size at most $s_{2}(n)$,\par


\[\begin{gathered}
\Pr_{(x,b)\sim\mathcal{E}}[C^{\prime}(x)=b]<\frac{1}{2}+\gamma(\Delta)\leq\frac{1}{2}+\gamma(n).
\end{gathered}\]


This completes the proof of the lemma.
$\blacktriangleleft$\par



% SOURCE-ID Thmtheorem65


\par\medskip\noindent\textbf{
\textbf{Lemma 65}\textbf{ }(Claim 7.3 of [16])\textbf{.}
}\par



\emph{For a uniformly random string $f=(f_{i})_{i\in[n]}\sim(\mathcal{U}_{\lambda})^{n}$, it holds with probability $1-2^{-n}$ that for all subsets $B\subseteq[n]$,}\par


\[\begin{gathered}
{\mathsf{K}}(f_{B})\geq|B|\cdot\lambda-2n\cdot\log n.
\end{gathered}\]



% END-SOURCE-ID Thmtheorem65


% SOURCE-ID Thmtheorem66


\par\medskip\noindent\textbf{
\textbf{Theorem 66}\textbf{ }(Restatement of Theorem 8)\textbf{.}
}\par



\emph{For some constant $c\in\mathbb{N}$ and $g:\mathbb{N}\to\mathbb{N}$ with $g(n)=n^{1/(\log\log n)^{O(1)}}$ for $n\in\mathbb{N}$, for any polynomials $\gamma^{-1},s_{1}$ such that $s_{1}(n)\geq\gamma^{-c}(n)$ for $n\in\mathbb{N}$, there is a randomized many-one reduction from ${\mathsf{SAT}}$ to ${\sf CGL}_{s_{1}}^{s_{1}\cdot g}[\gamma]$.}\par



% END-SOURCE-ID Thmtheorem66




\par\medskip\noindent\textbf{Proof.}\par




By Lemma 43, there is a reduction $R_{0}$ from ${\mathsf{SAT}}$ to ${\sf Gap}_{g_{0}}{\mathsf{MMSA}}_{\theta}^{\Delta}$ for some\par


\[\begin{gathered}
g_{0}(n)=n^{\frac{1}{(\log\log n)^{O(1)}}},
\end{gathered}\]


polynomial $\Delta(n)\geq n$, and $\theta(n)\leq n$. Let $k\in\mathbb{N}$ be the constant of Lemma 64, and define the constant $c\in\mathbb{N}$ such that $n^{c}>\Delta(n)^{k}$. For arbitrary polynomials $\gamma^{-1},s_{1}$, define $\gamma^{\prime}(n):=\gamma(2n\cdot\Delta(n))$ and $s_{1}^{\prime}(n):=s_{1}(2n\cdot\Delta(n))$ for $n\in\mathbb{N}$. We will apply Lemma 64 with parameters $s_{1}^{\prime},\gamma^{\prime}$, $g_{0}$, $\Delta$, and $\theta$.\par




More specifically, our reduction will first apply $R_{0}$ to obtain an instance $\varphi$ of ${\sf Gap}_{g_{0}}{\mathsf{MMSA}}_{\theta}^{\Delta}$ on $n$ variables, and then it will sample a random string $f=(f_{i})_{i\in[n]}\sim(\mathcal{U}_{\lambda})^{n}$. It will then apply $R$ as in Lemma 64 to obtain an instance $\mathcal{E}$ of ${\sf CGL}_{s_{1}}^{s_{2}}[\gamma]$ supported over $\{0,1\}^{2n\cdot\Delta(n)}\times\{0,1\}$, where\par


\[\begin{gathered}
\displaystyle s_{2}(n\cdot\Delta(n)) \displaystyle=\frac{s_{1}(2n\cdot\Delta(n))\cdot g_{0}(n)}{(\log n)^{3}} \\
\displaystyle=s_{1}(2n\cdot\Delta(n))\cdot g(2n\cdot\Delta(n)),
\end{gathered}\]


where we define $g(2n\cdot\Delta(n)):=g_{0}(n)/(\log n)^{3}$.\par




By Lemma 65, with probability $1-2^{-n}$ over the choice of $f$, $f$ satisfies the condition of Lemma 64. The desired statement is then immediate from the correctness of $R$ and the definition of ${\sf CGL}$.
$\blacktriangleleft$\par




Corollary 9 now follows by combining Theorem 8 with Lemma 33.\par






\par\medskip\noindent\textbf{
3.3 NP-hardness of improper learning in other settings}\par



% SOURCE-ID Thmtheorem67


\par\medskip\noindent\textbf{
\textbf{Theorem 67}\textbf{ }(Restatement of Theorem 12)\textbf{.}
}\par



\emph{There exist a distribution $D\in{\mathsf{P}}{\mathsf{SAMP}}$, an ${\mathsf{NP}}$-complete language $L$, and a polynomial $s_{1}$ such that, for all polynomials $s_{2},\varepsilon^{-1}$, there is a function $\lambda:\mathbb{N}\to\mathbb{N}$ such that $\lambda(n)=O(\log n)$ for $n\in\mathbb{N}$ and $L$ reduces to $D\text{-}{\sf oracle}\text{-}{\sf CGL}^{s_{2}}_{s_{1}}[\varepsilon,\lambda]$ under an honest half-Levin reduction.}\par



% END-SOURCE-ID Thmtheorem67



\par\medskip\noindent\textbf{Proof.}\par




Let $D=\{D_{\lambda}\}_{\lambda\in\mathbb{N}}\in{\mathsf{P}}{\mathsf{SAMP}}$ be the distribution family and $L$ the ${\mathsf{NP}}$-complete language of Lemma 41 with witness relation $V_{L}$. That is, each $D_{\lambda}$ is supported over truth-tables of functions $O:\{0,1\}^{\lambda}\to\{0,1\}$, and there is an oracle witness encryption scheme $({\mathsf{Enc}},{\mathsf{Dec}})$ for $L$ with respect to $D$ secure against programs of size and query complexity $2^{\lambda/\ell}$ and with advantage $2^{-\lambda/\ell}$ for some constant $\ell\in\mathbb{N}$.\par




Let $d\in\mathbb{N}$ be such that $n^{d}$ upper-bounds the running time of ${\mathsf{Dec}}^{(-)}$ on $L$-instances of length $n$, and define $s_{1}(n)=n^{2d}$. Let $s_{2}$ and $\varepsilon$ be given, and define $\lambda(n)=\ell\cdot\log(s_{2}(n)\cdot\varepsilon(n)^{-1})=O(\log n)$ for $n\in\mathbb{N}$.\par




We now give a reduction from $L$ to $D\text{-}{\sf oracle}\text{-}{\sf CGL}^{s_{2}}_{s_{1}}[\varepsilon,\lambda]$. Given an $L$-instance $z\in\{0,1\}^{n}$ for some $n\in\mathbb{N}$, the reduction outputs $\mathcal{E}^{(-)}$ as follows.\par



Given access to an oracle $O:\{0,1\}^{\lambda(n)}\to\{0,1\}$, sample $b\sim\mathcal{U}$, and let $x={\mathsf{Enc}}^{O}(1^{\lambda(n)},z,b)$. Output $(x,b)$.\par





First suppose $z\in L$. Let $w$ be such that $(z,w)\in V_{L}$. Consider the oracle circuit $C^{(-)}$ that, with $(z,w)$ hard-wired, given input $x$ and oracle $O$, simulates ${\mathsf{Dec}}^{O}(1^{\lambda(n)},x,z,w)$ and returns its output $b$. Note that $C^{(-)}$ has size at most $s_{1}(n)$. By the correctness guarantee of $({\mathsf{Enc}},{\mathsf{Dec}})$, for every $O$ in the support of $D_{\lambda(n)}$ and every $(x,b)$ in the support of $\mathcal{E}$, it holds that $C^{O}(x)=b$. Therefore, in this case, the reduction outputs a yes-instance of $D\text{-}{\sf oracle}\text{-}{\sf CGL}^{s_{2}}_{s_{1}}[\varepsilon,\lambda]$.\par




Now suppose $z\notin L$, and let $C^{(-)}$ be any oracle circuit of size at most $s_{2}(n)$. Note that\par


\[\begin{gathered}
2^{\lambda/\ell}=s_{2}(n)\cdot\varepsilon(n)^{-1}>s_{2}(n),
\end{gathered}\]


so $({\mathsf{Enc}},{\mathsf{Dec}})$ is secure against $C$. Suppose for some oracle $O$ that\par


\[\begin{gathered}
\Pr_{(x,b)\sim\mathcal{E}}[C^{O}(x)=b]\geq\frac{1}{2}+\varepsilon(n).
\end{gathered}\]


By the same reasoning as in Lemma 52,\par


\[\begin{gathered}
\left|\Pr[C^{O}({\mathsf{Enc}}^{O}(1^{\lambda(n)},z,1))=1]-\Pr[C^{O}({\mathsf{Enc}}^{O}(1^{\lambda(n)},z,0))=1]\right|\geq\varepsilon(n)>\frac{1}{2^{\lambda(n)/\ell}}.
\end{gathered}\]


By the security of the witness encryption scheme, this occurs with probability at most $2^{-\lambda(n)/\ell}<\varepsilon(n)$ over $O\sim D_{\lambda(n)}$. Thus, in this case, the reduction outputs a no-instance of $D\text{-}{\sf oracle}\text{-}{\sf CGL}^{s_{2}}_{s_{1}}[\varepsilon,\lambda]$.
$\blacktriangleleft$\par




By combining Theorem 12 with Lemma 34, we obtain the following.\par



% SOURCE-ID Thmtheorem68


\par\medskip\noindent\textbf{
\textbf{Corollary 68}\textbf{ }(Restatement of Corollary 13)\textbf{.}
}\par



\emph{For some polynomial $s$, for any constant $c\in\mathbb{N}$, it is ${\mathsf{NP}}$-hard under a randomized one-query reduction to agnostically oracle-PAC learn ${\mathsf{SIZE}}[s(n)]$ by ${\mathsf{SIZE}}[s(n)^{c}]$.}\par



% END-SOURCE-ID Thmtheorem68



Theorem 12 also implies ${\mathsf{NP}}$-hardness of learning RAMs in the improper setting, as indicated below.\par



% SOURCE-ID Thmtheorem69


\par\medskip\noindent\textbf{
\textbf{Corollary 69}\textbf{ }(Restatement of Corollary 14)\textbf{.}
}\par



\emph{For all polynomials $g$ and $\varepsilon^{-1}$, the problem ${\sf Gap}_{g,\varepsilon}{\sf RAM}\text{-}{\sf MINLT}$ is ${\mathsf{NP}}$-hard under a randomized many-one reduction.}\par



% END-SOURCE-ID Thmtheorem69




\par\medskip\noindent\textbf{Proof.}\par




We show that $D\text{-}{\sf oracle}\text{-}{\sf CGL}$ as in Theorem 12 reduces to ${\sf Gap}{\sf RAM}\text{-}{\sf MINLT}$. The main idea is to include the whole truth-table of an oracle $O$ sampled from $D_{\lambda(n)}$ in all the inputs to a RAM.\par




Let the polynomials $g$ and $\varepsilon^{-1}$ be given. Let $s_{1}$ be as in Theorem 12. Let $s_{2}(n)=g(s_{1}(n))^{4}$, and let $\lambda$ be as in Theorem 12 applied with $s_{2}$ and $\varepsilon$. Let $D$ be the distribution guaranteed by Theorem 12, and recall that $D_{\lambda(n)}$ is supported over truth-tables of length $2^{\lambda(n)}$, where $s_{2}(n)\cdot\varepsilon^{-1}(n)<2^{\lambda(n)}\leq{\mathsf{poly}}(n)$.\par




The reduction, given an instance $\mathcal{E}^{(-)}$ of $D\text{-}{\sf oracle}\text{-}{\sf CGL}_{s_{1}}^{s_{2}}[\varepsilon,\lambda]$, samples $O\sim D_{\lambda(n)}$ and then defines the following distribution $\mathcal{E}^{\prime}$:\par



Sample $(x,b)\sim\mathcal{E}^{O}$, let $x^{\prime}=(x,O)$, and output $(x^{\prime},b)$.\par



The reduction outputs $(\mathcal{E}^{\prime},1^{s_{1}(n)^{2}},1^{s_{1}(n)^{2}})$.\par




First suppose that $\mathcal{E}^{(-)}$ is a yes-instance of $D\text{-}{\sf oracle}\text{-}{\sf CGL}$. In this case, there is an oracle circuit $C^{(-)}$ of size at most $s_{1}(n)$ such that, for all $O$ in the support of $D_{\lambda(n)}$ and all $(x,b)$ in the support of $\mathcal{E}^{O}$, $C^{O}(x)=b$. Note that $C$ can be simulated by a program $M$ of size and running time at most $s_{1}(n)^{2}$. We obtain that for all $(x^{\prime},b)\in{\mathsf{supp}}(\mathcal{E}^{\prime})$,\par


\[\begin{gathered}
M^{x^{\prime}}=M^{x,O}=b,
\end{gathered}\]


so $(\mathcal{E}^{\prime},1^{s_{1}(n)^{2}},1^{s_{1}(n)^{2}})$ is a yes-instance of ${\sf Gap}_{g,\varepsilon}{\sf RAM}\text{-}{\sf MINLT}$.\par




Now suppose that $\mathcal{E}^{(-)}$ is a no-instance of $D\text{-}{\sf oracle}\text{-}{\sf CGL}$. Consider any random-access machine $M$ of size and running time at most $g(s_{1}(n)^{2})$. Note that $M$ can be simulated by an oracle circuit $C$ of size at most $s_{2}(n)$, where we may let the “$x$” part of the input $x^{\prime}=(x,O)\in\{0,1\}^{n}\times\{0,1\}^{2^{\lambda(n)}}$ be a standard (non-oracle) input to $C$. Then, since $\mathcal{E}^{(-)}$ is a no-instance of $D\text{-}{\sf oracle}\text{-}{\sf CGL}$, with probability at least $1-\varepsilon(n)$ over $O\sim D_{\lambda(n)}$,\par


\[\begin{gathered}
\Pr_{(x,b)\sim\mathcal{E}^{O}}\left[C^{O}(x)=b\right]<\frac{1}{2}+\varepsilon(n).
\end{gathered}\]


This implies, for such oracles $O$, that\par


\[\begin{gathered}
\Pr_{(x^{\prime},b)\sim\mathcal{E}^{\prime}}\left[M^{x^{\prime}}=b\right]<\frac{1}{2}+\varepsilon(n).
\end{gathered}\]


Thus, with probability at least $1-\varepsilon(n)$ over the randomness of the reduction,


$(\mathcal{E}^{\prime},1^{s_{1}(n)^{2}},1^{s_{1}(n)^{2}})$ is a no-instance of ${\sf Gap}_{g,\varepsilon}{\sf RAM}\text{-}{\sf MINLT}$, as desired.
$\blacktriangleleft$\par






\par\medskip\noindent\textbf{
3.4 Excluding Pessiland, Heuristica, and Minicrypt}\par




We show that if the hardness of learning can be based on the hardness of ${\mathsf{NP}}$ (under a non-adaptive reduction), then so can the existence of OWFs. Applying our main results connecting witness encryption and the ${\mathsf{NP}}$-hardness of learning, we obtain Theorems 16 and 17 as consequences.\par




\par\medskip\noindent\textbf{
3.4.1 Proof of Theorem 15
}\par




We will make use of the following lemmas.\par



% SOURCE-ID Thmtheorem70


\par\medskip\noindent\textbf{
\textbf{Lemma 70}\textbf{.}
}\par



\emph{Let $f=\{f_{z}\}_{z\in\{0,1\}^{*}}$ be a polynomial-time computable auxiliary input function. If infinitely-often one-way functions do not exist, then for every distribution family $D\in{\mathsf{P}}{\mathsf{SAMP}}$ and constant $b\in\mathbb{N}$, there exists a PPT machine $I$ such that for all sufficiently large $n\in\mathbb{N}$,}\par


\[\begin{gathered}
\Pr_{z\sim D_{n}}[\;I\text{ is an }n^{-b}\text{-inverter for }f_{z}\;]\geq 1-\frac{1}{n^{b}}.
\end{gathered}\]


\emph{Moreover, if $I$ does not $n^{-b}$-invert $f_{z}$ for some $z\in\{0,1\}^{n}$, then for any $y\in\{0,1\}^{*}$,}\par


\[\begin{gathered}
\Pr_{I}[\;I(z,y)\text{ outputs }\bot\;]\geq 1-2^{-n}.
\end{gathered}\]



% END-SOURCE-ID Thmtheorem70



\par\medskip\noindent\textbf{Proof.}\par




Assume that i.o. one-way functions do not exist. Let $S$ be a sampler for $D$. For a constant $b\in\mathbb{N}$, let $A$ be an $n^{-4b}$-inverter for the function $g$ defined as $g(r,x):=f_{S(r)}(x)$. Define the PPT machine $I$ as follows.\par



On input $(z,y)$, sample uniformly random strings $y_{1},...,y_{n^{8b}}$. Compute the fraction $\gamma$ of $i\in[n^{8b}]$ such that $A(z,y_{i})$ fails to output a pre-image of its input under $g$. If $\gamma$ is less than $2/n^{4b}$, output $A(z,y)$. Otherwise, output $\bot$.\par



We first claim that, with probability at least $1-n^{-2b}$ over $z\sim D$, $A(z,-)$ is an $n^{-b}$ inverter for $f_{z}$. Toward a contradiction, suppose the claim does not hold. That is, with probability greater than $n^{-2b}$ over $z\sim D_{n}$, $A(z,-)$ only inverts $f_{z}$ with success probability $1-n^{-b}$. But then the overall probability of $A$ inverting $g$, for uniformly random $r$ and $y$, is at most\par


\[\begin{gathered}
1-n^{-2b}+n^{-2b}\cdot(1-n^{-b})<1-n^{-4b},
\end{gathered}\]


contradicting our assumption on $A$. Moreover, by a Chernoff bound, the probability that $\gamma$ is greater than $2/n^{4b}$ is less than $2^{-n}$. By the definition of $I$ and a union bound, this proves the first part of the lemma.\par




For the second part, suppose that $A(z,-)$ does not $n^{-b}$-invert $f_{z}$. Then the expected value of $\gamma$ in the definition of $I$ is at least $n^{-b}$. By a Chernoff bound, the probability that $\gamma$ is less than $2/n^{4b}$ is less than $2^{-n}$. This completes the proof of the lemma.
$\blacktriangleleft$\par




We will require the following lemma, which states that if ${\mathsf{NP}}$ is easy on average for randomized algorithms, then agnostic PAC learning over efficiently samplable distributions is possible.\par



% SOURCE-ID Thmtheorem71


\par\medskip\noindent\textbf{
\textbf{Lemma 71}\textbf{ }([12])\textbf{.}
}\par



\emph{If ${\mathsf{Dist}}{\mathsf{NP}}\subseteq{\mathsf{Avg}}{\mathsf{BPP}}$, then there is a PPT machine that, for any polynomial $s$, for all sufficiently large $n\in\mathbb{N}$, agnostically learns ${\mathsf{SIZE}}[s(n)]$ in the improper setting over distributions in ${\mathsf{P}}{\mathsf{SAMP}}/{\mathsf{poly}}$.}\par



% END-SOURCE-ID Thmtheorem71



It remains to exhibit a reduction from ${\mathsf{SAT}}$ to inverting an auxiliary input one-way function, assuming ${\mathsf{NP}}$-hardness of learning. The following lemma is as in [2]. We include a proof for completeness.\par



% SOURCE-ID Thmtheorem72


\par\medskip\noindent\textbf{
\textbf{Lemma 72}\textbf{ }([2])\textbf{.}
}\par



\emph{Suppose there is a randomized non-adaptive honest reduction $R$ and a polynomial $s:\mathbb{N}\to\mathbb{N}$ such that, for every constant $c\in\mathbb{N}$, $R$ reduces ${\mathsf{SAT}}$ to agnostically PAC learning ${\mathsf{SIZE}}[s(n)]$ by ${\mathsf{SIZE}}[s(n)^{c}]$. Then there exist a poly-time computable auxiliary input function family $f=\{f_{\varphi}\}_{\varphi\in\{0,1\}^{*}}$, a constant $b\in\mathbb{N}$, and a PPT oracle machine $A$ satisfying the following: for any $n\in\mathbb{N}$, $\varphi\in\{0,1\}^{n}$, and PPT machine $I$ that $n^{-b}$-distributionally inverts $f_{\varphi}$,}\par


\[\begin{gathered}
A^{I}(\varphi)={\mathsf{SAT}}(\varphi).
\end{gathered}\]



% END-SOURCE-ID Thmtheorem72




\par\medskip\noindent\textbf{Proof.}\par




Let $R$ be the assumed reduction from ${\mathsf{SAT}}$ to agnostic PAC learning, let the polynomial $t_{R}$ denote its running time, and let the inverse polynomial $\varepsilon$ be such that $R$ requires an $(s,2\cdot\varepsilon)$-optimal hypothesis in response to each query. Define $\delta(n):=\varepsilon(n)/t_{R}(n)^{2}$. Define an auxiliary-input function\par


\[\begin{gathered}
f_{\varphi}(r_{0},i,r):=(r_{0},i,\mathcal{E}_{i}^{(\varphi,r_{0})}(r)[1]),
\end{gathered}\]


where $\mathcal{E}_{i}^{(\varphi,r_{0})}$ (henceforth denoted $\mathcal{E}_{i}$) is the $i^{\text{th}}$ query of $R$ on input $\varphi$ and randomness $r_{0}$, sampling a distribution $(X_{i},Y_{i})$. In particular, for $r\sim\mathcal{U}_{t_{R}}$, $\mathcal{E}_{i}(r)[1]$ represents a random sample $x\sim X_{i}$. Note that all such samplers $\mathcal{E}_{i}$ can be simulated uniformly in polynomial time given $(\varphi,r_{0},i)$, by the efficiency of $R$.\par




Let $I$ be a PPT machine that distributionally inverts $f_{\varphi}$. In particular, we require the success probability of $I$ to be such that, with probability $1-o(1)$ over $r_{0}\sim\mathcal{U}_{t_{R}(n)}$, for every $i\in[t_{R}(n)]$, $I$ $\delta(n)^{10}$-distributionally inverts $f_{\varphi}(r_{0},i,-)$.\par




For $i\in[t_{R}(n)]$, define hypotheses $h_{i}$ as follows:\par



On input $x$, invoke $I$ to find a pre-image $r$ under $f_{\varphi}(r_{0},i,-)$ of $(r_{0},i,x)$, and then return $y:=\mathcal{E}_{i}(r)[2]$,\par



and define a machine $A$ (taking oracle access to $I$) as follows:\par



On input $\varphi\in\{0,1\}^{n}$ and randomness $r_{0}\in\{0,1\}^{t_{R}(n)}$, simulate $R(\varphi;r_{0})$, obtaining queries $\mathcal{E}_{1},...,\mathcal{E}_{m}$. Answer the $i^{\text{th}}$ query of $R$ with the hypothesis $H_{i}$, which evaluates $h_{i}$ on $\delta(n)^{-8}$ independent choices of randomness (used by $I$) and then outputs the majority output $y\in\{0,1\}$. After answering all the queries this way and completing the simulation of $R$, accept $\varphi$ iff $R$ accepts.\par



For $i\in[t_{R}(n)]$, let $f_{opt}^{(i)}$ be a fixed function (of arbitrary complexity) that maximizes\par


\[\begin{gathered}
\Pr_{(x,y)\sim(X_{i},Y_{i})}[f_{opt}^{(i)}(x)=y].
\end{gathered}\]


We will argue that $H_{i}(x)$ predicts $y$ almost as well as $f_{opt}^{(i)}$ over $(X_{i},Y_{i})$.\par




For any $x\in{\mathsf{supp}}(X_{i})$, let the random variable $Y_{i}^{x}$ denote the conditional distribution $y\sim\mathcal{E}_{i}(r)[2]$ given that $\mathcal{E}_{i}(r)[1]=x$, for $r\sim\mathcal{U}_{t_{R}(n)}$. For simplicity, start by considering the hypotheses $\hat{h}_{i}$ and $\hat{H}_{i}$, which are the same as $h_{i}$ and $H_{i}$ except taking an oracle $\hat{I}$ that “perfectly” distributionally inverts $f_{\varphi}(r_{0},i,-)$. That is,\par


\[\begin{gathered}
\displaystyle\Delta\left(\;(r,f_{\varphi}(r_{0},i,r))\;,\;(\hat{I}(\varphi,f_{\varphi}(r_{0},i,r)),f_{\varphi}(r_{0},i,r))\;\right)=0,
\end{gathered}\]


for $r\sim\mathcal{U}_{t_{R}(n)}$. Observe that for $x\sim X_{i}$, outputs of $\hat{h}_{i}(x)$ are distributed exactly according to $Y_{i}^{x}$. In particular, for $r\sim\mathcal{U}_{t_{R}(n)}$ and $x=\mathcal{E}_{i}(r)[1]$,\par


\[\begin{gathered}
\displaystyle\Delta\left(\;(x,Y^{x}_{i})\;,\;(x,\hat{h}_{i}(x))\;\right) \displaystyle=\Delta\left(\;(x,\mathcal{E}_{i}(r)[2])\;,\;(x,\mathcal{E}_{i}(\hat{I}(\varphi,(r_{0},i,x)))[2])\;\right) \\
\displaystyle\leq\Delta\left(\;(x,r)\;,\;(x,\hat{I}(\varphi,(r_{0},i,x))))\;\right) \\
\displaystyle=\Delta\left(\;(r,f_{\varphi}(r_{0},i,r))\;,\;(\hat{I}(\varphi,f_{\varphi}(r_{0},i,r)),f_{\varphi}(r_{0},i,r))\;\right) \\
\displaystyle=0.
\end{gathered}\]


Now, for $i\in[t_{R}(n)]$ and strings $x$ in the support of $X_{i}$, define\par


\[\begin{gathered}
\alpha_{i}(x):=\max_{y\in\{0,1\}}\left\{\Pr[Y_{i}^{x}=y]\right\}.
\end{gathered}\]


That is, $y\in\{0,1\}$ is the label most likely to be associated with $x$ by $(X_{i},Y_{i})$, and $\alpha_{i}(x)$ is the corresponding conditional probability. Note that $f_{opt}^{(i)}(x)$ must always output this optimal label $y$, so $\alpha_{i}(x)$ is the success probability of $f_{opt}^{(i)}$ on input $x$. Moreover, by Eq. (2), for every $x\in{\mathsf{supp}}(X_{i})$,\par


\[\begin{gathered}
\displaystyle\Pr_{\hat{h}_{i}}\left[\hat{h}_{i}(x)=f_{opt}^{(i)}(x)\right]=\alpha_{i}(x).
\end{gathered}\]




We now fix $i$ and $x$ and break the subsequent analysis into two cases: either $f_{opt}^{(i)}$ predicts the label of $x$ with probability substantially bounded away from $1/2$, or not.\par





\par\medskip\noindent\textbf{Case I: $\alpha_{i}(x)\geq 1/2+\delta(n)^{2}$.}\par




For $j\in[\delta(n)^{-8}]$, let $z_{j}$ be the Bernoulli random variable that equals $1$ if $\hat{h}_{i}(x)=f_{opt}^{(i)}(x)$ in the $j^{\text{th}}$ iteration of $\hat{H}_{i}(x)$, and $0$ otherwise. Recall that $\hat{H}_{i}(x)=f_{opt}^{(i)}(x)$ iff the former event holds in a majority of trials $j\in[\delta(n)^{-8}]$. Let\par


\[\begin{gathered}
Z:=\sum_{j\in[\delta(n)^{-8}]}z_{j},
\end{gathered}\]


and note that, by Eq. (3),\par


\[\begin{gathered}
\mathop{\mathbb{E}}[Z]=\delta(n)^{-8}\cdot\alpha_{i}(x)\geq\delta(n)^{-8}\cdot\left(\frac{1}{2}+\delta(n)\right).
\end{gathered}\]


Then, by a Chernoff bound,\par


\[\begin{gathered}
\displaystyle\Pr\left[\hat{H}_{i}(x)\neq f_{opt}^{(i)}(x)\right] \displaystyle=\Pr\left[Z\leq\delta(n)^{-8}/2\right] \\
\displaystyle\leq\Pr\left[Z\leq\mathbb{E}[Z]\cdot(1-\delta(n)^{3})\right] \\
\displaystyle\leq 2^{-\Omega(n)}.
\end{gathered}\]





\par\medskip\noindent\textbf{Case II: $\alpha_{i}(x)<1/2+\delta(n)^{2}$.}\par




Now,\par


\[\begin{gathered}
\displaystyle\Pr_{y\sim Y_{i}^{x}}[\hat{H}_{i}(x)\neq y]-\Pr_{y\sim Y_{i}^{x}}[f_{opt}^{(i)}(x)\neq y] \displaystyle\leq\frac{1}{2}-(1-\alpha_{i}(x)) \\
\displaystyle\leq\delta(n)^{2}.
\end{gathered}\]




Thus, in both Cases I and II, it holds that\par


\[\begin{gathered}
\Pr_{y\sim Y_{i}^{x}}[\hat{H}_{i}(x)\neq y]-\Pr_{y\sim Y_{i}^{x}}[f_{opt}^{(i)}(x)\neq y]\leq\delta(n)^{2}.
\end{gathered}\]


Moreover,\par


\[\begin{gathered}
\Delta(\hat{H}_{i}(x),H_{i}(x))\leq\delta(n)^{-8}\cdot\delta(n)^{10}=\delta(n)^{2}.
\end{gathered}\]


for $x\sim X_{i}$. Overall, by a union bound and the definition of $\delta$,\par


\[\begin{gathered}
\Pr_{A}[A(\varphi)={\mathsf{SAT}}(\varphi)]\geq\frac{2}{3}-2\cdot t_{R}(n)\cdot\delta(n)^{2}-o(1)=\frac{2}{3}-o(1),
\end{gathered}\]


as desired.
$\blacktriangleleft$\par




We are now ready to finish the proof of Theorem 15, restated below.\par



% SOURCE-ID Thmtheorem73


\par\medskip\noindent\textbf{
\textbf{Theorem 73}\textbf{ }(Restatement of Theorem 15)\textbf{.}
}\par



\emph{Suppose there is a randomized non-adaptive honest reduction $R$ and a polynomial $s:\mathbb{N}\to\mathbb{N}$ such that, for every constant $c\in\mathbb{N}$, $R$ reduces ${\mathsf{SAT}}$ to agnostically PAC learning ${\mathsf{SIZE}}[s(n)]$ by ${\mathsf{SIZE}}[s(n)^{c}]$. Then, unless ${\mathsf{NP}}\subseteq{\mathsf{BPP}}$, there exist infinitely-often one-way functions.}\par



% END-SOURCE-ID Thmtheorem73




\par\medskip\noindent\textbf{Proof.}\par




Assume the non-existence of i.o. one-way functions. Let $R$ be the assumed randomized non-adaptive reduction from ${\mathsf{SAT}}$ to agnostic learning ${\mathsf{SIZE}}[s]$ running in polynomial time $t_{R}$. Let the inverse polynomial $\varepsilon$ be such that $R$ requires $(s,\varepsilon)$-optimal hypotheses in response to each query, and let the constant $b\in\mathbb{N}$ be as in Lemma 72. Combining Lemma 70 with Lemma 51, for every distribution family $D\in{\mathsf{P}}{\mathsf{SAMP}}$, there exists a PPT machine $I$ such that, for all sufficiently large $n\in\mathbb{N}$,\par


\[\begin{gathered}
\Pr_{z\sim D_{n}}[\;I\text{ is an }n^{-b}\text{-distributional inverter for }f_{z}\;]\geq 1-\frac{1}{n^{b}},
\end{gathered}\]


and if $I$ does not $n^{-b}$-invert $f_{z}$ for some $z\in\{0,1\}^{*}$, then for any $y\in\{0,1\}^{*}$,\par


\[\begin{gathered}
\Pr_{I}[\;I(z,y)\text{ outputs }\bot\;]\geq 1-2^{-n}.
\end{gathered}\]


Along with Lemma 72, this implies ${\mathsf{Dist}}{\mathsf{NP}}\subseteq{\mathsf{Avg}}{\mathsf{BPP}}$. Let $B$ be the PPT agnostic learner for ${\mathsf{SIZE}}[s]$ guaranteed by Lemma 71, with accuracy and confidence parameters $\delta$ chosen such that $\delta^{-1}(n)\geq\max\{t_{R}(n)^{2},\varepsilon^{-1}(n)\}$ for all $n\in\mathbb{N}$.\par




Now consider the PPT machine $R^{B}$ defined as follows.\par



On input $\varphi\in\{0,1\}^{n}$ and randomness $r_{0}\sim\mathcal{U}_{t_{R}(n)}$, simulate $R(\varphi;r_{0})$, obtaining queries $\mathcal{E}_{1},...,\mathcal{E}_{m}$, where each $\mathcal{E}_{i}$ samples a distribution $(X_{i},Y_{i})$. For each such $i\in[m]$, let $h_{i}$ be the output of $B$ on input $\mathcal{E}_{i}$. Answer the $i^{\text{th}}$ query of $R$ with $h_{i}$. Accept iff $R$ accepts.\par



By Lemma 71, for all sufficiently large $n\in\mathbb{N}$, for each $i\in[m]$, with probability $1-\delta(n)$ over the randomness of $B$, $B$ returns an $(s,\varepsilon)$-optimal hypothesis $h_{i}$ for $\mathcal{E}_{i}$.\par




By a union bound and the definition of $R$,\par


\[\begin{gathered}
\Pr_{R^{B}}[R^{B}(\varphi)={\mathsf{SAT}}(\varphi)]\geq\frac{2}{3}-\frac{1}{t_{R}(n)}.
\end{gathered}\]


We conclude that ${\mathsf{SAT}}\in{\mathsf{BPP}}$, as desired.
$\blacktriangleleft$\par






\par\medskip\noindent\textbf{
3.4.2 Proofs of Theorems 16 and 17
}\par




We will need the following lemma. For the reader’s convenience, we give a sketch of the proof in Appendix B. Though not originally stated for the infinitely-often setting, it is easily seen to follow from the same argument.\par



% SOURCE-ID Thmtheorem74


\par\medskip\noindent\textbf{
\textbf{Lemma 74}\textbf{ }([11])\textbf{.}
}\par



\emph{Suppose ${\mathsf{BPP}}$-secure witness encryption exists for ${\mathsf{NP}}$. Then, if infinitely-often one-way functions exist, there exists public-key encryption secure infinitely often against polynomial-time adversaries.}\par



% END-SOURCE-ID Thmtheorem74



Theorem 16 is now immediate from Theorem 4, Lemma 33, Theorem 15, and Lemma 74.\par




Theorem 17 follows from Theorem 4 and Theorem 16.\par







\par\medskip\noindent\textbf{
4 Concluding Remarks}\par




We have shown that the existence of secure witness encryption is equivalent to ${\mathsf{NP}}$-hardness of a certain learning problem (${\sf CGL}$) under restricted deterministic reductions. This may be taken as evidence that computational learning is indeed ${\mathsf{NP}}$-hard, if one believes in cryptography. At the same time, this result also suggests that one may want to consider a more general kind of reductions when trying to prove ${\mathsf{NP}}$-hardness of learning, as establishing the existence of secure witness encryption seems hard. On the other hand, actually proving ${\mathsf{NP}}$-hardness of learning for restricted deterministic reductions would imply a breakthrough result that public-key cryptography can be based on the worst-case assumption that ${\mathsf{NP}}\not\subseteq{\mathsf{BPP}}$.\par




There are many interesting research directions to explore in the context of cryptography and ${\mathsf{NP}}$-hardness of meta-complexity problems. For example,
Mazor and Pass [38] have recently showed that the existence of secure Indistinguishability Obfuscation (a cryptographic primitive that implies Witness Encryption) would mean that a polynomial-gap version of ${\mathsf{MCSP}}$ cannot be shown ${\mathsf{NP}}$-complete under Levin reductions, unless ${\mathsf{NP}}\subseteq{\mathsf{BPP}}$. Can one get similar results for some computational learning problems? What if one has secure Witness Encryption only? Another possible research direction is to find out whether existing tools from ${\mathsf{NP}}$-hardness of learning could unconditionally yield new non-trivial forms of witness encryption (eg. with security against restricted circuit classes).\par




Lastly, of course, the ${\mathsf{NP}}$-hardness of improper learning for ${\mathsf{P}}/{\mathsf{poly}}$ remains open. We have made some partial progress, and we have demonstrated cryptographic consequences of ${\mathsf{NP}}$-hardness under certain strong formulations, but we are optimistic that further progress can be made.\par





\par\medskip\noindent\textbf{References}\par



\begin{itemize}


\item[] 
[1]

Michael Alekhnovich, Samuel R. Buss, Shlomo Moran, and Toniann Pitassi.


Minimum propositional proof length is np-hard to linearly
approximate.


\emph{J. Symb. Log.}, 66(1):171–191, 2001.


doi:10.2307/2694916.





\item[] 
[2]

Benny Applebaum, Boaz Barak, and David Xiao.


On basing lower-bounds for learning on worst-case assumptions.


In \emph{49th Annual IEEE Symposium on Foundations of Computer
Science, FOCS 2008, October 25-28, 2008, Philadelphia, PA, USA}, pages
211–220. IEEE Computer Society, 2008.


doi:10.1109/FOCS.2008.35.





\item[] 
[3]

Boaz Barak, Oded Goldreich, Russell Impagliazzo, Steven Rudich, Amit Sahai,
Salil P. Vadhan, and Ke Yang.


On the (im)possibility of obfuscating programs.


\emph{J. ACM}, 59(2):6:1–6:48, 2012.


doi:10.1145/2160158.2160159.





\item[] 
[4]

Amos Beimel.


Secret-sharing schemes: A survey.


In Yeow Meng Chee, Zhenbo Guo, San Ling, Fengjing Shao, Yuansheng
Tang, Huaxiong Wang, and Chaoping Xing, editors, \emph{Coding and Cryptology
– Third International Workshop, IWCC 2011, Qingdao, China, May 30-June 3,
2011. Proceedings}, volume 6639 of \emph{Lecture Notes in Computer Science},
pages 11–46. Springer, 2011.


doi:10.1007/978-3-642-20901-7\_2.





\item[] 
[5]

Josh Cohen Benaloh and Jerry Leichter.


Generalized secret sharing and monotone functions.


In Shafi Goldwasser, editor, \emph{Advances in Cryptology – CRYPTO
’88, 8th Annual International Cryptology Conference, Santa Barbara,
California, USA, August 21-25, 1988, Proceedings}, volume 403 of \emph{Lecture
Notes in Computer Science}, pages 27–35. Springer, 1988.


doi:10.1007/0-387-34799-2\_3.





\item[] 
[6]

Avrim Blum, Merrick Furst, Michael Kearns, and Richard J. Lipton.


Cryptographic primitives based on hard learning problems.


In Douglas R. Stinson, editor, \emph{Advances in Cryptology —
CRYPTO’ 93}, pages 278–291, Berlin, Heidelberg, 1994. Springer Berlin
Heidelberg.





\item[] 
[7]

Marco L. Carmosino, Russell Impagliazzo, Valentine Kabanets, and Antonina
Kolokolova.


Agnostic learning from tolerant natural proofs.


In Klaus Jansen, José D. P. Rolim, David Williamson, and
Santosh S. Vempala, editors, \emph{Approximation, Randomization, and
Combinatorial Optimization. Algorithms and Techniques, APPROX/RANDOM 2017,
August 16-18, 2017, Berkeley, CA, USA}, volume 81 of \emph{LIPIcs}, pages
35:1–35:19. Schloss Dagstuhl – Leibniz-Zentrum für Informatik, 2017.


doi:10.4230/LIPICS.APPROX-RANDOM.2017.35.





\item[] 
[8]

Irit Dinur, Prahladh Harsha, and Guy Kindler.


Polynomially low error pcps with polyloglog n queries via modular
composition.


In Rocco A. Servedio and Ronitt Rubinfeld, editors, \emph{Proceedings
of the Forty-Seventh Annual ACM on Symposium on Theory of Computing, STOC
2015, Portland, OR, USA, June 14-17, 2015}, pages 267–276. ACM, 2015.


doi:10.1145/2746539.2746630.





\item[] 
[9]

Irit Dinur and Shmuel Safra.


On the hardness of approximating label-cover.


\emph{Inf. Process. Lett.}, 89(5):247–254, 2004.


doi:10.1016/J.IPL.2003.11.007.





\item[] 
[10]

Sanjam Garg, Craig Gentry, Shai Halevi, Mariana Raykova, Amit Sahai, and Brent
Waters.


Candidate indistinguishability obfuscation and functional encryption
for all circuits.


\emph{SIAM J. Comput.}, 45(3):882–929, 2016.


doi:10.1137/14095772X.





\item[] 
[11]

Sanjam Garg, Craig Gentry, Amit Sahai, and Brent Waters.


Witness encryption and its applications.


In Dan Boneh, Tim Roughgarden, and Joan Feigenbaum, editors, \emph{Symposium on Theory of Computing Conference, STOC’13, Palo Alto, CA, USA,
June 1-4, 2013}, pages 467–476. ACM, 2013.


doi:10.1145/2488608.2488667.





\item[] 
[12]

Halley Goldberg, Valentine Kabanets, Zhenjian Lu, and Igor C. Oliveira.


Probabilistic kolmogorov complexity with applications to average-case
complexity.


In Shachar Lovett, editor, \emph{37th Computational Complexity
Conference, CCC 2022, July 20-23, 2022, Philadelphia, PA, USA}, volume
234 of \emph{LIPIcs}, pages 16:1–16:60. Schloss Dagstuhl – Leibniz-Zentrum
für Informatik, 2022.


doi:10.4230/LIPICS.CCC.2022.16.





\item[] 
[13]

Dan Gutfreund and Amnon Ta-Shma.


Worst-case to average-case reductions revisited.


In Moses Charikar, Klaus Jansen, Omer Reingold, and José D. P.
Rolim, editors, \emph{Approximation, Randomization, and Combinatorial
Optimization. Algorithms and Techniques, 10th International Workshop,
APPROX 2007, and 11th International Workshop, RANDOM 2007, Princeton, NJ,
USA, August 20-22, 2007, Proceedings}, volume 4627 of \emph{Lecture Notes in
Computer Science}, pages 569–583. Springer, 2007.


doi:10.1007/978-3-540-74208-1\_41.





\item[] 
[14]

Johan Håstad, Russell Impagliazzo, Leonid A. Levin, and Michael Luby.


A pseudorandom generator from any one-way function.


\emph{SIAM J. Comput.}, 28(4):1364–1396, 1999.


doi:10.1137/S0097539793244708.





\item[] 
[15]

Shuichi Hirahara.


Non-disjoint promise problems from meta-computational view of
pseudorandom generator constructions.


In Shubhangi Saraf, editor, \emph{35th Computational Complexity
Conference, CCC 2020, July 28-31, 2020, Saarbrücken, Germany (Virtual
Conference)}, volume 169 of \emph{LIPIcs}, pages 20:1–20:47. Schloss Dagstuhl
– Leibniz-Zentrum für Informatik, 2020.


doi:10.4230/LIPICS.CCC.2020.20.





\item[] 
[16]

Shuichi Hirahara.


Np-hardness of learning programs and partial MCSP.


In \emph{63rd IEEE Annual Symposium on Foundations of Computer
Science, FOCS 2022, Denver, CO, USA, October 31 – November 3, 2022}, pages
968–979. IEEE, 2022.


doi:10.1109/FOCS54457.2022.00095.





\item[] 
[17]

Shuichi Hirahara.


Capturing one-way functions via np-hardness of meta-complexity.


In Barna Saha and Rocco A. Servedio, editors, \emph{Proceedings of the
55th Annual ACM Symposium on Theory of Computing, STOC 2023, Orlando, FL,
USA, June 20-23, 2023}, pages 1027–1038. ACM, 2023.


doi:10.1145/3564246.3585130.





\item[] 
[18]

Shuichi Hirahara and Mikito Nanashima.


Learning in pessiland via inductive inference.


In \emph{2023 IEEE 64th Annual Symposium on Foundations of Computer
Science (FOCS)}, pages 447–457, 2023.


doi:10.1109/FOCS57990.2023.00033.





\item[] 
[19]

Shuichi Hirahara and Mikito Nanashima.


One-way functions and zero knowledge.


In Bojan Mohar, Igor Shinkar, and Ryan O’Donnell, editors, \emph{Proceedings of the 56th Annual ACM Symposium on Theory of Computing, STOC
2024, Vancouver, BC, Canada, June 24-28, 2024}, pages 1731–1738. ACM,
2024.


doi:10.1145/3618260.3649701.





\item[] 
[20]

Justin Holmgren and Ron Rothblum.


Linear-size boolean circuits for multiselection.


In Rahul Santhanam, editor, \emph{39th Computational Complexity
Conference, CCC 2024, July 22-25, 2024, Ann Arbor, MI, USA}, volume 300
of \emph{LIPIcs}, pages 11:1–11:20. Schloss Dagstuhl – Leibniz-Zentrum
für Informatik, 2024.


doi:10.4230/LIPICS.CCC.2024.11.





\item[] 
[21]

Yizhi Huang, Rahul Ilango, and Hanlin Ren.


Np-hardness of approximating meta-complexity: A cryptographic
approach.


In Barna Saha and Rocco A. Servedio, editors, \emph{Proceedings of the
55th Annual ACM Symposium on Theory of Computing, STOC 2023, Orlando, FL,
USA, June 20-23, 2023}, pages 1067–1075. ACM, 2023.


doi:10.1145/3564246.3585154.





\item[] 
[22]

Russell Impagliazzo.


A personal view of average-case complexity.


In \emph{Proceedings of the Tenth Annual Structure in Complexity
Theory Conference, Minneapolis, Minnesota, USA, June 19-22, 1995}, pages
134–147. IEEE Computer Society, 1995.


doi:10.1109/SCT.1995.514853.





\item[] 
[23]

Russell Impagliazzo, Valentine Kabanets, and Ilya Volkovich.


Synergy between circuit obfuscation and circuit minimization.


In Nicole Megow and Adam D. Smith, editors, \emph{Approximation,
Randomization, and Combinatorial Optimization. Algorithms and Techniques,
APPROX/RANDOM 2023, September 11-13, 2023, Atlanta, Georgia, USA}, volume
275 of \emph{LIPIcs}, pages 31:1–31:21. Schloss Dagstuhl – Leibniz-Zentrum
für Informatik, 2023.


doi:10.4230/LIPICS.APPROX/RANDOM.2023.31.





\item[] 
[24]

Russell Impagliazzo and Leonid A. Levin.


No better ways to generate hard NP instances than picking uniformly
at random.


In \emph{31st Annual Symposium on Foundations of Computer Science, St.
Louis, Missouri, USA, October 22-24, 1990, Volume II}, pages 812–821.
IEEE Computer Society, 1990.


doi:10.1109/FSCS.1990.89604.





\item[] 
[25]

Russell Impagliazzo and Michael Luby.


One-way functions are essential for complexity based cryptography
(extended abstract).


In \emph{30th Annual Symposium on Foundations of Computer Science,
Research Triangle Park, North Carolina, USA, 30 October – 1 November 1989},
pages 230–235. IEEE Computer Society, 1989.


doi:10.1109/SFCS.1989.63483.





\item[] 
[26]

Russell Impagliazzo and Avi Wigderson.


\emph{P = BPP} if \emph{E} requires exponential circuits:
Derandomizing the XOR lemma.


In Frank Thomson Leighton and Peter W. Shor, editors, \emph{Proceedings of the Twenty-Ninth Annual ACM Symposium on the Theory of
Computing, El Paso, Texas, USA, May 4-6, 1997}, pages 220–229. ACM, 1997.


doi:10.1145/258533.258590.





\item[] 
[27]

Mitsuru Ito, Akira Saio, and Takao Nishizeki.


Multiple assignment scheme for sharing secret.


\emph{J. Cryptol.}, 6(1):15–20, 1993.


doi:10.1007/BF02620229.





\item[] 
[28]

Aayush Jain, Huijia Lin, and Amit Sahai.


Indistinguishability obfuscation from well-founded assumptions.


\emph{Commun. ACM}, 67(3):97–105, 2024.


doi:10.1145/3611095.





\item[] 
[29]

Valentine Kabanets and Jin-yi Cai.


Circuit minimization problem.


In F. Frances Yao and Eugene M. Luks, editors, \emph{Proceedings of
the Thirty-Second Annual ACM Symposium on Theory of Computing, May 21-23,
2000, Portland, OR, USA}, pages 73–79. ACM, 2000.


doi:10.1145/335305.335314.





\item[] 
[30]

Michael J. Kearns, Robert E. Schapire, and Linda Sellie.


Toward efficient agnostic learning.


\emph{Mach. Learn.}, 17(2-3):115–141, 1994.


doi:10.1007/BF00993468.





\item[] 
[31]

Ker-I Ko.


On the complexity of learning minimum time-bounded turing machines.


\emph{SIAM J. Comput.}, 20(5):962–986, 1991.


doi:10.1137/0220059.





\item[] 
[32]

Caleb Koch, Carmen Strassle, and Li-Yang Tan.


Properly learning decision trees with queries is np-hard.


In \emph{64th IEEE Annual Symposium on Foundations of Computer
Science, FOCS 2023, Santa Cruz, CA, USA, November 6-9, 2023}, pages
2383–2407. IEEE, 2023.


doi:10.1109/FOCS57990.2023.00146.





\item[] 
[33]

Ilan Komargodski, Tal Moran, Moni Naor, Rafael Pass, Alon Rosen, and Eylon
Yogev.


One-way functions and (im)perfect obfuscation.


In \emph{55th IEEE Annual Symposium on Foundations of Computer
Science, FOCS 2014, Philadelphia, PA, USA, October 18-21, 2014}, pages
374–383. IEEE Computer Society, 2014.


doi:10.1109/FOCS.2014.47.





\item[] 
[34]

Ilan Komargodski, Moni Naor, and Eylon Yogev.


Secret-sharing for NP.


In Palash Sarkar and Tetsu Iwata, editors, \emph{Advances in
Cryptology – ASIACRYPT 2014 – 20th International Conference on the Theory
and Application of Cryptology and Information Security, Kaoshiung, Taiwan,
R.O.C., December 7-11, 2014, Proceedings, Part II}, volume 8874 of \emph{Lecture Notes in Computer Science}, pages 254–273. Springer, 2014.


doi:10.1007/978-3-662-45608-8\_14.





\item[] 
[35]

Yanyi Liu, Noam Mazor, and Rafael Pass.


A note on zero-knowledge for NP and one-way functions.


\emph{Electron. Colloquium Comput. Complex.}, TR24-095, 2024.


URL: https://eccc.weizmann.ac.il/report/2024/095.





\item[] 
[36]

Yanyi Liu, Noam Mazor, and Rafael Pass.


On witness encryption and laconic zero-knowledge arguments.


\emph{Electron. Colloquium Comput. Complex.}, TR24-194, 2024.


URL: https://eccc.weizmann.ac.il/report/2024/194.





\item[] 
[37]

Zhenjian Lu, Noam Mazor, Igor C. Oliveira, and Rafael Pass.


Lower bounds on the overhead of indistinguishability obfuscation.


\emph{Electron. Colloquium Comput. Complex.}, TR24-146, 2024.


URL: https://eccc.weizmann.ac.il/report/2024/146.





\item[] 
[38]

Noam Mazor and Rafael Pass.


Gap MCSP is not (levin) np-complete in obfustopia.


In Rahul Santhanam, editor, \emph{39th Computational Complexity
Conference, CCC 2024, July 22-25, 2024, Ann Arbor, MI, USA}, volume 300
of \emph{LIPIcs}, pages 36:1–36:21. Schloss Dagstuhl – Leibniz-Zentrum
für Informatik, 2024.


doi:10.4230/LIPICS.CCC.2024.36.





\item[] 
[39]

Mikito Nanashima.


A theory of heuristic learnability.


In Mikhail Belkin and Samory Kpotufe, editors, \emph{Proceedings of
Thirty Fourth Conference on Learning Theory}, volume 134 of \emph{Proceedings
of Machine Learning Research}, pages 3483–3525. PMLR, August 2021.


URL: https://proceedings.mlr.press/v134/nanashima21a.html.





\item[] 
[40]

Noam Nisan and Avi Wigderson.


Hardness vs randomness.


\emph{J. Comput. Syst. Sci.}, 49(2):149–167, 1994.


doi:10.1016/S0022-0000(05)80043-1.





\item[] 
[41]

Igor C. Oliveira and Rahul Santhanam.


Conspiracies between learning algorithms, circuit lower bounds, and
pseudorandomness.


In Ryan O’Donnell, editor, \emph{32nd Computational Complexity
Conference, CCC 2017, July 6-9, 2017, Riga, Latvia}, volume 79 of \emph{LIPIcs}, pages 18:1–18:49. Schloss Dagstuhl – Leibniz-Zentrum für
Informatik, 2017.


doi:10.4230/LIPICS.CCC.2017.18.





\item[] 
[42]

Shien Jin Ong and Salil P. Vadhan.


Zero knowledge and soundness are symmetric.


In Moni Naor, editor, \emph{Advances in Cryptology – EUROCRYPT
2007, 26th Annual International Conference on the Theory and Applications of
Cryptographic Techniques, Barcelona, Spain, May 20-24, 2007, Proceedings},
volume 4515 of \emph{Lecture Notes in Computer Science}, pages 187–209.
Springer, 2007.


doi:10.1007/978-3-540-72540-4\_11.





\item[] 
[43]

Leonard Pitt and Leslie G. Valiant.


Computational limitations on learning from examples.


\emph{J. ACM}, 35(4):965–984, 1988.


doi:10.1145/48014.63140.





\item[] 
[44]

Ran Raz, Omer Reingold, and Salil P. Vadhan.


Extracting all the randomness and reducing the error in trevisan’s
extractors.


\emph{J. Comput. Syst. Sci.}, 65(1):97–128, 2002.


doi:10.1006/JCSS.2002.1824.





\item[] 
[45]

Rahul Santhanam.


Pseudorandomness and the minimum circuit size problem.


In Thomas Vidick, editor, \emph{11th Innovations in Theoretical
Computer Science Conference, ITCS 2020, January 12-14, 2020, Seattle,
Washington, USA}, volume 151 of \emph{LIPIcs}, pages 68:1–68:26. Schloss
Dagstuhl – Leibniz-Zentrum für Informatik, 2020.


doi:10.4230/LIPICS.ITCS.2020.68.





\item[] 
[46]

Leslie G. Valiant.


A theory of the learnable.


\emph{Commun. ACM}, 27(11):1134–1142, 1984.


doi:10.1145/1968.1972.





\item[] 
[47]

Andrew Chi-Chih Yao.


Theory and applications of trapdoor functions (extended abstract).


In \emph{23rd Annual Symposium on Foundations of Computer Science,
Chicago, Illinois, USA, 3-5 November 1982}, pages 80–91. IEEE Computer
Society, 1982.


doi:10.1109/SFCS.1982.45.





\end{itemize}




\par\medskip\noindent\textbf{
Appendix A WE from Indistinguishability Obfuscation}\par




In this section, we give a sketch of the construction showing that indistinguishability obfuscation for ${\mathsf{P}}/{\mathsf{poly}}$ implies witness encryption for ${\mathsf{NP}}$.\par



% SOURCE-ID Thmtheorem75


\par\medskip\noindent\textbf{
\textbf{Definition 75}\textbf{ }(Indistinguishability obfuscation (iO))\textbf{.}
}\par



\emph{A PPT algorithm $i\mathcal{O}$ is an indistinguishability obfuscator for a circuit class $\mathcal{C}$ if the following conditions are met.}\par


\begin{itemize}


\item[] 
$\blacksquare$ 


\emph{Correctness:}\emph{ for all }$C\in\mathcal{C}$\emph{ and all inputs }$x$\emph{, letting }$C^{\prime}=i\mathcal{O}(C)$\emph{, it holds that }$C^{\prime}(x)=C(x)$\emph{.}\par





\item[] 
$\blacksquare$ 


\emph{Security:}\emph{ for any efficient adversary }$A$\emph{, for any circuits }$C_{0}$\emph{ and }$C_{1}$\emph{ computing the same function, }$i\mathcal{O}(C_{0})$\emph{ and }$i\mathcal{O}(C_{1})$\emph{ are indistinguishable for }$A$\emph{.}\par





\end{itemize}



% END-SOURCE-ID Thmtheorem75


% SOURCE-ID Thmtheorem76


\par\medskip\noindent\textbf{
\textbf{Theorem 76}\textbf{ }([10])\textbf{.}
}\par



\emph{For any $L\in{\mathsf{NP}}$, iO for ${\mathsf{P}}/{\mathsf{poly}}$ implies WE for $L$.}\par



% END-SOURCE-ID Thmtheorem76



\par\medskip\noindent\textbf{Proof sketch..}\par




Let $V$ be a poly-time verifier for $L$. For an instance $x$ of $L$ and a bit $b\in\{0,1\}$, define\par


\[\begin{gathered}
F_{x,b}(w)=\begin{cases}b&V(x,w)=1\\
0&\text{otherwise}\end{cases}
\end{gathered}\]


Consider the witness encryption scheme defined as follows:\par


\begin{itemize}


\item[] 
$\blacksquare$ 


${\mathsf{Enc}}(x,b)=i\mathcal{O}(F_{x,b})$;\par





\item[] 
$\blacksquare$ 


${\mathsf{Dec}}(c,x,w)=c(w)$.\par





\end{itemize}


By the correctness of the iO, if $V(x,w)=1$, then for $C:={\mathsf{Enc}}(x,b)=i\mathcal{O}(F_{x,b})$, it holds that ${\mathsf{Dec}}(C,x,w)=C(w)=b$. Therefore, the correctness of WE is satisfied.\par




We leave the proof of security to the reader, but we note that if $x\notin L$, then $F_{x,0}$ and $F_{x,1}$ are the same function.
$\blacktriangleleft$\par






\par\medskip\noindent\textbf{
Appendix B Public-key Encryption from WE and OWFs}\par




In this section, we give a sketch of the construction showing that witness encryption, together with the existence of a one-way function, implies the possibility of public-key encryption.\par



% SOURCE-ID Thmtheorem77


\par\medskip\noindent\textbf{
\textbf{Definition 77}\textbf{ }(Public-key encryption)\textbf{.}
}\par



\emph{A public-key encryption scheme is a triple of algorithms $({\mathsf{Gen}},{\mathsf{Enc}},{\mathsf{Dec}})$ meeting the following conditions.}\par


\begin{itemize}


\item[] 
$\blacksquare$ 


\emph{Correctness:}\emph{ for all sufficiently large }$\lambda\in\mathbb{N}$\emph{ and any bit }$b\in\{0,1\}$\emph{,}\par


\[\begin{gathered}
\Pr_{(sk,pk)\sim{\mathsf{Gen}}(1^{\lambda})}[\;{\mathsf{Dec}}(sk,{\mathsf{Enc}}(pk,b))=b\;]=1.
\end{gathered}\]





\item[] 
$\blacksquare$ 


\emph{Security:}\emph{ for all sufficiently large }$\lambda\in\mathbb{N}$\emph{ and any efficient adversary }$A$\emph{,}\par


\[\begin{gathered}
\Pr_{(sk,pk)\sim{\mathsf{Gen}}(1^{\lambda})\;;\;b\sim\mathcal{U}}[\;A(pk,{\mathsf{Enc}}(pk,b))=b\;]\leq\frac{1}{2}+{\sf negl}(\lambda).
\end{gathered}\]





\end{itemize}



% END-SOURCE-ID Thmtheorem77


% SOURCE-ID Thmtheorem78


\par\medskip\noindent\textbf{
\textbf{Theorem 78}\textbf{ }([11])\textbf{.}
}\par



\emph{Suppose that WE exists for every language in ${\mathsf{NP}}$ and that a one-way function exists. Then, public-key encryption exists.}\par



% END-SOURCE-ID Thmtheorem78



\par\medskip\noindent\textbf{Proof sketch..}\par




By [14], we may assume the existence of a PRG $G:\{0,1\}^{\lambda}\to\{0,1\}^{2\lambda}$ for all sufficiently large $\lambda\in\mathbb{N}$. Consider the language\par


\[\begin{gathered}
L=\{t\in\{0,1\}^{2\lambda}\mid\exists s\in\{0,1\}^{\lambda},\;t=G(s)\},
\end{gathered}\]


and let $({\mathsf{Enc}}_{WE},{\mathsf{Dec}}_{WE})$ be a witness encryption scheme for $L$. Define a public-key encryption scheme as follows.\par


\begin{itemize}


\item[] 
$\blacksquare$ 


${\mathsf{Gen}}(1^{\lambda})$: sample $s\sim\mathcal{U}_{\lambda}$ and let $t=G(s)$. Let $sk=s$ and $pk=t$.\par





\item[] 
$\blacksquare$ 


${\mathsf{Enc}}(pk,b)$: return $(\;{\mathsf{Enc}}_{WE}(pk,b)\;,\;pk\;)$.\par





\item[] 
$\blacksquare$ 


${\mathsf{Dec}}(sk,c^{\prime}=(c,pk))$: return the output of ${\mathsf{Dec}}_{WE}(c,pk,sk)$.\par





\end{itemize}


By the correctness of the witness encryption, for any pair $(sk,pk)=(s,t)$ in the support of ${\mathsf{Gen}}$, it holds that $s$ is a witness for $t\in L$. Therefore, ${\mathsf{Dec}}(sk,{\mathsf{Enc}}(pk,b))={\mathsf{Dec}}_{WE}({\mathsf{Enc}}_{WE}(t,b),t,s)=b$. This shows the correctness of the PKE scheme.\par




We leave the proof of security to the reader.
$\blacktriangleleft$
\par




\end{document}
